\section{Surjectivity III: assembly; low weight}\label{sec:pf:onto3}

\providecommand{\ontothreeidx}{\operatorname{idx}}
\providecommand{\ontothreeTw}{\mathrm{Tw}}
\providecommand{\ontothreeON}{\sfname{ON}}
\providecommand{\ontothreeONw}{\sfname{ON}_{\ast}}
\providecommand{\ontothreeLsix}{\sfname{L6}}
\providecommand{\ontothreeClosed}{\sfname{Closed}}
\providecommand{\ontothreeEv}{\mathrm{Ev}}
\providecommand{\ontothreeMore}{\mathrm{Aux}}
\providecommand{\ontothreePC}{\mathrm{PC}}
\providecommand{\ontothreeAdel}{\mathcal{A}}
\providecommand{\ontothreeSt}[1]{\mathcal{S}_{\mathrm{#1}}}
\providecommand{\ontothreeerr}{\mathrm{err}}
\providecommand{\ontothreeerrOn}{\mathrm{err}_{\mathrm{On}}}
\providecommand{\ontothreeerrl}{\mathrm{err}_{\mathrm{l}}}
\providecommand{\ontothreeerrP}{\mathrm{err}_{\mathrm{P}}}
\providecommand{\ontothreeerrG}{\mathrm{err}_{\mathrm{G}}}
\providecommand{\ontothreeDF}{\operatorname{DF}}
\providecommand{\ontothreesize}{\operatorname{size}}
\providecommand{\ontothreePP}{\mathsf{P}}
\providecommand{\ontothreeRR}{\mathsf{R}}

This section proves the first half of \cref{prop:families:points}: the bundle $P$ of classical results on
unitary groups implies the two point inequalities $\sfname{SurjectiveOntoX}$ (the inequality $\sfname{ON}$ at the
first side of the points of the families of $\sfname{FamHyp}$) and $\sfname{LambdaLowerBound}_2$ (the inequality
$\sfname{L6}$ at the second side). The second half of \cref{prop:families:points}, the implication towards
$\sfname{ConstantFunctionFits}$, is
\lref{GlobalTheory.BigTop.Fun.Close.fourier_coefficient_fit_of_join_alt}.

\paragraph{The idea.}
The inequality $\sfname{ON}$ asserts that at the twisted places $w$ of a point the theta-type maps
$\theta_{j,\mu(w)}$ send the lattices $T_l(j,\mu(w))$ almost onto the lattices $M_j$: the generalised index
$I_j(\mu(w),w)$ is bounded by an explicit local budget, and the weighted sum of the excesses over all twisted
places is bounded globally. For the points of the families this is reduced to a single global inequality,
because the error terms of a point are the canonical ones (\cref{pf:onto3:canonical}). The global inequality is
obtained from a \emph{record object} attached to the point: a system of lattices in spaces carrying a
representation of a finite group $G$, which is a product of mirabolic subgroups of $\GL_r$ over the residue fields
at the two type primes. The local index at a place is bounded by an argument of Mackey type
(\cref{pf:onto3:GG}): for a pair $(U,\psi)$ satisfying Mackey's criterion, the $\psi$-eigenspace of $U$ captures
a fraction $1/[G:U]$ of the index of any pair of $G$-stable lattices. The remaining inputs are estimates for the
record object (``leaves''): depth bounds at the various kinds of primes, share bounds for eigenspace dimensions,
twist bounds, exponent bounds and numerical conditions. They are supplied by nine closed statements about the
families of $\sfname{FamHyp}$; two of them, the geometric statement and the exponent statement, carry most of the
arithmetic and are proved in \lref{GlobalTheory.BigOn.geoSt_of_pageSt_alt} (the page statement implies the
geometric statement) and \lref{GlobalTheory.T.expSt_of_base_change_and_crystalline_lattices} (the Galois-theoretic
inputs of $P$ imply the exponent statement); the bound for the index of theta lattices along runs, which enters the
geometric statement and the models of \cref{pf:onto3:share}, is \lref{GlobalTheory.FinalJoin.runHolds_of_parOK}.

The inequality $\sfname{L6}$ asserts that the scalars $\lambda$ by which the theta-type maps of the base type of
the second side multiply hermitian majorants are bounded below by an explicit main term. At the base infinite
place the scalars are eigenvalues of a theta lift composed with its adjoint, and the bound is an inner product
formula for the twisted lift (the adelic constant bound $\ontothreeAdel$ of \cref{pf:onto3:adelc}); at the other
infinite places it follows from the Galois conjugates of the rational theta lift.

\paragraph{Organisation.}
\Cref{sec:pf:onto3:setting} recalls the objects and the two statements. \Cref{sec:pf:onto3:main} states the main
theorem and reduces it to \cref{pf:onto3:ontoX,pf:onto3:lamlow}. \Cref{sec:pf:onto3:local} contains the
lattice theory: generalised indices, Mackey pairs, the mirabolic pairs and the local index bound.
\Cref{sec:pf:onto3:record} introduces record objects and proves the global bound
(\cref{pf:onto3:oncore}). \Cref{sec:pf:onto3:ties} contains the definitions which tie record objects to the runs of
the core of a point (record algebras, scalars, form, trace and depth ties, pairs of runs, record conditions and the
closed statements over record objects). \Cref{sec:pf:onto3:path} contains the reduction of $\sfname{SurjectiveOntoX}$
to nine statements about record objects and the proofs of those statements that belong here.
\Cref{sec:pf:onto3:low} proves the low-weight bound. \Cref{sec:pf:onto3:runw} proves the results on run weights, page
conditions, theta lattices of records and generating families which are used in the index bound along runs and in
the geometric statement, and \cref{sec:pf:onto3:remarks} collects remarks on the formal statements.

Throughout, ``closedly'' and ``eventually'' have the meaning of \cref{def:families:closed} and
\cref{sec:families:nonempty}; all statements are proved for all families of the class $\sfname{FamHyp}$
with small exponent and large constants, and (as noted after \cref{def:families:closed}) they say nothing
about classes without large members.

% ---------------------------------------------------------------------------
\subsection{The objects and the two statements}\label{sec:pf:onto3:setting}

\paragraph{Places and generalised indices.}
Let $K$ be a number field with an involutive automorphism $c$ (in the applications $K$ is the totally complex
coefficient field $\mathit{LL}$ of \cref{sec:families:theta} and $c$ is complex conjugation). A finite place
$w$ of $K$ is a non-zero prime ideal of $\OO_K$; we write $p_w$ for its residue characteristic,
$e_w=v_w(p_w)\ge1$ for its absolute ramification index, $\ell_w=\log\Nm w/[K:\Q]>0$, and $\bar w=c(w)$ for the
conjugate place; $p_{\bar w}=p_w$, $e_{\bar w}=e_w$ and $\ell_{\bar w}=\ell_w$. For a finite-dimensional
$K$-vector space $Z$ and $\OO_K$-submodules $A,B\subseteq Z$ put
\begin{equation}\label{eq:pf:onto3:idx}
\ontothreeidx_w(A,B)=\len\bigl(A_w/(A\cap B)_w\bigr)-\len\bigl(B_w/(A\cap B)_w\bigr)\in\Z ,
\end{equation}
where $X_w$ is the localisation of $X$ at $w$, lengths are lengths of $\OO_{K,w}$-modules, and an infinite length
is replaced by $0$. A \emph{lattice in} a subspace $W\subseteq Z$ is a finitely generated $\OO_K$-submodule
spanning $W$ over $K$; a \emph{full lattice} is a lattice in $Z$. For lattices $A,B$ in the same subspace all
lengths in \eqref{eq:pf:onto3:idx} are finite and $\ontothreeidx_w(A,B)$ is the $w$-adic valuation of the
generalised index $[A:B]$. For a subspace $W$ and an $\OO_K$-submodule $A$ we write $W\cap A$ for
$A\cap W$ (an $\OO_K$-submodule of $W$), and $aA$ for the image of $A$ under multiplication by $a\in K$.

\paragraph{Indefinite data.}
A \emph{base} consists of a finite index set $J$, a finite set $\mathrm{Ix}$, a set of \emph{types} with a base
type $\mu_0$, and a base infinite place $v_0$ of $K$ (and further data not used here). A \emph{side} over a base
provides, for $j\in J$ and every type $t$, a dimension $d_Z(j,t)$ and an $\OO_K$-submodule
$T_l(j,t)\subseteq K^{d_Z(j,t)}$ (together with hermitian data not used here). An \emph{indefinite datum} $S$
over a side consists of a map $\mu$ from finite places to types, a finite set $\ontothreeTw(S)$ of finite places
(the \emph{twisted places}), an integer $D(S)\ge0$, for each $j\in J$ a dimension $d_W(j)$, an
$\OO_K$-submodule $M_j\subseteq K^{d_W(j)}$ and matrices $\theta_{j,t}\in M_{d_W(j)\times d_Z(j,t)}(K)$, real
numbers $\lambda_{j,t,p,v}$ ($j\in J$, $t$ a type, $p\in\mathrm{Ix}$, $v$ an infinite place) and real numbers
$\ontothreeerr_0(w,j)$ ($w$ a finite place, $j\in J$), and further data (metrics at the infinite places, subspaces,
scalars $\kappa_v$) which play no role here. Its \emph{local indices} are
\begin{equation}\label{eq:pf:onto3:Iw}
I_j(t,w)=\ontothreeidx_w\bigl(M_j,\ \theta_{j,t}(T_l(j,t))\bigr).
\end{equation}
We say that $S$ has \emph{rank rows} if $D(S)>0$ and $d_W(j)=D(S)$ for all $j$.

\paragraph{Scales, numbers and the predicates.}
A \emph{scale} $Z$ over a base consists of functions $F_0(w,j)$ and $J_0(t,w,j)$ and real numbers $E_{1u}$,
$l_{3m}(j)$, $l_{3u}$, $l_{6m}(j)$, $l_{6u}$, $x$, $b$, $l_a$, $d_e$. For $C,C_H\in\R$ its \emph{numbers}
$N_Z(C,C_H)$ are $F_0$, $J_0$, $E_1(j)=C\,E_{1u}$, $l_3(j)=l_{3m}(j)+C\,l_{3u}$,
$l_6(j)=l_{6m}(j)-C\,l_{6u}$, $hh=x^{C_H}b$, $l_a$ and $d_e$.

\begin{definition}\label{pf:onto3:defON}
Let $S$ be an indefinite datum, $N$ numbers over the same base and $\ontothreeerr$ a real function of
$(w,j)$. Then $\ontothreeONw(S,N,\ontothreeerr)$ is the conjunction of:
\begin{enumerate}[label=(ON\arabic*)]
\item for every $w\in\ontothreeTw(S)$,
\[
\sum_{j\in J}\bigl(I_j(\mu(w),w)+I_j(\mu(w),\bar w)\bigr)\le
\Bigl(\sum_{j\in J}\bigl(2e_wF_0(w,j)+e_w\ontothreeerr(w,j)\bigr)+e_w\frac{2}{p_w^2}\Bigr)D(S);
\]
\item $\ontothreeerr(w,j)\ge0$ for all $w\in\ontothreeTw(S)$ and $j\in J$;
\item for every $j\in J$,
$\displaystyle\sum_{w\in\ontothreeTw(S)}\ell_w\bigl(e_w\ontothreeerr(w,j)+e_wJ_0(\mu(w),w,j)\bigr)\le E_1(j)$.
\end{enumerate}
$S$ satisfies $\ontothreeON(N)$ if $\ontothreeONw(S,N,\ontothreeerr_0)$ holds for its own error function
$\ontothreeerr_0$, and $\ontothreeLsix(N)$ if $l_6(j)\le\log\lambda_{j,\mu_0,p,v}$ for all $j$, $p$ and infinite
places $v$. Here $\log$ is the real logarithm extended by $\log y=0$ for $y\le0$.
\end{definition}

The predicates $\ontothreeON$ and $\ontothreeLsix$ are the inequalities $\sfname{ON}$ and $\sfname{L6}$ of
\cref{sec:families:nonempty} (with $\sum_j$ in (ON1) and the factor $e_w$ written out).

\paragraph{Families and points.}
Families $\mathcal F\in\sfname{FamP}(\eps)$, the quantifier $\ontothreeClosed(\mathcal C,\cdot)$ and largeness of
a family with respect to staged thresholds $K$ are as in \cref{def:families:closed}: $\ontothreeClosed(\mathcal
C,Q)$ means that there is $\eps_1>0$ such that for every $\eps\in(0,\eps_1]$ there are thresholds $K$ with
$Q(\mathcal F,C_H)$ for every $\mathcal F\in\mathcal C\cap\sfname{FamP}(\eps)$ and every $C_H$ at which
$\mathcal F$ is $K$-large. Largeness requires the ``pre'' constants of $\mathcal F$ to be at least $K_{\mathrm{pre}}$,
the ``middle'' constants (among them $C_H$ and the size constant $C_{\mathrm{sz}}$ of $\mathcal F$) to be at least
$K_{\mathrm{mid}}$ of the outer constants, and the last constants to be at least $K_{\mathrm{last}}$ of the outer and
middle ones. We use three closure properties \cite{LeanSource2026}: if $\ontothreeClosed(\mathcal C,Q_1)$ and
$\ontothreeClosed(\mathcal C,Q_2)$ then $\ontothreeClosed(\mathcal C,Q_1\wedge Q_2)$
(\lref{Formalism.MainNotions.Closed.and}: take the smaller $\eps_1$ and the pointwise larger thresholds);
if $\ontothreeClosed(\mathcal C,Q_1)$ and $Q_1(\mathcal F,C_H)\Rightarrow Q_2(\mathcal F,C_H)$ for all
$\eps>0$, $\mathcal F\in\sfname{FamP}(\eps)$ and $C_H$, then $\ontothreeClosed(\mathcal C,Q_2)$
(\lref{Formalism.MainNotions.Closed.imp}); and a property holding for all families and all $C_H$ is closed
(\lref{Formalism.MainNotions.Closed.of_all}).

The class $\sfname{FamHyp}$ consists of the families satisfying the conjunction of: fixed numerical constants
and host conditions; the rank windows \eqref{eq:families:rankval}; the condition that the first modulus $q_0$ of
every member is smaller than the scale parameter $X'$; eventually, tameness and distinctness of the characters at the
places above the type primes; eventually, a condition on the chosen data $G_{x,i}$ of the members; eventual
divisibility of the constants $D_\Gamma$ of the members by an explicit Minkowski-type number; box and congruence
conditions on the host primes; $C_{YL}\le9$; eventually, every prime divisor of $D_\Gamma$ is at most $2fN(m)+1$;
the twist goodness $\sfname{TwGood}$; the existence of $C_H$ with $\sfname{ChosenE}(\mathcal F,3,C_H)$; and the size
conditions $\sfname{TuSize}$ and $\sfname{DGamSize}$ (the conditions of these names in \cref{sec:families}).

A \emph{point} $p$ of a family $\mathcal F$ over $(x,i)$, $x>x_1$, $i\in I_x$, is an element of
$\mathrm{Pt}(\mathcal F,x,i)$ (\cref{sec:families:nonempty}); it consists of a row-side family (a finite group
$G_p$, three-sided families of forms, subgroups $P_0,P_s\le G_p$ with characters $\psi_0,\psi_s$, two indefinite
data $S_0,S_1$, \dots), a core, a real datum $R$, a face, weight families $\kappa$, vectors $v_0$, ties
$\mathrm{fw}$, and about forty compatibility fields. From $p$ one obtains two indefinite data $S_{\mathrm f}$ (the
first side) and $S_{\mathrm t}$ (the second side) over a common base, which coincide with $S_0$ and $S_1$, a scale
$Z_p$ and the numbers $N_p(C,C_H)=N_{Z_p}(C,C_H)$. Among the fields of a point is the identity
\begin{equation}\label{eq:pf:onto3:erreq}
\ontothreeerr_0(w,j)=\frac{\max\Bigl(0,\ \dfrac{1}{D}\displaystyle\sum_{j'\in J}\bigl(I_{j'}(\mu(w),w)+
I_{j'}(\mu(w),\bar w)\bigr)-\sum_{j'\in J}2e_wF_0(w,j')-e_w\frac{2}{p_w^2}\Bigr)}{e_w\,|J|}
\end{equation}
for the error function of $S_{\mathrm f}$, all finite places $w$ and all $j$, where $D=D(S_{\mathrm f})$ and
$F_0$ is that of $Z_p$ (division by $0$ gives $0$). We call the right-hand side the \emph{canonical error}.

The \emph{point conditions} $\ontothreePC(p)$ are seventeen conditions on a point: existence conditions at the
infinite places, three evaluation and action conditions, an extension condition for the theta data, the one-map
rows for all indices, the primitive rows, separation of types, a condition on the family of forms, frame
conditions, the division rows, a tie between the characters $\psi$, the identity $D_0=d_Z$ for the base type, the
export row, a goodness condition on the datum, the condition that every place at which the complex of $p$ has a
single index for the numbers $N_p(0,0)$ is twisted, and the compatibility conditions. The \emph{auxiliary
conditions} $\ontothreeMore(p)$ are the conjunction of three conditions: every place at which the complex of $p$ has
a single index is twisted; the core ties of the two-sided reduction of $p$ (\cref{sec:pf:onto3:path}), which
contain a canonical-core condition and a half-error condition; and the record ties of $p$. For a family
$\mathcal F$ and a property $Q$ of points,
\[
\ontothreeEv(\mathcal F,Q):\iff\exists x_0\ \forall x>x_1\text{ with }x\ge x_0\ \forall i\in I_x\
\forall p\in\mathrm{Pt}(\mathcal F,x,i):\ \ontothreePC(p)\Rightarrow Q(x,i,p).
\]

\begin{definition}[The two statements]\label{pf:onto3:stmts}
For a family $\mathcal F$ and $C_H\in\R$ let $Q_{\mathrm{ON}}(\mathcal F,C_H)$ be the property: if $\mathcal F$
satisfies $\sfname{DIVinN}$, then there is $C_0\ge0$ with
\[
\ontothreeEv\Bigl(\mathcal F,\ p\mapsto\bigl(\ontothreeMore(p)\Rightarrow\forall C\ge C_0:\
S_{\mathrm f}\text{ satisfies }\ontothreeON(N_p(C,C_H))\bigr)\Bigr);
\]
and let $Q_{\mathrm{L6}}(\mathcal F,C_H)$ be the property $\ontothreeEv\bigl(\mathcal F,\ p\mapsto\forall C\in\R:\
S_{\mathrm t}\text{ satisfies }\ontothreeLsix(N_p(C,C_H))\bigr)$. Then $\sfname{SurjectiveOntoX}$ is the statement
$\ontothreeClosed(\sfname{FamHyp},Q_{\mathrm{ON}})$, and $\sfname{LambdaLowerBound}_2$ is the statement
$\ontothreeClosed(\sfname{FamHyp},Q_{\mathrm{L6}})$.
\end{definition}

Here $\sfname{DIVinN}(\mathcal F)$ is the condition of \cref{sec:families:nonempty}: there are $C,x_0$ such that
for $x\ge x_0$, $i\in I_x$ and every lift index $j$ of $D_{x,i}$, the weighted divisibility defect
$\sum_w\ell^{P}_w e_w\eta_j(w)$ over the places $w$ of the support of $L_{x,i}$ with $\mu_w\ne\mu_0$ is at most
$D_0Cx^{\eps}(B_0+N_c)$. It is not used in the proof (\cref{pf:onto3:rem:divinn}).

% ---------------------------------------------------------------------------
\subsection{The main theorem}\label{sec:pf:onto3:main}

The bundle $P=\sfname{UnitaryGroupThirdPlaceHypotheses}$ of \cref{sec:families:theta} is a conjunction of
classical results, each of which is a theorem of the formal development. We use the following ones, in the
formulations of the development:
\begin{enumerate}[label=(P\arabic*),leftmargin=*]
\item\label{pf:onto3:P:green} Green's parametrisation of the cuspidal characters of $\GL_d(\F_g)$ \cite{Green1955};
\item\label{pf:onto3:P:clozel} Clozel's theorem that regular algebraic cuspidal representations are cohomological
\cite{Clozel1990};
\item\label{pf:onto3:P:hecke} the $\ell$-adic characters attached to algebraic Hecke characters, including the
global Artin maps \cite{Serre1968};
\item\label{pf:onto3:P:cheb} Chebotarev's density theorem with primes of degree one;
\item\label{pf:onto3:P:galrep} Galois representations attached to polarised regular algebraic cuspidal
representations \cite{ClozelHarrisLabesse2011,ChenevierHarris2013};
\item\label{pf:onto3:P:bc} base change for definite unitary groups \cite{Labesse2011};
\item\label{pf:onto3:P:bs} the realisation of crystalline lattices by prismatic $F$-crystals
\cite{BhattScholze2023};
\item\label{pf:onto3:P:above} the local properties above $\ell$ of the $\ell$-adic characters of (P3);
\item\label{pf:onto3:P:decay} rapid decay of cusp forms on unitary groups \cite{HarishChandra1968,MoeglinWaldspurger1995};
\item\label{pf:onto3:P:llc} the existence of local Langlands parameters (the proposition $\sfname{RecTExists}$ of
\cref{sec:galois});
\item\label{pf:onto3:P:ds} the classification of the discrete series of $\Uni(p,q)$ \cite{HarishChandra1966};
\item\label{pf:onto3:P:cw} Casselman--Wallach globalisation \cite{Casselman1989,Wallach1992};
\item\label{pf:onto3:P:ws} the convergence of theta integrals and the Siegel--Weil identity for unitary dual pairs
in Weil's range (Weil--Siegel formula) \cite{Weil1965,Ichino2004};
\item\label{pf:onto3:P:sw} a second version of the Siegel--Weil formula in the convergent range;
\item\label{pf:onto3:P:dm} the local zeta integrals of the doubling method \cite{GelbartPiatetskiShapiroRallis1987};
\item\label{pf:onto3:P:udm} the unramified computation of the doubling integral for $\GL_n$
\cite{GelbartPiatetskiShapiroRallis1987,LapidRallis2005};
\item\label{pf:onto3:P:bz} Bernstein--Zelevinsky finiteness \cite{BernsteinZelevinsky1976};
\item\label{pf:onto3:P:fd} Harish-Chandra's formula for the formal degrees of the discrete series
\cite{HarishChandra1966}.
\end{enumerate}

\begin{theorem}[\cref{prop:families:points}, first part]\label{pf:onto3:main}
The proposition $P=\sfname{UnitaryGroupThirdPlaceHypotheses}$ implies
$\sfname{SurjectiveOntoX}\wedge\sfname{LambdaLowerBound}_2$.
\formallabel{GlobalTheory.Fc.S0.ontoX_and_lamLow_of_unitaryGroupHypotheses'}
\end{theorem}

\begin{proof}
The first conjunct is \cref{pf:onto3:ontoX}, applied to \ref{pf:onto3:P:green}--\ref{pf:onto3:P:cw}. The
proposition $\sfname{RecTExists}$ \ref{pf:onto3:P:llc} is passed through
\lref{UpperTheory.One.M.New.llOf_alt}, which takes twelve fields of $P$ and returns the first one,
$\sfname{RecTExists}$. The second conjunct is \cref{pf:onto3:lamlow}, applied to \ref{pf:onto3:P:decay} and
\ref{pf:onto3:P:ws}--\ref{pf:onto3:P:fd}.
\end{proof}

% ---------------------------------------------------------------------------
\subsection{Local index bounds}\label{sec:pf:onto3:local}

Throughout this subsection $K$ is a number field, $Z$ a finite-dimensional $K$-vector space, $w$ a finite place
of $K$ and $v=v_w$ the normalised valuation. All $\OO_K$-submodules are submodules of $Z$ (or of a second
$K$-vector space where indicated).

\subsubsection{Generalised indices}

We use the following standard facts on lattices over the discrete valuation ring $\OO_{K,w}$
\cite[Ch.~III, \S1]{Serre1979}: for lattices $A,B,C$ in a subspace $W$ one has
$\ontothreeidx_w(A,C)=\ontothreeidx_w(A,B)+\ontothreeidx_w(B,C)$; if $B\subseteq A$ then
$\ontothreeidx_w(A,B)=\len(A_w/B_w)\ge0$; and $\ontothreeidx_w(B,aB)=\dim W\cdot v(a)$ for $a\in K^\times$.

\begin{lemma}\label{pf:onto3:idxbasic}
Let $W'\subseteq W\subseteq Z$ be subspaces.
\begin{enumerate}[label=(\roman*)]
\item If $A,B,B'$ are lattices in $W$ and $B\subseteq B'$, then
$\ontothreeidx_w(A,B')\le\ontothreeidx_w(A,B)$.
\item If $A,B$ are lattices in $W$ with $B\subseteq A$, then $W'\cap A$ is a lattice in $W'$ and
$0\le\ontothreeidx_w(W'\cap A,W'\cap B)\le\ontothreeidx_w(A,B)$.
\item If $A,B$ are lattices in $W$ and $a\in K^\times$, then $aB$ is a lattice in $W$ and
$\ontothreeidx_w(A,aB)=\ontothreeidx_w(A,B)+\dim W\cdot v(a)$.
\item If $A,B$ are full lattices and $a\in K^\times$, then $A\cap aB$ is a full lattice.
\item If $\varphi\colon K^d\to Z$ is an injective $K$-linear map and $A,B\subseteq K^d$ are
$\OO_K$-submodules, then $\ontothreeidx_w(\varphi A,\varphi B)=\ontothreeidx_w(A,B)$.
\end{enumerate}
\formallabel{Formalism.EigenvalueLevelCeiling.idxMono}
\formallabel{Formalism.EigenvalueLevelCeiling.index_monotone}
\formallabel{Formalism.EigenvalueLevelCeiling.idxCutLe}
\formallabel{Formalism.EigenvalueLevelCeiling.idxScaleIn}
\formallabel{Formalism.EigenvalueLevelCeiling.latInter}
\formallabel{Formalism.EigenvalueLevelCeiling.lattice_intersection}
\formallabel{Formalism.EigenvalueLevelCeiling.idxTransport}
\end{lemma}

\begin{proof}
(i) By additivity $\ontothreeidx_w(A,B)-\ontothreeidx_w(A,B')=\ontothreeidx_w(B',B)=\len(B'_w/B_w)\ge0$.
(ii) $W'\cap A$ is finitely generated since $\OO_K$ is noetherian, and it spans $W'$ because every vector of
$W'$ has a non-zero integral multiple in $A$. The map $(W'\cap A)_w/(W'\cap B)_w\to A_w/B_w$ is injective, which
gives the inequalities. (iii) $aB$ is a lattice in $W$ and $\ontothreeidx_w(A,aB)=\ontothreeidx_w(A,B)+
\ontothreeidx_w(B,aB)$. (iv) $aB$ is a full lattice, so it contains $dA$ for some non-zero $d\in\OO_K$; hence
$dA\subseteq A\cap aB\subseteq A$. (v) $\varphi$ induces isomorphisms of all modules in
\eqref{eq:pf:onto3:idx}.
\end{proof}

\begin{lemma}\label{pf:onto3:threeforms}
Let $B\subseteq A$ be full lattices, $a\in K^\times$ with $v(a)=E\ge0$ and $r\in\R$. Then
\[
\ontothreeidx_w(W\cap A,W\cap B)\le E\dim W+r\ \text{ for all subspaces }W
\quad\iff\quad \ontothreeidx_w(A,\,A\cap a^{-1}B)\le r .
\]
\formallabel{Formalism.EigenvalueLevelCeiling.threeForms}
\formallabel{Formalism.EigenvalueLevelCeiling.three_forms}
\end{lemma}

\begin{proof}
All quantities depend only on the localisations at $w$, and $(A\cap a^{-1}B)_w=A_w\cap a^{-1}B_w$. By the
elementary divisor theorem over $\OO_{K,w}$ there is a basis $f_1,\dots,f_n$ of $A_w$ and integers $n_i\ge0$ with
$B_w=\bigoplus\pi^{n_i}\OO_{K,w}f_i$, $\pi$ a uniformiser. Then $A_w\cap a^{-1}B_w=\bigoplus\pi^{\max(0,n_i-E)}
\OO_{K,w}f_i$, so $\ontothreeidx_w(A,A\cap a^{-1}B)=\sum_i\max(0,n_i-E)$.
($\Rightarrow$) Take $W$ spanned by the $f_i$ with $n_i>E$. Then $\ontothreeidx_w(W\cap A,W\cap B)=\sum_{n_i>E}n_i
=E\dim W+\sum_i\max(0,n_i-E)$, and the hypothesis gives the claim.
($\Leftarrow$) Put $C=a^{-1}B$, so that $B\subseteq A\cap C$. For a subspace $W$, by additivity and
\cref{pf:onto3:idxbasic}(ii),
$\ontothreeidx_w(W\cap A,W\cap B)=\ontothreeidx_w(W\cap A,W\cap A\cap C)+\ontothreeidx_w(W\cap A\cap C,W\cap B)$.
The first term is at most $\ontothreeidx_w(A,A\cap C)\le r$, since $(W\cap A)/(W\cap A\cap C)$ injects into
$A/(A\cap C)$; the second is at most $\ontothreeidx_w(W\cap C,W\cap B)=\ontothreeidx_w(W\cap a^{-1}B,W\cap B)=
E\dim W$ by (i) and (iii) of \cref{pf:onto3:idxbasic}.
\end{proof}

\subsubsection{Mackey pairs}

Let $G$ be a finite group and $\rho\colon G\to\GL(Z)$ a representation over $K$. For a subgroup $U\le G$ and a
function $\psi\colon G\to\OO_K$ put
\[
Z^\psi=\{z\in Z:\ \rho(u)z=\psi(u)z\ \text{for all }u\in U\},\qquad
e_\psi=\sum_{u\in U}\psi(u^{-1})\rho(u)\in\End_K(Z).
\]

\begin{definition}\label{pf:onto3:defmackey}
$(U,\psi)$ is a \emph{Mackey pair} in $G$ if $\psi(1)=1$, $\psi(uu')=\psi(u)\psi(u')$ for $u,u'\in U$, and for
every $h\in G\smallsetminus U$ there is $u\in U$ with $h^{-1}uh\in U$ and $\psi(h^{-1}uh)\ne\psi(u)$.
\end{definition}

If $\psi|_U$ takes values in the roots of unity, this is Mackey's criterion for the irreducibility of
$\Ind_U^G\psi|_U$ in the form ``$\psi$ and its conjugate by $h$ differ on $U\cap hUh^{-1}$ for $h\notin U$''
\cite[\S7.4]{Serre1977}; we only use the elementwise statement.

\begin{lemma}\label{pf:onto3:orth}
Let $(U,\psi)$ be a Mackey pair. Then for $h\in G\smallsetminus U$ and $z\in Z^\psi$ one has
$e_\psi(\rho(h)z)=0$. Moreover $e_\psi(Z)\subseteq Z^\psi$, $e_\psi z=|U|z$ for $z\in Z^\psi$, and $e_\psi$
maps every $G$-stable $\OO_K$-submodule into itself.
\formallabel{Formalism.EigenvalueLevelCeiling.orthogonality_of_mackey}
\end{lemma}

\begin{proof}
Note that $\psi(u^{-1})=\psi(u)^{-1}$ for $u\in U$. Let $h\notin U$ and choose $u_0\in U$ with $h^{-1}u_0h\in U$
and $\psi(h^{-1}u_0h)\ne\psi(u_0)$; put $s=e_\psi(\rho(h)z)$. Substituting $u\mapsto uu_0$ and using
$\psi(u^{-1})=\psi(u_0)\psi((uu_0)^{-1})$ gives $e_\psi(\rho(u_0)\rho(h)z)=\psi(u_0)s$. On the other hand
$\rho(u_0)\rho(h)z=\rho(h)\rho(h^{-1}u_0h)z=\psi(h^{-1}u_0h)\rho(h)z$, so the same vector equals
$\psi(h^{-1}u_0h)s$. Hence $(\psi(u_0)-\psi(h^{-1}u_0h))s=0$ and $s=0$. For $u'\in U$, the substitution
$u\mapsto u'^{-1}u$ gives $\rho(u')e_\psi x=\psi(u')e_\psi x$, so $e_\psi x\in Z^\psi$; and for $z\in Z^\psi$ each
term of $e_\psi z$ equals $z$. The last assertion holds because $\psi$ takes values in $\OO_K$.
\end{proof}

\begin{lemma}\label{pf:onto3:ggcore}
Let $R$ be a commutative ring, $n\in R$, $\iota$ a finite set, and $A'\subseteq A$, $B'\subseteq B$, $C$
submodules of an $R$-module. Let $s_i,t_i$ ($i\in\iota$) be $R$-linear endomorphisms with $s_i(A)\subseteq B$,
$s_i(A')\subseteq B'$, $t_i(B)\subseteq A$, $t_i(B')\subseteq A'$, $t_is_j=\delta_{ij}\,n$ on $A$, and
$C\cap A\subseteq nA$. Then, with lengths in $\N\cup\{\infty\}$,
\[
|\iota|\cdot\len(A/A')\le\len(B/B')+|\iota|\cdot\len\bigl(A/(C\cap A)\bigr).
\]
\formallabel{Formalism.EigenvalueLevelCeiling.gg_core}
\end{lemma}

\begin{proof}
Put $Q=A/A'$. The maps $s_i,t_i$ induce $\bar s_i\colon Q\to B/B'$ and $\bar t_i\colon B/B'\to Q$ with
$\bar t_i\bar s_j=\delta_{ij}n$. The image of $T=(\bar t_i)_i\colon B/B'\to Q^\iota$ contains $(nQ)^\iota$: the
family $(ny_i)_i$ is the image of $\sum_j\bar s_j(y_j)$. Hence $|\iota|\len(nQ)\le\len(B/B')$. Since $C\cap
A\subseteq nA$, the module $Q/nQ$ is a quotient of $A/(C\cap A)$, and $\len(Q)=\len(nQ)+\len(Q/nQ)$.
\end{proof}

\begin{lemma}[Mackey index inequality]\label{pf:onto3:GG}
Let $(U,\psi)$ be a Mackey pair in $G$ and $B'\subseteq B$ full lattices stable under $\rho(G)$. Then
\[
[G:U]\cdot\ontothreeidx_w\bigl(Z^\psi\cap B,\ Z^\psi\cap B'\bigr)\le\ontothreeidx_w(B,B')+\dim Z\cdot v(|U|).
\]
\formallabel{Formalism.EigenvalueLevelCeiling.lemmaGG}
\formallabel{Formalism.EigenvalueLevelCeiling.lemmaGG_of}
\end{lemma}

\begin{proof}
Choose representatives $h_i$ of the cosets $i\in G/U$, and put $N=|U|$, $W=Z^\psi$. By \cref{pf:onto3:orth},
for $a\in W$ we have $e_\psi\rho(h_i)^{-1}\rho(h_j)a=e_\psi\rho(h_i^{-1}h_j)a=\delta_{ij}Na$, because
$h_i^{-1}h_j\in U$ only for $i=j$. In particular the map $W^{G/U}\to Z$, $(z_i)\mapsto\sum_i\rho(h_i)z_i$, is
injective (apply $e_\psi\rho(h_i)^{-1}$), so that $[G:U]\dim W\le\dim Z$.
Apply \cref{pf:onto3:ggcore} over $R=\OO_K$ to the localisations at $w$ of $A=W\cap B$, $A'=W\cap B'$,
$B$, $B'$ and $C=N(W\cap B)$, with $s_i=\rho(h_i)$ and $t_i=e_\psi\circ\rho(h_i)^{-1}$ and $n=N$: the inclusions
hold because $B,B'$ are $G$-stable and $e_\psi$ maps them into themselves and into $W$
(\cref{pf:onto3:orth}); and $C\subseteq NA$. All lengths are finite, $B'$ and $W\cap B'$ being lattices in the
same spaces as $B$ and $W\cap B$ (\cref{pf:onto3:idxbasic}). By \cref{pf:onto3:idxbasic}(iii),
$\len(A_w/C_w)=\dim W\cdot v(N)$. Thus
$[G:U]\ontothreeidx_w(W\cap B,W\cap B')\le\ontothreeidx_w(B,B')+[G:U]\dim W\cdot v(N)$, and $[G:U]\dim W\le\dim
Z$ with $v(N)\ge0$ gives the claim.
\end{proof}

\begin{lemma}\label{pf:onto3:mackeyprod}
Let $G=\prod_{v\in\Upsilon}G_v$ be a finite product of finite groups and $(U_v,\psi_v)$ Mackey pairs in $G_v$.
Then $(\prod_vU_v,\ h\mapsto\prod_v\psi_v(h_v))$ is a Mackey pair in $G$.
\formallabel{Formalism.EigenvalueLevelCeiling.mackey_product}
\end{lemma}

\begin{proof}
The first two conditions hold componentwise. If $h\notin\prod U_v$, choose $v_0$ with $h_{v_0}\notin U_{v_0}$ and
$u_0\in U_{v_0}$ for $h_{v_0}$ as in \cref{pf:onto3:defmackey}; the element $u$ equal to $u_0$ at $v_0$ and to $1$
elsewhere satisfies $h^{-1}uh\in\prod U_v$ and $\psi(h^{-1}uh)=\psi_{v_0}(h_{v_0}^{-1}u_0h_{v_0})\ne
\psi_{v_0}(u_0)=\psi(u)$, since $\psi_v(1)=1$.
\end{proof}

\paragraph{Mirabolic pairs.}
Let $k$ be a finite field, $n\ge0$, and let $\mathrm{Mir}_{n+1}(k)\le\GL_{n+1}(k)$ be the mirabolic subgroup of
matrices whose last row is $(0,\dots,0,1)$, and $N_{n+1}(k)\le\mathrm{Mir}_{n+1}(k)$ the group of upper
unitriangular matrices. For an additive character $\chi\colon k\to\OO_K^\times$ put
$\psi_\chi(h)=\chi\bigl(\sum_{i=1}^{n}h_{i,i+1}\bigr)$ (rows and columns numbered $1,\dots,n+1$).

\begin{lemma}\label{pf:onto3:mirabolic}
If $\chi$ is non-trivial, then $(N_{n+1}(k),\psi_\chi)$ is a Mackey pair in $\mathrm{Mir}_{n+1}(k)$. Consequently,
for finitely many finite fields $k_v$ and non-trivial characters $\chi_v$, the pair
$\bigl(\prod_vN_{n+1}(k_v),\ h\mapsto\prod_v\psi_{\chi_v}(h_v)\bigr)$ is a Mackey pair in
$\prod_v\mathrm{Mir}_{n+1}(k_v)$.
\formallabel{Formalism.EigenvalueLevelCeiling.mirabolic_is_mackey}
\formallabel{Formalism.EigenvalueLevelCeiling.mirabolic_mackey_pair}
\formallabel{Formalism.EigenvalueLevelCeiling.isMackey_pi_upper_unipotent_prod_character}
\end{lemma}

\begin{proof}
$\psi_\chi(1)=\chi(0)=1$, and for unitriangular $u,u'$ the superdiagonal of $uu'$ is the sum of those of $u$ and
$u'$, so $\psi_\chi$ is multiplicative on $N=N_{n+1}(k)$. Let $h\in\mathrm{Mir}_{n+1}(k)\smallsetminus N$. Call a row
$j$ of $h$ \emph{unitriangular} if $h_{jj}=1$ and $h_{jb}=0$ for $b<j$. The last row is unitriangular and not all
rows are, so there is a largest index $s\le n$ whose row is not unitriangular; all rows $j>s$ are unitriangular.
Choose $l\le s$ with $h_{sl}\ne\delta_{ls}$ (take $l=s$ if $h_{ss}\ne1$, and otherwise $l<s$ with $h_{sl}\ne0$).
Since row $s+1$ is unitriangular and $l+1\le s+1$, $h_{s+1,l+1}=\delta_{ls}$ and $h_{s+1,b}=0$ for $b\le l$.
For $t\in k$ let $E_{ab}$ be the matrix units and put
\[
g=1+t\sum_ah_{al}E_{a,s+1},\qquad g'=1+t\sum_bh_{s+1,b}E_{lb}.
\]
Then $gh=hg'$ (both equal $h+t\,(h_{al}h_{s+1,b})_{a,b}$), so $h^{-1}gh=g'$. The matrix $g$ is unitriangular,
since $h_{al}=0$ for $a\ge s+1$ ($a>s$ and $l<a$); and $g'$ is unitriangular, since $h_{s+1,b}=0$ for $b\le l$.
The only superdiagonal entries are $g_{s,s+1}=th_{sl}$ and $g'_{l,l+1}=th_{s+1,l+1}=t\delta_{ls}$, so
$\psi_\chi(g)=\chi(th_{sl})$ and $\psi_\chi(h^{-1}gh)=\chi(t\delta_{ls})$. Put $c=\delta_{ls}-h_{sl}\ne0$, choose $x$
with $\chi(x)\ne1$ and $t=x/c$; then $t\delta_{ls}=th_{sl}+x$ and $\psi_\chi(h^{-1}gh)=\psi_\chi(g)\chi(x)\ne
\psi_\chi(g)$. The last assertion follows from \cref{pf:onto3:mackeyprod}.
\end{proof}

\begin{remark}
For $\chi$ with values in $\C^\times$, the first assertion is Mackey's criterion for the irreducibility of
$\Ind_{N}^{\mathrm{Mir}}\psi_\chi$, a classical fact about the mirabolic group over a finite field
\cite{GelfandKazhdan1975,BernsteinZelevinsky1976}.
\end{remark}

\subsubsection{The local bound}

\begin{lemma}\label{pf:onto3:psigood}
Let $Z_\psi,V_0\subseteq Z$ be subspaces, $L=Z_\psi\cap V_0$, and let $M,T,T_K$ be full lattices with $T\subseteq
M$ and $T_0\subseteq Z_\psi$ a finitely generated $\OO_K$-submodule. Let $\sigma\in K^\times$ with $T\subseteq\sigma
T_K$, $u\ge1$ an integer with $u\,(Z_\psi\cap T_K)\subseteq T_0$, and $a\in K^\times$ with $v(a)=E\ge0$. Put
$N=M\cap a^{-1}T$ and assume $\ontothreeidx_w(M,N)\le r$ and, for an integer $\delta>0$,
$\delta\,\ontothreeidx_w(Z_\psi\cap M,Z_\psi\cap N)\le\ontothreeidx_w(M,N)+\dim Z\cdot v(u)$. Then $V_0\cap T_0$ is
a lattice in $L$ and
\[
\ontothreeidx_w(L\cap M,\ V_0\cap T_0)\le\bigl(E-v(\sigma)+v(u)\bigr)\dim L+\frac{r+\dim Z\cdot v(u)}{\delta}.
\]
\formallabel{Formalism.EigenvalueLevelCeiling.idxPsiGood_of}
\formallabel{Formalism.EigenvalueLevelCeiling.idxPsiGood}
\end{lemma}

\begin{proof}
$Z_\psi\cap T_K$ is a lattice in $Z_\psi$ (\cref{pf:onto3:idxbasic}(ii)), hence so is $T_0$, which lies between
$u(Z_\psi\cap T_K)$ and $Z_\psi$; therefore $V_0\cap T_0=L\cap T_0$ is a lattice in $L$. By
\cref{pf:onto3:idxbasic}(iv) $N$ is a full lattice. We claim
$u\,a\sigma^{-1}(L\cap N)\subseteq V_0\cap T_0$. Let $z\in L\cap N$. Then $z=a^{-1}t$ with $t\in T\subseteq\sigma
T_K$, so $k:=a\sigma^{-1}z\in T_K$; since $z\in L$, also $k\in L\subseteq Z_\psi$. Hence $k\in Z_\psi\cap T_K$ and
$uk\in T_0$; and $uk\in V_0$. By \cref{pf:onto3:idxbasic}(i) and (iii),
\begin{align*}
\ontothreeidx_w(L\cap M,V_0\cap T_0)&\le\ontothreeidx_w\bigl(L\cap M,ua\sigma^{-1}(L\cap N)\bigr)\\
&=\ontothreeidx_w(L\cap M,L\cap N)+\dim L\,\bigl(v(u)+E-v(\sigma)\bigr).
\end{align*}
By \cref{pf:onto3:idxbasic}(ii) applied in $Z_\psi$ (note $L\cap(Z_\psi\cap X)=L\cap X$),
$\ontothreeidx_w(L\cap M,L\cap N)\le\ontothreeidx_w(Z_\psi\cap M,Z_\psi\cap N)\le(r+\dim Z\cdot v(u))/\delta$.
\end{proof}

\begin{proposition}\label{pf:onto3:onlocal}
Let $(U,\psi)$ be a Mackey pair in a finite group $G$ acting on $Z$ by $\rho$, let $e\in\End_K(Z)$ be an idempotent
with image $V_0$, and put $Z_g=Z^\psi\cap V_0$. Let $M,T,T_K$ be full lattices with $T\subseteq M$, $M$ and $T$ stable
under $\rho(G)$, $\sigma\in K^\times$ with $T\subseteq\sigma T_K$, $T_0\subseteq Z^\psi$ finitely generated with
$|U|(Z^\psi\cap T_K)\subseteq T_0$, and $a\in K^\times$, $E\in\N$, $r\in\R$ with $v(a)=E$ and
$\ontothreeidx_w(W\cap M,W\cap T)\le E\dim W+r$ for all subspaces $W$. Let $X_3\subseteq X_4$ be lattices in
$Z_g$ such that $e(T_0)$ is a lattice in $Z_g$ and $\ontothreeidx_w(Z_g\cap M,X_3)\le\ontothreeidx_w(Z_g\cap
M,e(T_0))+\eps$. Let $\varphi\colon K^d\to Z$ be injective with image $Z_g$ and $M',T'\subseteq K^d$ with
$\varphi(M')=Z_g\cap M$ and $\varphi(T')=X_4$. Then
\[
\ontothreeidx_w(M',T')\le\bigl(E-v(\sigma)+v(|U|)\bigr)d+\frac{r+\dim Z\cdot v(|U|)}{[G:U]}+\eps.
\]
\formallabel{Formalism.EigenvalueLevelCeiling.onLocal}
\formallabel{Formalism.EigenvalueLevelCeiling.onLocal_of}
\end{proposition}

\begin{proof}
Put $N=M\cap a^{-1}T$. By \cref{pf:onto3:threeforms}, $\ontothreeidx_w(M,N)\le r$. $N$ is a full lattice and is
$\rho(G)$-stable, since $M$ and $T$ are. \Cref{pf:onto3:GG} applied to $N\subseteq M$ gives
$[G:U]\ontothreeidx_w(Z^\psi\cap M,Z^\psi\cap N)\le\ontothreeidx_w(M,N)+\dim Z\cdot v(|U|)$, and
\cref{pf:onto3:psigood} (with $Z_\psi=Z^\psi$, $u=|U|$, $\delta=[G:U]$) gives
$\ontothreeidx_w(Z_g\cap M,V_0\cap T_0)\le(E-v(\sigma)+v(|U|))\dim Z_g+(r+\dim Z\cdot v(|U|))/[G:U]$. Since $e$ is
the identity on $V_0$, $V_0\cap T_0\subseteq e(T_0)$. By \cref{pf:onto3:idxbasic}(i), applied in $Z_g$ (where
$Z_g\cap M$ is a lattice by \cref{pf:onto3:idxbasic}(ii)),
\[
\ontothreeidx_w(Z_g\cap M,X_4)\le\ontothreeidx_w(Z_g\cap M,X_3)\le\ontothreeidx_w(Z_g\cap M,e(T_0))+\eps
\le\ontothreeidx_w(Z_g\cap M,V_0\cap T_0)+\eps .
\]
Finally $\ontothreeidx_w(M',T')=\ontothreeidx_w(Z_g\cap M,X_4)$ by \cref{pf:onto3:idxbasic}(v), and $\dim Z_g=d$.
\end{proof}

\begin{proposition}[Weighted local conditions]\label{pf:onto3:lcw}
Let $G,\rho,(U,\psi),e,V_0,Z_g$ be as in \cref{pf:onto3:onlocal}. Let $M,T_1,T,T_K$ be full lattices with
$T_1\subseteq T\subseteq M$, $M$ and $T$ stable under $\rho(G)$, $\sigma\in K^\times$ with $T\subseteq\sigma T_K$,
$T_0\subseteq Z^\psi$ finitely generated with $|U|(Z^\psi\cap T_K)\subseteq T_0$, $T_J$ an $\OO_K$-submodule with
$T_0\subseteq T_J$ such that $e(T_0)$ and $e(T_J)$ are lattices in $Z_g$, and $a\in K^\times$, $E\in\N$, $r\in\R$
with $v(a)=E$ and $\ontothreeidx_w(W\cap M,W\cap T_1)\le E\dim W+r$ for all subspaces $W$. Let $X$ be a
$K$-vector space, $\theta'\colon X\to Z$ and $e_X,\mathrm{pr}\in\End_K(X)$ linear, and $M_t\subseteq X$ an
$\OO_K$-submodule, with $\theta'e_X=e\theta'$ and $x-\mathrm{pr}(x)\in\ker\theta'$ for $x\in e_X(X)$. Let $X_3$ be a
lattice in $Z_g$ such that $\mathrm{Im}=e\theta'(M_t)$ is a lattice in $Z_g$, $(X_3)_w\subseteq\mathrm{Im}_w$ and
$\ontothreeidx_w(Z_g\cap M,X_3)\le\ontothreeidx_w(Z_g\cap M,e(T_J))+\eps$. Let $\varphi\colon K^d\to Z$ be
injective with image $Z_g$, and $M',T'\subseteq K^d$ with $\varphi(M')=Z_g\cap M$ and
$\varphi(T')=\theta'\mathrm{pr}\,e_X(M_t)$. Then
\[
\ontothreeidx_w(M',T')\le\bigl(E-v(\sigma)+v(|U|)\bigr)d+\frac{r+\dim Z\cdot v(|U|)}{[G:U]}+\eps.
\]
\formallabel{Formalism.EigenvalueLevelCeiling.local_conditions_weighted_at_place}
\formallabel{Formalism.EigenvalueLevelCeiling.local_conditions_weighted}
\formallabel{Formalism.EigenvalueLevelCeiling.local_conditions_family}
\end{proposition}

\begin{proof}
(a) \emph{The initial lattice may be enlarged:} for every subspace $W$, $W\cap T_1\subseteq W\cap T\subseteq W\cap
M$ and \cref{pf:onto3:idxbasic}(i),(ii) give $\ontothreeidx_w(W\cap M,W\cap T)\le\ontothreeidx_w(W\cap M,W\cap
T_1)\le E\dim W+r$.
(b) \emph{Theta identity:} $\theta'\mathrm{pr}\,e_X(M_t)=e\theta'(M_t)=\mathrm{Im}$. Indeed, for $m\in M_t$ the
vector $e_X m-\mathrm{pr}(e_Xm)$ lies in $\ker\theta'$, so $\theta'(\mathrm{pr}(e_Xm))=\theta'(e_Xm)=e(\theta'm)$.
(c) \emph{Twist chain:} since $T_0\subseteq T_J$, \cref{pf:onto3:idxbasic}(i) gives
$\ontothreeidx_w(Z_g\cap M,e(T_J))\le\ontothreeidx_w(Z_g\cap M,e(T_0))$, so
$\ontothreeidx_w(Z_g\cap M,X_3)\le\ontothreeidx_w(Z_g\cap M,e(T_0))+\eps$.
(d) Put $X_3'=X_3\cap\mathrm{Im}$, a lattice in $Z_g$ (both are lattices in $Z_g$). Since $(X_3)_w\subseteq
\mathrm{Im}_w$ we have $(X_3')_w=(X_3)_w$, hence $\ontothreeidx_w(Z_g\cap M,X_3')=\ontothreeidx_w(Z_g\cap M,X_3)$.
Now \cref{pf:onto3:onlocal}, applied with $T$ (by (a)), $X_3'\subseteq X_4=\mathrm{Im}$ (by (c)) and $T'$ (by (b)),
gives the claim.
\end{proof}

% ---------------------------------------------------------------------------
\subsection{Record objects and the global bound}\label{sec:pf:onto3:record}

\subsubsection{Record objects}

Let $S$ be an indefinite datum over $K$ (with its base, side and index set $J$), $G$ a finite group and
$c=(C,C_{b1},C_{b2})\in\R^3$ a triple of constants; put $C_b=C_{b1}+C_{b2}+4$ and
\begin{equation}\label{eq:pf:onto3:CE}
C_E(c)=4C^2+25C+C_b+44 .
\end{equation}

\begin{definition}[Record objects]\label{pf:onto3:defobj}
A \emph{record object} $O$ over $(S,G,c)$ consists of:
\begin{enumerate}[label=(\alph*)]
\item for each $j\in J$: an integer $n_j\ge0$, the space $Z_j=K^{n_j}$ with a representation $\rho_j$ of $G$, a
subspace $V_0(j)\subseteq Z_j$, an endomorphism $e_j$ of $Z_j$, an $\OO_K$-submodule $M_s(j)\subseteq Z_j$, a
$K$-linear map $\varphi_j\colon K^{d_W(j)}\to Z_j$ and an element $n_K(j)\in K$;
\item a subgroup $U\le G$ and a function $\psi\colon G\to\OO_K$;
\item for each finite place $v$ and $j\in J$: $\OO_K$-submodules $\Theta_1(v,j)$, $\Theta(v,j)$, $\Theta_K(v,j)$,
$\Theta_0(v,j)$, $\Theta_J(v,j)$ and $X_3(v,j)$ of $Z_j$; elements $\sigma(v,j),a(v,j)\in K$; an integer
$E(v,j)\ge0$; real numbers $r_1(v,j),r_2(v,j),\eps_0(v,j)$; a space $X_{v,j}=K^{n_X(v,j)}$ with linear maps
$\theta'_{v,j}\colon X_{v,j}\to Z_j$ and $e^X_{v,j},\mathrm{pr}_{v,j}\in\End(X_{v,j})$, and an $\OO_K$-submodule
$M_t(v,j)\subseteq X_{v,j}$;
\item real numbers $J_p(j,\ell)$, $w_t(j)$, $\mathrm{th}(\ell)$, $\mathrm{ad}(\ell)$ ($j\in J$, $\ell\in\N$) and
$k$, $m$, $f$, $b_p$, $S_c$, $\eta$, $T$, $D_F$, $X_2$, $x$, $\eps$, $v_d$, $r_r$, $d_{e2}$, $N_c$; integers
$p_{\mathrm{de}}$, $T_2$, $\gamma_1$, $\gamma_2$, $q_1$, $d_1\ge0$; and a finite set $s_m$ of natural numbers.
\end{enumerate}
\formallabel{Formalism.L2.Indef.Tree.OnObj}
\end{definition}

We use the following derived quantities, for $j\in J$, finite places $v$ and $\ell\in\N$:
\begin{gather*}
Z_j^\psi=\{z:\rho_j(u)z=\psi(u)z\ \forall u\in U\},\quad Z_g(j)=Z_j^\psi\cap V_0(j),\quad d_0(j)=\dim Z_j^\psi,\\
e_J(v,j)=e_j(\Theta_J(v,j)),\quad \mathrm{Im}(v,j)=e_j\theta'_{v,j}(M_t(v,j)),\quad
\mathrm{I}(v,j)=\theta'_{v,j}\mathrm{pr}_{v,j}e^X_{v,j}(M_t(v,j)),\\
r(v,j)=r_2(v,j)\ \text{if }p_v=p_{\mathrm{de}},\quad r(v,j)=r_1(v,j)\ \text{otherwise},\\
e_0(v,j)=\frac{E(v,j)-v_v(\sigma(v,j))}{e_v},\quad v_U(v)=\frac{v_v(|U|)}{e_v},\quad
n_n(v,j)=\frac{v_v(n_K(j))}{e_v},\\
x_e=x^{\eps},\quad b=e^{x_e},\quad L=\log b=x_e,\quad B_0=fmbL,\quad E_{1u}(O)=x_e(B_0+N_c),\\
S_e=S_c+\eta,\quad Y=C\,D_F,\quad A_1(j)=C\bigl(\log k\cdot m\,w_t(j)+fkS_e\bigr),\quad A_2(j)=Ck\,w_t(j),\\
s_B(j)=w_t(j)+fS_e,\quad s_C(j)=Cm(2+v_d)(w_t(j)+fS_e),\quad \beta=fkb_p,\quad y=\frac{1}{4(fk)^3},\\
R_b(j)=2n_jfmb,\quad g_q(v)=2fr_r^2\ \text{if }p_v\in\{\gamma_1,\gamma_2\},\quad g_q(v)=0\ \text{otherwise},\\
\mathrm{Bad}=\{\gamma_1,\gamma_2,q_1,d_1\}\cup s_m,\qquad E_J=30T_2mbx_e^2+2mb,\\
a_E(j,\ell)=\frac{A_1(j)}{\ell}+A_2(j)\mathrm{th}(\ell)+T\mathrm{ad}(\ell)+[\ell>X_2]\frac{A_2(j)}{\ell^2},
\end{gather*}
where $[\cdot]$ is $1$ if the condition holds and $0$ otherwise. The \emph{error terms} are
$\ontothreeerrG(j,\ell)=2a_E(j,\ell)+2$ if $\ell\le Y$ and $\ontothreeerrG(j,\ell)=0$ otherwise;
\begin{equation}\label{eq:pf:onto3:errP}
\ontothreeerrP(j,\ell)=\begin{cases}
4T+4fr_r^2 & \ell\in\{\gamma_1,\gamma_2\},\\
2s_C(j)+2T+2 & \ell\in s_m\text{ (and not the above)},\\
2(s_B(j)+T+1)+4\beta[\ell=d_1] & \ell\in\{q_1,d_1\}\text{ (and not the above)},\\
\ontothreeerrG(j,\ell) & \text{otherwise};
\end{cases}
\end{equation}
$\ontothreeerrl(v,j)=\ontothreeerrP(j,p_v)$, and
\begin{equation}\label{eq:pf:onto3:errOn}
\ontothreeerrOn(w,j)=\ontothreeerrl(w,j)+\frac{r(w,j)/e_w+r(\bar w,j)/e_{\bar w}}{[G:U]\,D(S)} .
\end{equation}

\begin{definition}[Leaves]\label{pf:onto3:leaves}
Let $O$ be a record object over $(S,G,c)$ and $N$ numbers over the base of $S$. ``At pairs'' means: for every
$w\in\ontothreeTw(S)$, every $j\in J$ and both $v\in\{w,\bar w\}$. The \emph{leaves} of $O$ are the following
conditions.
\begin{enumerate}[label=(L\arabic*),leftmargin=*]
\item\label{pf:onto3:L:mackey} $(U,\psi)$ is a Mackey pair in $G$.
\item\label{pf:onto3:L:sup} Each $M_s(j)$ is a full lattice in $Z_j$ stable under $\rho_j(G)$.
\item\label{pf:onto3:L:proj} $e_j^2=e_j$ and $e_j(Z_j)=V_0(j)$.
\item\label{pf:onto3:L:coord} $\varphi_j$ is injective with image $Z_g(j)$ and $\varphi_j(M_j)=Z_g(j)\cap
M_s(j)$; and at pairs $\varphi_j\bigl(\theta_{j,\mu(w)}T_l(j,\mu(w))\bigr)=\mathrm{I}(v,j)$.
\item\label{pf:onto3:L:span} At pairs, $\Theta(v,j)$ and $\Theta_K(v,j)$ are full lattices.
\item\label{pf:onto3:L:lemS} At pairs, $\Theta(v,j)\subseteq M_s(j)$, $\Theta(v,j)$ is $\rho_j(G)$-stable,
$\sigma(v,j)\ne0$ and $\Theta(v,j)\subseteq\sigma(v,j)\Theta_K(v,j)$.
\item\label{pf:onto3:L:idx} At pairs, $\Theta_1(v,j)$ is a full lattice contained in $\Theta(v,j)$, $a(v,j)\ne0$,
$v_v(a(v,j))=E(v,j)$, and $\ontothreeidx_v(W\cap M_s(j),W\cap\Theta_1(v,j))\le E(v,j)\dim W+r(v,j)$ for every
subspace $W\subseteq Z_j$. Moreover $0\le d_{e2}\le1$, $r_1,r_2\ge0$ everywhere, and for $i=1,2$, every $j$
and every finite set $S'$ of finite places, $\sum_{v\in S'}\ell_vr_i(v,j)\le n_jfmb\,d_{e2}$.
\item\label{pf:onto3:L:data} At pairs, $\Theta_0(v,j)$ is finitely generated and contained in $Z_j^\psi$, and
$|U|\,(Z_j^\psi\cap\Theta_K(v,j))\subseteq\Theta_0(v,j)$.
\item\label{pf:onto3:L:rank} At pairs, $e_j(\Theta_0(v,j))$ and $e_J(v,j)$ are lattices in $Z_g(j)$, and
$\Theta_0(v,j)_v\subseteq\Theta_J(v,j)_v$ (localisations at $v$).
\item\label{pf:onto3:L:ident} At pairs, $\theta'_{v,j}e^X_{v,j}=e_j\theta'_{v,j}$ and
$x-\mathrm{pr}_{v,j}(x)\in\ker\theta'_{v,j}$ for $x\in e^X_{v,j}(X_{v,j})$.
\item\label{pf:onto3:L:tw2} At pairs, $X_3(v,j)$ and $\mathrm{Im}(v,j)$ are lattices in $Z_g(j)$ and
$X_3(v,j)_v\subseteq\mathrm{Im}(v,j)_v$.
\item\label{pf:onto3:L:tw1} At pairs, $\eps_0(v,j)\ge0$ and
$\ontothreeidx_v(Z_g(j)\cap M_s(j),X_3(v,j))\le\ontothreeidx_v(Z_g(j)\cap M_s(j),e_J(v,j))+e_v\eps_0(v,j)$;
and $\sum_j(\eps_0(w,j)+\eps_0(\bar w,j))\le2/p_w^2$ for $w\in\ontothreeTw(S)$.
\item\label{pf:onto3:L:share} For every $j$ there is $e\le y$ with $(1-e)n_j\le|U|\,d_0(j)$.
\item\label{pf:onto3:L:indexU} $(1-y)|U|\le[G:U]$.
\item\label{pf:onto3:L:bad} For every $j$ there is $\delta\in[0,\frac14]$ with $(1-\delta)d_0(j)\le D(S)$.
\item\label{pf:onto3:L:unip} $1\le r_r$, and for $w\in\ontothreeTw(S)$: $v_w(|U|)=e_w\,fr_r(r_r-1)/2$ if
$p_w\in\{\gamma_1,\gamma_2\}$, and $v_w(|U|)=0$ otherwise.
\item\label{pf:onto3:L:good} (Good depth.) $\mathrm{th},\mathrm{ad}\ge0$; $\sum_{p\in Q}\mathrm{th}(p)\log p\le
C/k$ and $T\sum_{p\in Q}\mathrm{ad}(p)\log p\le Cmb$ for every finite set $Q$ of primes; and for
$w\in\ontothreeTw(S)$ with $p_w\notin s_m\cup\{q_1,d_1,\gamma_1,\gamma_2\}$ and all $j$: if $p_w\le Y$, then
$e_0(v,j)\le F_0(w,j)+a_E(j,p_w)+1$ for $v\in\{w,\bar w\}$; if $p_w>Y$, then
$e_0(w,j)=e_0(\bar w,j)=0\le F_0(w,j)$.
\item\label{pf:onto3:L:ram} (Ramified depth.) For $w\in\ontothreeTw(S)$ with $p_w\in\{q_1,d_1\}$ and all $j$:
$e_0(v,j)\le s_B(j)+T+1$ for $v\in\{w,\bar w\}$; $F_0(w,j)=0$ if $p_w=q_1$; $-2\beta\le F_0(w,j)\le0$ if
$p_w=d_1$.
\item\label{pf:onto3:L:small} (Small primes.) For $w\in\ontothreeTw(S)$ with $p_w\in s_m$ and all $j$:
$e_0(v,j)\le s_C(j)+T+1$ for $v\in\{w,\bar w\}$, and $F_0(w,j)=0$.
\item\label{pf:onto3:L:level} (Level primes.) For $w\in\ontothreeTw(S)$ with $p_w\in\{\gamma_1,\gamma_2\}$ and
all $j$: $e_0(v,j)\le2T-n_n(v,j)$ for $v\in\{w,\bar w\}$, and $n_n(w,j)+n_n(\bar w,j)=-2F_0(w,j)$.
\item\label{pf:onto3:L:page} (Page.) $1\le C$, $1\le k$, $1\le m$, $8\le f$, $0\le T,S_c,\eta,b_p,v_d,N_c$;
$\gamma_1,\gamma_2,q_1,d_1$ are primes with $\gamma_1\ne\gamma_2$, $q_1\ne d_1$ and
$\{q_1,d_1\}\cap\{\gamma_1,\gamma_2\}=\emptyset$, none of them lies in $s_m$; and $1\le x$.
\item\label{pf:onto3:L:sizes} (Sizes.) $0\le w_t(j)\le2fb$ for all $j$; $1<Y$, $\log Y\le2L$, $Y\le mb$,
$kS_eL\le mb$, $k\le X_2$, $\log k\le L$ and $1\le L$.
\item\label{pf:onto3:L:x} $\log x\le x_e$ and $T_2\le f$.
\item\label{pf:onto3:L:badc} (Bad constants.) $C_{b1},C_{b2}\ge0$ and for all $j$:
$(2(s_B(j)+T+1)+4\beta)\log(q_1d_1)\le C_{b1}fmb$, $(\sum_{p\in s_m}\log p)(2s_C(j)+2T+2)\le C_{b2}fmbL$ and
$(4T+4fr_r^2)\log(\gamma_1\gamma_2)\le4b$.
\item\label{pf:onto3:L:exp} (Exponents.) $J_p\ge0$; $\sum_{p\in Q}J_p(j,p)\log p\le E_J$ for every $j$ and every
finite set $Q$ of primes; and $J_0(\mu(w),w,j)\le J_p(j,p_w)$ for $w\in\ontothreeTw(S)$ and all $j$.
\item\label{pf:onto3:L:E1} For the constant $C_E$ under consideration, $E_1(j)=C_E\,E_{1u}(O)$ for all $j$.
\end{enumerate}
\end{definition}

\begin{lemma}[Places above a prime]\label{pf:onto3:place3}
For every number field $K$ and every prime $p$ there is a finite place $w$ of $K$ with $p_w=p$; in particular there is a
finite place with residue characteristic $3$.
\formallabel{Formalism.L2.Indef.Tree.exists_place3}
\formallabel{Formalism.L2.Indef.Tree.exists_place_lemma}
\end{lemma}

\begin{proof}
$p\Z$ is a maximal ideal of $\Z$ and $\Z\to\OO_K$ is injective and integral, so by the lying-over theorem there is a
maximal ideal $Q$ of $\OO_K$ with $Q\cap\Z=p\Z$. It is non-zero (it contains $p$), and the residue field $\OO_K/Q$ has
characteristic $p$ since $p\in Q$.
\end{proof}


The leaves \ref{pf:onto3:L:good}--\ref{pf:onto3:L:level}, \ref{pf:onto3:L:exp} and \ref{pf:onto3:L:E1} involve the
numbers $N$ (through $F_0$, $J_0$ and $E_1$). The leaves \ref{pf:onto3:L:idx}, \ref{pf:onto3:L:tw1} and the first
two parts of \ref{pf:onto3:L:exp} are also used in \emph{existential form}: there exist $(\Theta_1,a,r_1,r_2,d_{e2})$,
resp.\ $\eps_0$, resp.\ $J_p$, such that the record object with these data replaced satisfies the leaf. For
\ref{pf:onto3:L:exp} the existential form includes in addition $J_0(\mu(w),w,j)\ge0$ for $w\in\ontothreeTw(S)$; we
call it the \emph{exponent bound}.

\subsubsection{Numerical lemmas}

\begin{lemma}\label{pf:onto3:numerics}
\begin{enumerate}[label=(\roman*)]
\item (Shares.) Let $U,\delta>0$, $n,d_0\ge0$, $0\le e\le y\le\frac18$ and $0\le\delta'\le\frac14$ be real numbers
with $(1-e)n\le Ud_0$, $(1-y)U\le\delta$ and $(1-\delta')d_0\le d$. Then $n/\delta\le2d$.
\item (Depth cases.) Let $e,e',F\in\R$. If $e,e'\le F+a+1$, then $e+e'-2F\le2a+2$. If $e=e'=0\le F$, then
$e+e'-2F\le0$. If $e,e'\le s+T+1$ and $-2\beta\le F\le0$, then $e+e'-2F\le2(s+T+1)+4\beta$. If $e,e'\le s+T+1$ and
$F=0$, then $e+e'-2F\le2s+2T+2$. If $e\le2T-n$, $e'\le2T-n'$ and $n+n'=-2F$, then $e+e'-2F+2g\le4T+2g$.
\item (Budget.) If $\delta,D>0$, $B\ge0$ and $n/\delta\le3D$, then $4nB/(\delta D)\le12B$.
\end{enumerate}
\formallabel{Formalism.EigenvalueLevelCeiling.share_numeric}
\formallabel{Formalism.EigenvalueLevelCeiling.depth_cases}
\formallabel{Formalism.EigenvalueLevelCeiling.budget_numeric}
\end{lemma}

\begin{proof}
(i) The hypotheses give $\frac78n\le Ud_0$, $\frac78U\le\delta$ and $\frac34d_0\le d$, so
$n\le\frac87\cdot\frac43\cdot\frac87\,\delta d=\frac{256}{147}\delta d\le2\delta d$. (ii) and (iii) are immediate.
\end{proof}

\begin{lemma}[Prime sums]\label{pf:onto3:primes}
Let $Y>1$, $X_2>0$ and let $Q$ be a finite set of primes. Then $\sum_{p\in Q,\,p\le Y}\log p/p\le\log Y$,
$\sum_{p\in Q,\,p\le Y}\log p\le Y\log4$, and $\sum_{p\in Q,\,p>X_2}\log p/p^2\le2.04/X_2$.
\formallabel{Formalism.EigenvalueLevelCeiling.primeBoundsExact}
\formallabel{Formalism.EigenvalueLevelCeiling.chebyshev_theta_log_four}
\end{lemma}

\begin{proof}
The first bound is \lref{sum_log_p_over_p}, which states it in this form (compare Mertens's estimate
\cite{Mertens1874}). The second follows from Chebyshev's bound $\sum_{p\le Y}\log p\le Y\log 4$, a consequence of
$\prod_{p\le n}p\le4^n$ \cite[Ch.~2]{MontgomeryVaughan2007}, which is available in Mathlib. The third is
\lref{prime_log_over_sq_tail}, which states it in this form.
\end{proof}

\begin{lemma}[Sums over places]\label{pf:onto3:placesum}
Let $S$ be a finite set of finite places of $K$.
\begin{enumerate}[label=(\roman*)]
\item If $X\colon\N\to\R_{\ge0}$, then $\sum_{w\in S}\ell_we_wX(p_w)\le\sum_{p\in p(S)}\log p\cdot X(p)$, where
$p(S)$ is the set of residue characteristics of the places of $S$.
\item Let $Y>1$, $X_2>0$, $A_1,A_2,k_w,T\ge0$ and $\mathrm{th},\mathrm{ad}\ge0$ with
$\sum_{p\in Q}\mathrm{th}(p)\log p\le C_{\mathrm{th}}$ and $T\sum_{p\in Q}\mathrm{ad}(p)\log p\le C_{\mathrm{ad}}$
for all finite sets $Q$ of primes, and put $g(\ell)=2\bigl(A_1/\ell+k_w\mathrm{th}(\ell)+T\mathrm{ad}(\ell)+
[\ell>X_2]A_2/\ell^2\bigr)+2$ for $\ell\le Y$ and $g(\ell)=0$ otherwise. Then for every finite set $Q$ of primes,
$\sum_{p\in Q}\log p\cdot g(p)\le2\bigl(A_1\log Y+k_wC_{\mathrm{th}}+C_{\mathrm{ad}}+2.04A_2/X_2+Y\log4\bigr)$.
\item In the situation of (ii), let $\mathrm{Bad}$ be a finite set, $\mathrm{eb}\colon\N\to\R_{\ge0}$, and let
$\ontothreeerrl$ be a function on $S$ with $\ontothreeerrl(w)\le\mathrm{eb}(p_w)$ if $p_w\in\mathrm{Bad}$ and
$\ontothreeerrl(w)\le g(p_w)$ otherwise. Then
$\sum_{w\in S}\ell_we_w\ontothreeerrl(w)\le2\bigl(A_1\log Y+k_wC_{\mathrm{th}}+C_{\mathrm{ad}}+2.04A_2/X_2+
Y\log4\bigr)+\sum_{p\in\mathrm{Bad}}\log p\cdot\mathrm{eb}(p)$.
\end{enumerate}
\formallabel{Formalism.EigenvalueLevelCeiling.place_sum}
\formallabel{Formalism.EigenvalueLevelCeiling.error_sum_good_places_exact}
\formallabel{Formalism.EigenvalueLevelCeiling.local_sum_exact}
\end{lemma}

\begin{proof}
(i) Group the places by residue characteristic. For a prime $p$, $\ell_we_w=e_wf_w\log p/[K:Q]$ for $w\mid p$,
and $\sum_{w\mid p}e_wf_w=[K:\Q]$; hence the places of $S$ above $p$ contribute at most $\log p\cdot X(p)$.
(ii) Expand $\log p\cdot g(p)$ and apply \cref{pf:onto3:primes} to the primes of $Q$ which are $\le Y$, resp.\
$\le Y$ and $>X_2$: the four sums are at most $A_1\log Y$, $k_wC_{\mathrm{th}}$, $C_{\mathrm{ad}}$ and
$2.04A_2/X_2$, and $\sum_{p\le Y}2\log p\le2Y\log4$. (iii) Since $g\ge0$, (i) applied to the function equal to
$\mathrm{eb}$ on $\mathrm{Bad}$ and to $g$ elsewhere bounds the sum by
$\sum_{p\in p(S)\cap\mathrm{Bad}}\log p\,\mathrm{eb}(p)+\sum_{p\in p(S)\smallsetminus\mathrm{Bad}}\log p\,g(p)$,
and (ii) applies to the second sum.
\end{proof}

\begin{lemma}[Absorption]\label{pf:onto3:absorb}
Let $C\ge1$, $C_b\ge0$, $1\le L\le x_e$, $f,m\ge1$, $b>0$, $k\ge1$ with $\log k\le L$, $0\le w\le2fb$, $S_e\ge0$
with $kS_eL\le mb$, $Y>1$ with $\log Y\le2L$ and $Y\le mb$, $k\le X_2$, and $\mathrm{bad}\le C_bfmbL$. Then
\begin{multline*}
2\Bigl(C(\log k\cdot mw+fkS_e)\log Y+Ckw\frac Ck+Cmb+2.04\frac{Ckw}{X_2}+Y\log4\Bigr)+\mathrm{bad}+12fmb\\
\le(4C^2+25C+C_b+12)\,x_e\,fmbL .
\end{multline*}
\formallabel{Formalism.EigenvalueLevelCeiling.absorb_local_conditions_exact}
\end{lemma}

\begin{proof}
Bound the terms in the bracket separately, using $\log Y\le2L$: $C\log k\cdot mw\log Y\le2CL\log k\cdot mw$;
$CfkS_e\log Y\le2Cf(kS_eL)\le2Cfmb$; $Ckw\cdot C/k=C^2w\le2C^2fbm$; $Cmb\le Cfmb$;
$2.04Ckw/X_2\le2.04Cw\le4.08Cfbm$; $Y\log4\le1.3863\,mb\le1.3863\,Cfmb$. Hence twice the bracket is at most
$4CL\log k\cdot mw+(4C^2+17C)fmb$. Further $4CL\log k\cdot mw\le4CL^2m\cdot2fb\le8C\,x_e\,fmbL$,
$(4C^2+17C)fmb\le(4C^2+17C)x_efmbL$ and $12fmb\le12x_efmbL$ (as $x_eL\ge1$), and $\mathrm{bad}\le C_bx_efmbL$.
\end{proof}

\begin{lemma}\label{pf:onto3:badsum}
Let $O$ satisfy \ref{pf:onto3:L:good} (non-negativity of $\mathrm{th},\mathrm{ad}$), \ref{pf:onto3:L:page},
\ref{pf:onto3:L:sizes}, \ref{pf:onto3:L:x} and \ref{pf:onto3:L:badc}. Then $0\le y\le\frac18$,
$\ontothreeerrP(j,\ell)\ge0$ for all $j,\ell$, and $\sum_{p\in\mathrm{Bad}}\log p\cdot\ontothreeerrP(j,p)\le
C_bfmbL$ for all $j$.
\formallabel{Formalism.L2.Indef.Tree.bad_sum}
\formallabel{Formalism.L2.Indef.Tree.errorTerm_nonneg}
\formallabel{Formalism.L2.Indef.Tree.y_bounds}
\end{lemma}

\begin{proof}
$fk\ge8$ gives $y\le1/32$. All terms in \eqref{eq:pf:onto3:errP} are non-negative by \ref{pf:onto3:L:page},
\ref{pf:onto3:L:sizes}, $L=x_e$ and the non-negativity of $\mathrm{th},\mathrm{ad}$. By \ref{pf:onto3:L:page} the
primes $\gamma_1,\gamma_2,q_1,d_1$ are distinct and not in $s_m$, so $\ontothreeerrP(j,\gamma_i)=4T+4fr_r^2$,
$\ontothreeerrP(j,q_1)=2(s_B(j)+T+1)$, $\ontothreeerrP(j,d_1)=2(s_B(j)+T+1)+4\beta$ and
$\ontothreeerrP(j,p)=2s_C(j)+2T+2$ for $p\in s_m$. By \ref{pf:onto3:L:badc} the four groups contribute at most
$4b$, $C_{b1}fmb$ (using $\beta\log q_1\ge0$) and $C_{b2}fmbL$; finally $4b\le4fmbL$ and $C_{b1}fmb\le
C_{b1}fmbL$ since $f,m,L\ge1$.
\end{proof}

\begin{lemma}[Global form]\label{pf:onto3:globalform}
Let $D>0$ be an integer, $S$ a finite set of finite places, $\delta>0$, and let, for $w\in S$, $v\in\{w,\bar w\}$
and $j\in J$, real numbers $I_j(v)$, $e_0(v,j)$, $r_0(v,j)$, $\eps_0(v,j)$, $v_U(v)$, $d_K(j)$, $F_0(w,j)$,
$\ontothreeerrl(w,j)$, $g_q(w)$ and $J_x(w,j)$ be given, together with $E_A(j),R_b(j),E_J(j)$, such that for all
$w\in S$, $v\in\{w,\bar w\}$ and $j$:
\begin{enumerate}[label=(\alph*)]
\item $I_j(v)\le e_v\bigl((e_0(v,j)+v_U(v))D+(r_0(v,j)+d_K(j)v_U(v))/\delta+\eps_0(v,j)\bigr)$;
\item $d_K(j)/\delta\le3D$; $0\le v_U(w)=v_U(\bar w)$ and $4v_U(w)\le g_q(w)$;
\item $e_0(w,j)+e_0(\bar w,j)-2F_0(w,j)+2g_q(w)\le\ontothreeerrl(w,j)$;
\item $\sum_j(\eps_0(w,j)+\eps_0(\bar w,j))\le2/p_w^2$;
\item $\sum_{w\in S}\ell_we_w\ontothreeerrl(w,j)\le E_A(j)$, $\sum_{w\in S'}\ell_we_wr_0(w,j)\le R_b(j)$ for every
finite set $S'$, and $\sum_{w\in S}\ell_wJ_x(w,j)\le E_J(j)$.
\end{enumerate}
Then, with $\ontothreeerr(w,j)=\ontothreeerrl(w,j)+(r_0(w,j)+r_0(\bar w,j))/(\delta D)$, for all $w\in S$
\[
\sum_j\bigl(I_j(w)+I_j(\bar w)\bigr)\le\Bigl(\sum_j\bigl(2e_wF_0(w,j)+e_w\ontothreeerr(w,j)\bigr)+
e_w\frac{2}{p_w^2}\Bigr)D,
\]
and $\sum_{w\in S}\ell_w\bigl(e_w\ontothreeerr(w,j)+J_x(w,j)\bigr)\le E_A(j)+2R_b(j)/(\delta D)+E_J(j)$ for all $j$.
\formallabel{Formalism.EigenvalueLevelCeiling.local_conditions_global_form}
\formallabel{Formalism.EigenvalueLevelCeiling.onExport2_of}
\end{lemma}

\begin{proof}
Fix $w\in S$ and $j$; recall $e_{\bar w}=e_w$. Adding (a) for $v=w,\bar w$ and using (b),
\[
I_j(w)+I_j(\bar w)\le e_w\Bigl(\bigl(e_0(w,j)+e_0(\bar w,j)+8v_U(w)\bigr)D+\frac{r_0(w,j)+r_0(\bar w,j)}{\delta}
+\eps_0(w,j)+\eps_0(\bar w,j)\Bigr),
\]
and $e_0(w,j)+e_0(\bar w,j)+8v_U(w)\le e_0(w,j)+e_0(\bar w,j)+2g_q(w)\le2F_0(w,j)+\ontothreeerrl(w,j)$ by (b), (c).
Summing over $j$ and using (d) together with $e_w\cdot2/p_w^2\le e_w\cdot2/p_w^2\cdot D$ gives the first claim. For
the second, the map $w\mapsto\bar w$ is injective and preserves $\ell_w$ and $e_w$, so
$\sum_{w\in S}\ell_we_wr_0(\bar w,j)=\sum_{w\in\bar S}\ell_we_wr_0(w,j)\le R_b(j)$; together with (e) this gives
the bound.
\end{proof}

\subsubsection{The global bound}

\begin{theorem}[Global bound for record objects]\label{pf:onto3:oncore}
Let $S$ be an indefinite datum with rank rows, $N$ numbers over its base, $O$ a record object over $(S,G,c)$
satisfying the leaves \ref{pf:onto3:L:mackey}--\ref{pf:onto3:L:E1} for $N$ and a constant $C_E\ge C_E(c)$. Then
$\ontothreeONw(S,N,\ontothreeerrOn)$ holds for the error function \eqref{eq:pf:onto3:errOn} of $O$. Consequently, if
$O$ satisfies the leaves with \ref{pf:onto3:L:idx}, \ref{pf:onto3:L:tw1} and the first two parts of
\ref{pf:onto3:L:exp} in existential form, then $\ontothreeONw(S,N,\ontothreeerr)$ holds for some error function
$\ontothreeerr$.
\formallabel{Formalism.L2.Indef.Tree.on_core}
\formallabel{GlobalTheory.TraceClassRank.on_core_atPlaces}
\formallabel{GlobalTheory.OnAtPlacesBase.on_of_has_atPlaces}
\formallabel{Formalism.L2.Indef.Tree.local_bound}
\end{theorem}

\begin{proof}
\emph{Reduction.} Replacing $\Theta_J(v,j)$ by $\Theta_J(v,j)+\Theta_0(v,j)$ for all $v,j$ does not change
$\Theta_J(v,j)_v$ (by \ref{pf:onto3:L:rank}), hence does not change the indices in \ref{pf:onto3:L:tw1}, and
$e_j(\Theta_J+\Theta_0)=e_J+e_j(\Theta_0)$ is again a lattice in $Z_g(j)$; after the replacement,
$\Theta_0(v,j)\subseteq\Theta_J(v,j)$. The existential form follows by applying the first assertion to the
record object with the data provided by the existential leaves. So assume $\Theta_0\subseteq\Theta_J$ at pairs.

\emph{Step 1 (local bound).} Let $w\in\ontothreeTw(S)$, $j\in J$, $v\in\{w,\bar w\}$ and $t=\mu(w)$. We apply
\cref{pf:onto3:lcw} at the place $v$ to $Z=Z_j$, $\rho=\rho_j$, $(U,\psi)$ \ref{pf:onto3:L:mackey}, $e=e_j$
\ref{pf:onto3:L:proj}, $M=M_s(j)$ \ref{pf:onto3:L:sup}, $T_1,T,T_K,T_0,T_J=\Theta_1,\Theta,\Theta_K,\Theta_0,
\Theta_J$ at $(v,j)$ (\ref{pf:onto3:L:span}--\ref{pf:onto3:L:rank}), $\sigma=\sigma(v,j)$, $a=a(v,j)$,
$E=E(v,j)$, $r=r(v,j)$ \ref{pf:onto3:L:idx}, $X=X_{v,j}$ with $\theta',e^X,\mathrm{pr},M_t$
(\ref{pf:onto3:L:ident}, \ref{pf:onto3:L:tw2}), $X_3=X_3(v,j)$ and $\eps=e_v\eps_0(v,j)$ \ref{pf:onto3:L:tw1}, and
$\varphi=\varphi_j$, $M'=M_j$, $T'=\theta_{j,t}T_l(j,t)$ \ref{pf:onto3:L:coord}. Since $d_W(j)=D=D(S)$, dividing
by $e_v$ gives
\begin{equation}\label{eq:pf:onto3:loc}
I_j(t,v)\le e_v\Bigl((e_0(v,j)+v_U(v))D+\frac{r(v,j)/e_v+n_jv_U(v)}{[G:U]}+\eps_0(v,j)\Bigr).
\end{equation}

\emph{Step 2 (shares).} By \ref{pf:onto3:L:share}, \ref{pf:onto3:L:indexU}, \ref{pf:onto3:L:bad},
\cref{pf:onto3:badsum} ($y\le\frac18$) and \cref{pf:onto3:numerics}(i) with $U=|U|$, $\delta=[G:U]$, $d=D$:
$n_j/[G:U]\le2D\le3D$.

\emph{Step 3 (unipotent valuations).} For $w\in\ontothreeTw(S)$, $v_{\bar w}(|U|)=v_w(|U|)$ and $e_{\bar
w}=e_w$, so $v_U(\bar w)=v_U(w)\ge0$; by \ref{pf:onto3:L:unip}, $v_U(w)=fr_r(r_r-1)/2$ or $0$, so $4v_U(w)\le
2fr_r^2=g_q(w)$ at the level primes and $4v_U(w)=0=g_q(w)$ otherwise.

\emph{Step 4 (depth).} For $w\in\ontothreeTw(S)$ and $j$,
$e_0(w,j)+e_0(\bar w,j)-2F_0(w,j)+2g_q(w)\le\ontothreeerrl(w,j)$. This follows from
\cref{pf:onto3:numerics}(ii), according to the residue characteristic $p=p_w$: for $p\in\{\gamma_1,\gamma_2\}$
from \ref{pf:onto3:L:level} (with $g=g_q(w)=2fr_r^2$, giving $4T+4fr_r^2$); for $p\in s_m$ from
\ref{pf:onto3:L:small} ($g_q=0$); for $p\in\{q_1,d_1\}$ from \ref{pf:onto3:L:ram} (with $\beta$ replaced by $0$
if $p=q_1$); otherwise from \ref{pf:onto3:L:good}, where for $p\le Y$ the bound $2a_E(j,p)+2$ is
$\ontothreeerrG(j,p)$ and for $p>Y$ the bound is $0=\ontothreeerrG(j,p)$.

\emph{Step 5 (sums).} By \cref{pf:onto3:placesum}(iii) with $\mathrm{eb}=\ontothreeerrP(j,\cdot)$ (non-negative
by \cref{pf:onto3:badsum}), $A_1=A_1(j)$, $A_2=k_w=A_2(j)$, $C_{\mathrm{th}}=C/k$, $C_{\mathrm{ad}}=Cmb$
\ref{pf:onto3:L:good}:
\begin{multline*}
\sum_{w\in\ontothreeTw(S)}\ell_we_w\ontothreeerrl(w,j)\le E_A(j):=2\Bigl(A_1(j)\log Y+A_2(j)\frac Ck+Cmb+
2.04\frac{A_2(j)}{X_2}+Y\log4\Bigr)\\
+\sum_{p\in\mathrm{Bad}}\log p\cdot\ontothreeerrP(j,p).
\end{multline*}
For every finite set $S'$, $\sum_{v\in S'}\ell_vr(v,j)\le\sum_{v\in S'}\ell_v(r_1(v,j)+r_2(v,j))\le
2n_jfmb\,d_{e2}\le R_b(j)$ by \ref{pf:onto3:L:idx}. By \ref{pf:onto3:L:exp} and \cref{pf:onto3:placesum}(i),
$\sum_{w\in\ontothreeTw(S)}\ell_we_wJ_0(\mu(w),w,j)\le\sum_{w}\ell_we_wJ_p(j,p_w)\le E_J$.

\emph{Step 6 (global form).} \Cref{pf:onto3:globalform}, applied with $\delta=[G:U]$, $r_0(v,j)=r(v,j)/e_v$,
$d_K(j)=n_j$, $J_x(w,j)=e_wJ_0(\mu(w),w,j)$, $E_J(j)=E_J$ and the bounds of Steps 1--5, gives (ON1) for
$\ontothreeerrOn$, and
\[
\sum_{w\in\ontothreeTw(S)}\ell_w\bigl(e_w\ontothreeerrOn(w,j)+e_wJ_0(\mu(w),w,j)\bigr)\le
E_A(j)+\frac{2R_b(j)}{[G:U]D}+E_J .
\]
(ON2) holds since $\ontothreeerrP\ge0$ (\cref{pf:onto3:badsum}) and $r\ge0$.

\emph{Step 7 (absorption).} Put $B=fmbL$. By Step 2 and \cref{pf:onto3:numerics}(iii),
$2R_b(j)/([G:U]D)\le12fmb$. By \cref{pf:onto3:badsum} the sum over $\mathrm{Bad}$ in $E_A(j)$ is at most $C_bB$,
and \cref{pf:onto3:absorb} (with $w=w_t(j)$ and the conditions \ref{pf:onto3:L:page}, \ref{pf:onto3:L:sizes},
$L=x_e$) gives $E_A(j)+12fmb\le(4C^2+25C+C_b+12)x_eB$. Since $T_2\le f$ \ref{pf:onto3:L:x}, $x_e=L\ge1$ and
$f\ge1$, $E_J=30T_2mbx_e^2+2mb\le32x_eB$. Hence the right-hand side of Step 6 is at most
$(4C^2+25C+C_b+44)x_eB\le C_Ex_e(B+N_c)=C_E\,E_{1u}(O)=E_1(j)$, using $C_E\ge C_E(c)\ge0$, $N_c\ge0$ and
\ref{pf:onto3:L:E1}. This is (ON3).
\end{proof}

\subsubsection{The canonical error}

\begin{lemma}\label{pf:onto3:canonical}
Let $T$ be a finite set, $J$ a finite set and $I,F,\ontothreeerr,\ontothreeerr_0,J_0$ real functions on $T\times J$,
$e,c,\ell$ real functions on $T$ with $e>0$ and $\ell\ge0$, and $D>0$, $E,E'$ real with $2E\le E'$. Assume
$J_0\ge0$ and
\[
\ontothreeerr_0(w,j)=\frac{\max\bigl(0,\ \frac1D\sum_{j'}I(w,j')-\sum_{j'}2e(w)F(w,j')-c(w)\bigr)}{e(w)|J|}
\qquad(w\in T,\ j\in J).
\]
If $\ontothreeerr$ satisfies (1) $\sum_jI(w,j)\le(\sum_j(2e(w)F(w,j)+e(w)\ontothreeerr(w,j))+c(w))D$ for $w\in T$,
(2) $\ontothreeerr\ge0$ and (3) $\sum_{w\in T}\ell(w)(e(w)\ontothreeerr(w,j)+e(w)J_0(w,j))\le E$ for all $j$, then
$\ontothreeerr_0$ satisfies (1), (2), and (3) with $E'$ in place of $E$.
\formallabel{GlobalTheory.BigOn.canonical_error_real}
\end{lemma}

\begin{proof}
Write $X_w=\frac1D\sum_{j'}I(w,j')-\sum_{j'}2e(w)F(w,j')-c(w)$, so that $\sum_je(w)\ontothreeerr_0(w,j)=\max(0,X_w)$
(if $J=\emptyset$ all claims are trivial or follow from the hypotheses). (1) for $\ontothreeerr_0$ is $X_w\le
\max(0,X_w)$ multiplied by $D$. (2) is clear. For (3), fix $j$. By (1) and (2) for $\ontothreeerr$,
$\max(0,X_w)\le\sum_{j'}e(w)\ontothreeerr(w,j')$, so
\[
|J|\sum_{w\in T}\ell(w)e(w)\ontothreeerr_0(w,j)=\sum_{w\in T}\ell(w)\max(0,X_w)\le\sum_{j'}\sum_{w\in T}
\ell(w)e(w)\ontothreeerr(w,j')\le|J|\,E ,
\]
where the last step uses (3) for each $j'$ and $J_0\ge0$. Also $\sum_w\ell(w)e(w)J_0(w,j)\le E$ by (3) and (2).
Adding, the left-hand side of (3) for $\ontothreeerr_0$ is at most $2E\le E'$.
\end{proof}

\begin{corollary}\label{pf:onto3:onwith}
Let $S$ be an indefinite datum with $D(S)>0$ whose error function is the canonical error
\eqref{eq:pf:onto3:erreq} with respect to a scale $Z$, and assume $E_{1u}(Z)\ge0$ and $J_0(\mu(w),w,j)\ge0$ for
$w\in\ontothreeTw(S)$. If $\ontothreeONw(S,N_Z(C,C_H),\ontothreeerr)$ holds for some function
$\ontothreeerr$, then $S$ satisfies $\ontothreeON(N_Z(C',C_H'))$ for all $C'\ge2C$ and all $C_H'$.
\formallabel{GlobalTheory.BigOn.on_of_onWith_canonical}
\end{corollary}

\begin{proof}
Apply \cref{pf:onto3:canonical} with $T=\ontothreeTw(S)$, $I(w,j)=I_j(\mu(w),w)+I_j(\mu(w),\bar w)$, $F=F_0$,
$e(w)=e_w>0$, $c(w)=e_w\cdot2/p_w^2$, $\ell(w)=\ell_w\ge0$, $D=D(S)$, $E=C\,E_{1u}(Z)$ and $E'=C'E_{1u}(Z)$ (so
$2E\le E'$ as $E_{1u}(Z)\ge0$). The numbers $F_0,J_0$ do not depend on $(C,C_H)$, and $E_1(j)=C\,E_{1u}(Z)$.
\end{proof}

% ---------------------------------------------------------------------------
\subsection{Records of runs and ties}\label{sec:pf:onto3:ties}

This subsection contains the definitions which relate a record object of a point to the runs of the core of the
point: the record algebras of a top structure, the scalars of a run, the form, trace and depth ties, pairs of runs
and their conditions, the record conditions, and the closed statements over points and record objects. We use the
following notions from \cref{sec:pf:onto1:setting,sec:pf:onto1:record}: a datum $D$ with ranks $f,r,m,k,t,n_V$, host
field $E$ (a CM field with conjugation $a\mapsto\bar a$), host primes $g_0,g_1$, integers $q,d,Q_S$, the auxiliary
primes $p_{\mathrm{de}},p'_{\mathrm{de}}$, the finite set $s_m$ of small primes of the host and the coefficient field
$\mathit{LL}$; size data $L$ with scale $x$, exponent $\eps$, $\mathrm L=x^{\eps}$ and $b=e^{\mathrm L}$; weights $w$
with entries $a^\pm_\tau(i)$, integer $n_w$, degree $\ontothreeDF(w)=\sum_{\tau,i}(a^+_\tau(i)+a^-_\tau(i))$ and size
$\ontothreesize(w)$; form providers $B$ with coefficient field $K_0$, spaces $\mathcal F_w$ of forms and lattices
$\mathcal M_w$ of integral forms; valued form algebras $X$ with maps $\iota_w$, theta forms $\mathcal T(w)$ and short
theta forms $\mathcal T^{\mathrm{sh}}(w)$; scalar quadruples and theta lattices $\Theta(S,\kappa,w_L,w_F,T)$; and top
structures $T$ with laws $h$ and their form providers $B_T$. For an index $j$ in range we write $\delta_j$ for the
integer $\delta$ of \cref{pf:onto1:sizeq}, and $\mathcal X_j$ for the moduli space of $T$ at $j$ with coefficient
datum $c$, coefficient field $c.F=K_0$ and ring of rules $R_j$; the forms of weight $w$ of $B_T$ are the rules in a
$K_0$-subspace $C(j,w)\subseteq R_j$, and $\mathcal M_w$ consists of the forms with non-negative valuation at all the
valuations of the core. For $w=(a^\pm_\tau,n)$ put $D_+(w)=\sum_{\tau,i}a^+_\tau(i)$, $D_-(w)=\sum_{\tau,i}a^-_\tau(i)$,
so that $\ontothreeDF(w)=D_+(w)+D_-(w)$.

\subsubsection{Record algebras}

\begin{definition}[Record algebras]\label{pf:onto3:recpj}
Let $T$ be a top structure with laws $h$ and $j$ in range.
\begin{enumerate}[label=(\roman*)]
\item The \emph{record algebra} $X_{T,j}$ is the valued form algebra over $B_T$ at $j$ whose algebra of forms is the
monoid algebra of $R_j$ over the monoid of degrees $(a^+,a^-,n)$, with $\iota_w(F)$ the rule $F$ placed in degree
$(a^\pm_\tau,n_w)$ of $w$; whose points at a place $\lambda$ of $K_0$ are the integral points of $T$ at $j$ over
$\lambda$, a point $y$ acting on a form $\sum_mF_m$ by $v_{\lambda,y}(\sum_mF_m)=\inf_mv_y(F_m)$, where $v_y$ is the
order of rules at $y$; whose theta forms and short theta forms of weight $w$ are the intersections with $C(j,w)$ of
the $\OO_{K_0}$-spans of the theta rules and of the short theta rules of $T$ at $(j,w)$; and whose chart data are those
of $T$ at $j$.
The integrality condition holds by the points condition of the laws.
\item The \emph{enlargement by all valued points} $E_{\mathrm{all}}$ of $X_{T,j}$ has as points at $\lambda$ the
valuations $f\mapsto\nu_e(f(y))$, where $\mathcal O$ is a complete discrete valuation ring with an embedding of
$K_0$ into its fraction field lying over $\lambda$ with ramification index $e$, $y$ is an integral point of
$\mathcal X_j$ over $\mathcal O$, and $\nu_e$ is the valuation of $\mathcal O$ divided by $e$; the class condition
imposed on $y$ in the formal definition is the true proposition. Such a point acts on forms as in (i).
\item $X^+_{T,j}=X_{T,j}\oplus E_{\mathrm{all}}$ (\cref{pf:onto1:record-def}), and for a second norm $n$ of $B_T$,
$X^{+,n}_{T,j}$ is $X^+_{T,j}$ with the $L^2$-norms replaced by those of $n$.
\end{enumerate}
\formallabel{Formalism.L2.Pj.recPJ}
\formallabel{Formalism.L2.Pj.extAll}
\formallabel{Formalism.L2.Pj.recPlus}
\formallabel{GlobalTheory.RunIn.recPlusN}
\end{definition}

\subsubsection{Scalars of a run}

Let $\mathrm{EX}$ be an error datum: elements $\pi_0,\pi_1\in\OO_E$ (with further data not used here) such that
$\pi_i\bar\pi_i=g_i$. Let $B$ be a form provider whose ring $\OO_{K_0}$ is identified with the ring $\OO$ of the
datum, and assume that the map $\OO\to\mathit{LL}$ of the datum is injective; this gives a ring homomorphism
$\iota_E\colon\OO_E\to\OO_{K_0}$. For a weight $w$ and $a,b,b'\in\N$ put
\[
\mathrm{nn}(w)=\prod_{i=0,1}\iota_E(\pi_i)^{D_+(w)}\iota_E(\bar\pi_i)^{D_-(w)},\qquad
\mathrm{pp}(a,b,b')=Q_S^a\,q^b\,d^{b'},\qquad c_{1/2}(w)=\bigl\lceil\ontothreeDF(w)/2\bigr\rceil .
\]

\begin{definition}[Scalars]\label{pf:onto3:scal}
For a weight $w$ put $s(w)=\mathrm{nn}(w)\,\mathrm{pp}(\ontothreeDF(w),c_{1/2}(w),c_{1/2}(w))\in\OO_{K_0}$. For a prime
$p$, weights $\kappa,w_L,w_F$, $T\in\N$ and $\delta=\delta_j$, the scalar quadruple $S_{\mathrm{rec}}=\mathrm{scal}(B,
\mathrm{EX},p,\kappa,w_L,w_F,T)$ is
\begin{gather*}
s_W=s(w_L),\qquad s_F=z_p(w_F)\,s(w_F),\\
s_0=z_p(w_0)\,\mathrm{nn}(w_L)^{\delta}\,\mathrm{nn}(w_F)\,\pi_\delta,\qquad
s_K=z_p(\kappa)\,\mathrm{nn}(w_L)^{T}\,\mathrm{nn}(w_F)\,\pi_T,
\end{gather*}
where $w_0=w_F+\delta w_L$, $z_p(w)=p^{\ontothreeDF(w)+2frn_w}$ and, for $a\in\N$,
$\pi_a=\mathrm{pp}\bigl(a\ontothreeDF(w_L)+\ontothreeDF(w_F),\,c',c'\bigr)$ with $c'=ac_{1/2}(w_L)+c_{1/2}(w_F)$; and the
\emph{record scalar} is $\sigma_{\mathrm{rec}}(\kappa,w_L,w_F,T)=\mathrm{nn}(\kappa)\,(g_0g_1)^{frn_{w_F}}\,\pi_T$.
In particular the first entry of $S_{\mathrm{rec}}$ is $s(w_L)$; it depends on $w_L$ only.
\formallabel{Formalism.MainNotions.DIV.Cons.scalOfP}
\formallabel{GlobalTheory.RunIn.sOf}
\formallabel{GlobalTheory.RunIn.sW_eq_sOf}
\formallabel{GlobalTheory.FormTies.sigRec}
\end{definition}

\begin{lemma}[Size of the first scalar]\label{pf:onto3:sof}
For every weight $w$ we have $s(w)\ne0$, and for every embedding $\sigma\colon K_0\to\C$:
\begin{enumerate}[label=(\roman*)]
\item $\mathrm{nn}(w)\ne0$ and $\log|\sigma(\mathrm{nn}(w))|=\frac12\ontothreeDF(w)\log(g_0g_1)$;
\item $\mathrm{pp}(a,b,b')\ge1$ and $\log\mathrm{pp}(a,b,b')=a\log Q_S+b\log q+b'\log d$; moreover
$\log\mathrm{pp}\bigl(\ontothreeDF(w),c_{1/2}(w),c_{1/2}(w)\bigr)\le(\ontothreeDF(w)+1)\log(f^2qd)$;
\item $0\le\log|\sigma(s(w))|$, $\frac12\ontothreeDF(w)\log\min(g_0,g_1)\le\log|\sigma(s(w))|$,
$\frac12\ontothreeDF(w)\log(g_0g_1)\le\log|\sigma(s(w))|$ and $\log|\sigma(s(w))|\le\lambda^*\ontothreesize(w)$,
\end{enumerate}
where $\lambda^*=\log(g_0g_1)+\log(f^2qd)$ (\cref{pf:onto1:sizeq}).
\formallabel{GlobalTheory.RunIn.sOf_ne_zero_and_log_bounds}
\formallabel{GlobalTheory.log_norm_normaliser}
\formallabel{GlobalTheory.T.log_norm_normaliser_eq_half_degree_log}
\formallabel{GlobalTheory.RunIn.pPow_size}
\formallabel{GlobalTheory.T.ramifiedPrimePowers_size_bounds}
\end{lemma}

\begin{proof}
(i) Let $\rho=\sigma\circ\iota_E\colon\OO_E\to\C$; it is injective, since $\iota_E$ and $\sigma$ are. As $E$ is a CM
field, $\rho(\bar a)=\overline{\rho(a)}$ for $a\in\OO_E$, so $|\rho(\pi_i)|^2=\rho(\pi_i\bar\pi_i)=g_i$ and
$|\rho(\pi_i)|=|\rho(\bar\pi_i)|=g_i^{1/2}$. Hence $|\sigma(\mathrm{nn}(w))|=\prod_ig_i^{(D_+(w)+D_-(w))/2}=
(g_0g_1)^{\ontothreeDF(w)/2}>0$. (ii) $Q_S,q,d\ge1$ (host conditions), so $\mathrm{pp}\ge1$ and the logarithm is
additive. The host condition $Q_S\le n^2$, where $f=2n$, gives $\log Q_S\le\log f^2$, and
$c_{1/2}(w)\le\ontothreeDF(w)+1$; hence $\ontothreeDF(w)\log Q_S+c_{1/2}(w)(\log q+\log d)\le(\ontothreeDF(w)+1)
(\log f^2+\log q+\log d)$. (iii) By (i) and (ii), $s(w)\ne0$ and $\log|\sigma(s(w))|=\frac12\ontothreeDF(w)
\log(g_0g_1)+\log\mathrm{pp}(\cdots)$ with $\log\mathrm{pp}(\cdots)\ge0$; this gives the first three lower bounds, as
$g_0,g_1\ge1$ and $\min(g_0,g_1)\le g_0g_1$. For the upper bound, $\frac12\ontothreeDF(w)\log(g_0g_1)+\log\mathrm{pp}
\le(\ontothreeDF(w)+1)\lambda^*\le\ontothreesize(w)\lambda^*$, because $\ontothreesize(w)=\ontothreeDF(w)+
f\max(q_j,1)rn_w+1\ge\ontothreeDF(w)+1$ and $\log(g_0g_1),\log(f^2qd)\ge0$.
\end{proof}

\subsubsection{Form ties and trace ties}

Fix a top structure $T$ with laws $h$, an index $j=\iota+1$ in range, a weight $\kappa$, a number field $K\supseteq
K_0$, a record object $O$ over $(S,P,c)$ with $K$ as base field, a place $v$ of $K$, an index $i$ of $S$, a
$K$-linear isomorphism $z\colon K^{n_Z(i)}\to K\otimes_{K_0}\mathcal F_\kappa$ and maps $\mathrm{act}(g)\colon
\mathcal F_\kappa\to\mathcal F_\kappa$ ($g\in P$). For a set $\Sigma\subseteq\mathcal F_\kappa$ let $\langle\Sigma
\rangle_z$ be the $\OO_K$-span of the vectors $z^{-1}(1\otimes F)$, $F\in\Sigma$.

\begin{definition}[Form tie]\label{pf:onto3:formtie}
Let $\mathrm{EX},p,w_L,w_F,T_n$ be as in \cref{pf:onto3:scal} and $X=X_{T,j}$. The \emph{form tie} at $(v,i)$ holds if
\begin{enumerate}[label=(\roman*)]
\item $M_s(i)=\langle\mathcal M_\kappa\rangle_z$, and $\rho_i(g)z^{-1}(1\otimes F)=z^{-1}(1\otimes\mathrm{act}(g)F)$
for all $g\in P$ and $F\in\mathcal F_\kappa$;
\item $\Theta(v,i)=\langle\Sigma_\Theta\rangle_z$, where $\Sigma_\Theta$ is the set of forms $F$ with $\iota_\kappa(F)\in
\Theta(S_{\mathrm{rec}},\kappa,w_L,w_F,T_n)$ and $S_{\mathrm{rec}}=\mathrm{scal}(B_T,\mathrm{EX},p,\kappa,w_L,w_F,T_n)$;
\item $\sigma(v,i)$ is the image in $K$ of $\sigma_{\mathrm{rec}}(\kappa,w_L,w_F,T_n)$;
\item $\Theta_K(v,i)=\langle\mathcal T^{\mathrm{sh}}(\kappa)\rangle_z$.
\end{enumerate}
\formallabel{GlobalTheory.FormTies.FormTieK}
\formallabel{GlobalTheory.FormTies.spanK}
\end{definition}

\begin{definition}[Eigen-traces of theta rules]\label{pf:onto3:thetaeig}
Let $U$ be a subgroup of a group $P$, $\psi\colon P\to\OO_K$ a function and $\ell(g)$ ($g\in P$) $\OO_E$-linear
automorphisms of the lattice module of $T$ at $j$. A family $(F_u)_{u\in U}$ of forms of weight $\kappa$ is an
\emph{orbit family} if there are a unimodular generic lattice $\Lambda$, trivial at the bad integer
$N_{\mathrm{bad}}=Q_Sqd\prod_pp$ (product over the level primes of the host) and at the level of the auxiliary datum
$d_\delta(j,n_\kappa)$ of $T$ and satisfying the level condition for some $N$, an admissible base tuple $x_0$, a
harmonic polynomial $P_1$ for $\Lambda$ and the weight $\kappa$, and admissible tuples $x(u)=x_0\cdot\ell(u)$, such
that $F_u$ is the short theta rule attached to $(d_\delta(j,n_\kappa),\Lambda,N,x(u),P_1)$ for every $u\in U$. The
\emph{eigen-trace lattice} $\Theta^{\mathrm{eig}}(z,U,\psi,\ell)$ is the $\OO_K$-span of the vectors
$\sum_{u\in U}\psi(u)\,z^{-1}(1\otimes F_u)$ for all orbit families $(F_u)$.
\formallabel{GlobalTheory.OnTies.thetaEig}
\formallabel{GlobalTheory.OnTies.traceSpan}
\formallabel{GlobalTheory.OnTies.IsOrbitFamily}
\end{definition}

Let now $R$ be a real datum of a point with a family $A$ of forms on a finite group $G$, a face $F$ and a tie
$\mathrm{fw}$ of the weight families $\kappa_\iota$ with $R$; we write $\mathrm{cf}_\iota\colon\mathcal F_{\kappa_\iota}
\to\mathit{LL}^{n_W(\iota)}$ for the cusp-form map of the tie (the coordinates of a cusp form in the forms of the face
family $R.\mathrm{FW}(\iota)$), and $\mathrm{tf}=(\mathrm{cf}_\iota\otimes\mathit{LL})\circ z$. For a subgroup
$U\le G$, a character $\psi\colon U\to\mathit{LL}^\times$, a phase $\mathrm{ph}\in\C$ and a type $t$, the real datum
provides: the \emph{theta base lattice} $\Theta^{\mathrm{base}}(U,\psi,\mathrm{ph},v,p)$ and the \emph{theta tuple
lattice} $\Theta^{\mathrm{tup}}(U,\psi,\mathrm{ph},t)$ in $\mathit{LL}^{n_W(\iota)}$ (the $\OO_{\mathit{LL}}$-spans of
the $\psi$-level traces, multiplied by $\mathrm{ph}$, of the theta series attached to the Weil data which are kernels
at the places over $p$ of type $\mu(v)$, resp.\ to the tuple of Weil data of type $t$); the space of source vectors
$\mathit{LL}^{d(t,U,\psi)}$ of the $(U,\psi)$-level part of the multiplicity space of $A$ at $(\iota,t)$, with its
lattice $M^{\mathrm{src}}(t,U,\psi)$, its level operator $\mathrm{lv}(t,U,\psi)$ attached to the idempotent of $A$ and its
hermitian matrix $H(t,U,\psi)$; and the notion of a \emph{class sum map} $\theta\colon\mathit{LL}^{d(t,U,\psi)}\to
\mathit{LL}^{n_W(\iota)}$ (the adelic form of $\theta(x)$ is $\mathrm{ph}$ times the theta form of the source form
of $x$). Class sums at $(t,U,\psi,\mathrm{ph})$ are \emph{rational} if a class sum map exists.

\begin{definition}[Trace tie and level compatibility]\label{pf:onto3:ontie}
In this situation let $\mathrm{inc}\colon P\to G$ be a homomorphism, $\ell(g)$ ($g\in G$) automorphisms of the lattice
module, $U\le G$, $\psi\colon U\to\mathit{LL}^\times$, $\mathrm{ph}\in\C$ and $y\colon K^{n_X(v,i)}\to
\mathit{LL}^{d(\mu(v),U,\psi)}$ an isomorphism ($K=\mathit{LL}$). The \emph{trace tie} at $(v,i)$ holds if the form tie
holds (with $j=\iota+1$, $\kappa=\kappa_\iota$) and:
\begin{enumerate}[label=(\roman*)]
\item $\mathrm{act}(g)$ is the level action of $\ell(\mathrm{inc}(g))$ on $\mathcal F_\kappa$, for $g\in P$;
\item $U(O)=\mathrm{inc}^{-1}(U)$, and $\psi_O(g)=\psi(\mathrm{inc}(g))$ for $\mathrm{inc}(g)\in U$;
\item $\mathrm{tf}\circ e_i=e^{\mathrm{FW}}\circ\mathrm{tf}$, where $e^{\mathrm{FW}}$ is the operator of the face
family attached to the coefficients of the idempotent of $A$, and $V_0(i)=e_i(K^{n_Z(i)})$;
\item $\Theta_0(v,i)=\Theta^{\mathrm{eig}}(z,U(O),\psi_O,\ell\circ\mathrm{inc})$;
\item $\Theta_J(v,i)=\mathrm{tf}^{-1}\bigl(\Theta^{\mathrm{base}}(U,\psi,\mathrm{ph},v,p)\bigr)$ and
$X_3(v,i)=e_i\,\mathrm{tf}^{-1}\bigl(\Theta^{\mathrm{tup}}(U,\psi,\mathrm{ph},\mu(v))\bigr)$;
\item $M_t(v,i)=y^{-1}\bigl(M^{\mathrm{src}}(\mu(v),U,\psi)\bigr)$ and $y\circ e^X_{v,i}=\mathrm{lv}(\mu(v),U,\psi)\circ
y$;
\item if class sums at $(\mu(v),U,\psi,\mathrm{ph})$ are rational, then $\mathrm{tf}\circ\theta'_{v,i}\circ y^{-1}$ is a
class sum map;
\item $y\circ\mathrm{pr}_{v,i}=\pi\circ y$, where $\pi$ is the $H(\mu(v),U,\psi)$-orthogonal projection onto the
$H$-orthogonal complement of $y(\ker\theta'_{v,i})$.
\end{enumerate}
The \emph{level compatibility} of $O$ with a family of top structures $(T_v)_v$ (with laws, the same core, and primes
$p(v)$, weights $w_L(v,\iota),w_F(v,\iota)$, an integer $T_n$, level isomorphisms $\ell_\iota$ and phases
$\mathrm{ph}_\iota$) holds if there are $z_\iota$, $\mathrm{act}_\iota$ and $y_{v,\iota}$ such that for every place $v$
and every $\iota$ the trace tie holds at $(v,\iota)$ for $T_v$, with $U=P_0$ and $\psi=\psi_0^{-1}$ of the face, and the
\emph{action pin} holds for $T_v$ at $\iota$: there is an action $\mathrm{act}$ of $G$ on $\mathcal F_{\kappa_\iota}$
which is the level action of $\ell_\iota$ and satisfies $\mathrm{cf}_\iota(\mathrm{act}(g)F)=\rho^{\mathrm{FW}}_\iota(g)
\,\mathrm{cf}_\iota(F)$.
\formallabel{GlobalTheory.OnTies.OnTie}
\formallabel{GlobalTheory.RealIndexBound.OnObjLevelCompatibility}
\formallabel{GlobalTheory.OnTies.ActionPin}
\end{definition}

\subsubsection{Depth data and the depth tie}

\begin{definition}[Depth data]\label{pf:onto3:depdata}
A \emph{depth page} consists of $T_n\in\N$, primes $g_1,g_2,q_1,d$, a finite set $s_m$, ranks $f,r$ and a function
$v_d$ on primes. A prime is \emph{ramified} if it lies in $s_m\cup\{q_1,d\}$ and \emph{bad} if it is ramified or equal to
$g_1$ or $g_2$. Put $n'=8fr+4$, $c_P(\mathrm{DF},p)=\lceil v_d(p)\mathrm{DF}\rceil$ for ramified $p$ and $0$ otherwise,
\[
\nu_N(\mathrm{DF},\nu,p)=\begin{cases}\mathrm{DF}/(g-1)&p=g\in\{g_1,g_2\},\\ \nu+c_P(\mathrm{DF},p)&\text{otherwise},
\end{cases}
\]
$e_\Lambda(e,\nu,p)=\max\bigl(0,\lceil e\,T_n(\nu+[p\ \text{bad}])\rceil\bigr)$, and
$Z_1(\mathrm{DF}_L,\mathrm{DF}_F)=\prod_{p}p^{T_nc_P(\mathrm{DF}_L,p)+c_P(\mathrm{DF}_F,p)}$ (product over the primes in
$s_m$ and over $q_1$ and $d$). A \emph{depth datum} consists of a depth page, weight vectors $a_L(j),a_F(j)$, integers
$n_W(j),D_+(j),D_-(j)$ and elements $\pi_1,\pi_2\in\OO_E$. For a core $T$ of $D$ with moduli spaces $\mathcal X_j$,
level $M$ and coefficient datum $c$:
\begin{enumerate}[label=(\roman*)]
\item The \emph{local depth} of a weight vector $a$ at a place $\lambda$ of $c.F$ is
\[
\nu_W(\lambda;k,a,n)=\max\Bigl(0,\ \sup_{(\mathcal O,e,y)}\ \inf_{C}\ \nu_C(y;e,k,a,n)\Bigr)\in[0,\infty],
\]
where $(\mathcal O,e,y)$ runs over complete discrete valuation rings $\mathcal O$ with an embedding of $c.F$ into their
fraction field over $\lambda$ with ramification index $e$ and algebraic integral points $y$ of $\mathcal X_j$ over
$\mathcal O$, $C$ over the integral charts at $y$, and $\nu_C(y;e,k,a,n)$ is the infimum, over harmonic test
polynomials of weight $a$ in $n$ variables and integral power series, of the normalised valuations of the associated
jet values. For a depth datum, $\nu_W(\lambda,j)=\nu_W(\lambda;M,a_L(j),n_W(j))$.
\item The \emph{normalised depth} is $\nu^*(\lambda,j)=\nu_N\bigl(\ontothreeDF(a_L(j)),\nu_W(\lambda,j)^\flat/e_\lambda,
p_\lambda\bigr)$, where $e_\lambda$ is the absolute ramification index, $p_\lambda$ the residue characteristic and
$\nu^\flat=\nu$ for $\nu<\infty$, $\infty^\flat=0$.
\item For a number field $K$ over $c.F$ and a place $v$ of $K$ above $\lambda_v$, the \emph{exponent} is
$E(v,j)=e(v|\lambda_v)\,e_\Lambda\bigl(e_{\lambda_v},\nu^*(\lambda_v,j),p_{\lambda_v}\bigr)$.
\item With $\pi_i,\bar\pi_i$ mapped to $K$: $\mathrm{nn}_K(j)=\prod_{i=1,2}\bar\pi_i^{D_+(j)}\pi_i^{D_-(j)}$ and
\[
\sigma(j)=\mathrm{nn}_K(j)\,(g_1g_2)^{frn'}\,Z_1\bigl(\ontothreeDF(a_L(j)),\ontothreeDF(a_F(j))\bigr).
\]
\end{enumerate}
The \emph{depth datum of a run} with core $T$, error datum $\mathrm{EX}$, $T_n\ge1$, a function $v_d$ and weight families
$w_L,w_F,\kappa$ has the page $(T_n,g_0,g_1,q,d,s_m,f,r,v_d)$ of the host, $a_L(j)=a(w_L(j))$, $a_F(j)=a(w_F(j))$,
$n_W(j)=n_{w_L(j)}$, $D_\pm(j)=D_\pm(\kappa(j))$, and $\pi_i$ equal to $\bar\pi_i^{\mathrm{EX}}$ if this element lies in the
half ideal of $T$ at $i$ and to $\pi^{\mathrm{EX}}_i$ otherwise. For parameters $T_p$ of size data the \emph{depth datum of
the parameters} is the depth datum of the run with $T_n=T_O$, $v_d(p)=v_p(Q_S)$ for $p\in s_m$ and $v_d(p)=\frac12$
otherwise, and the run weights $(w_L,w_F,\kappa)=(lR,fR,\kappa R)$ of \cref{sec:pf:onto3:runw}.
\formallabel{Formalism.L2.Pj.depRun}
\formallabel{GlobalTheory.RunDepth.depOf}
\formallabel{Formalism.L2.Pj.nuW}
\formallabel{Formalism.L2.Pj.nuStarAt}
\formallabel{Formalism.L2.Pj.EAt}
\formallabel{Formalism.L2.Pj.sigAt}
\end{definition}

\begin{lemma}[Crude depth bound]\label{pf:onto3:crude}
Assume that integral charts exist at every algebraic integral point of $\mathcal X_j$ and that the charts are
saturated. Let $a$ be a dominant weight vector and $n$ at least the number of variables of the harmonic test
polynomials. Then for every place $\lambda$ of $c.F$ and every $k$,
\[
\nu_W(\lambda;k,a,n)\le\max\Bigl(0,\ \frac{e_\lambda\,\ontothreeDF(a)}{p_\lambda-1}\Bigr).
\]
\formallabel{Formalism.L2.Pj.nuW_le_crude}
\end{lemma}

\begin{proof}
Since both sides are maxima with $0$, it suffices to bound $\inf_C\nu_C(y;e,k,a,n)$ for every $(\mathcal O,e,y)$ as in
\cref{pf:onto3:depdata}(i). By the first assumption there is an integral chart $C$ at $y$, and the infimum is at most
$\nu_C(y;e,k,a,n)$. By \lref{Formalism.L2.Pj.crude_inf}, which bounds the jet values of a test polynomial built
from the saturated chart directions, $\nu_C(y;e,k,a,n)\le e_\lambda\ontothreeDF(a)/(p_\lambda-1)$ whenever $a$ is dominant
and $n$ is large enough.
\end{proof}

\begin{definition}[Numerical datum and depth tie]\label{pf:onto3:runtie}
The \emph{numerical datum} $\mathrm{num}(D,L)$ of size data consists of the ranks of $D$, the scale of $L$, the host
axioms with the host primes $g_0,g_1$ and the host set $T_T$, the integers $\lfloor\Delta\rfloor$ ($\Delta$ the
constant of the spaces of $D$) and $\lfloor h_0(L)\rfloor$, and the Harish-Chandra parameters of $D$ at the real places
attached to the pairs $z$. Let $T$ be a core with error datum $\mathrm{EX}$, $T_p$ parameters of $(D,L)$, $K$ a number
field over $c.F$ and $E$ with an involution, $O$ a record object over $K$ and $\iota$ a map from the index set of $O$ to
$\N$. The \emph{depth tie} of $(O,\iota)$ holds if: $\iota(j)$ is in range for every $j$; the involution of $\OO_E$ is the
host conjugation, also after mapping to $K$; for all $v$ and $j$, $E(v,j)$, $\sigma(v,j)$ and $n_K(j)$ of $O$ are the
quantities $E(v,\iota j)$, $\sigma(\iota j)$, $\mathrm{nn}_K(\iota j)$ of the depth datum of $T_p$; $T(O)=T_O$;
$(\gamma_1,\gamma_2,q_1,d_1,s_m)$ of $O$ are $(g_0,g_1,q,d,s_m)$ of the host; $k,m,f$ of $O$ are the ranks of $D$; $x,\eps,
S_c,\eta,b_p,X_2$ of $O$ are $x$, $\eps$, $S_c$, $\eta$, $b'$ and $X''(n)$ of $L$; $v_d(O)=\log_2Q_S$; $w_t(j)=2fb$ for all
$j$; $\mathrm{th}(\ell)=\mathrm{th}_k(\ell)$ and $\mathrm{ad}(\ell)=0$ for all $\ell$; and $D_F\ge\ontothreeDF(lR(\iota j))$
for all $j$. Here $\mathrm{th}_k(\ell)=\ell^{-i_0}$, $i_0$ the least odd integer $i\ge3$ with $\ell^i>k^{3/2}$ (and
$\mathrm{th}_k(\ell)=0$ if there is none).
\formallabel{GlobalTheory.RunDepth.numOf}
\formallabel{GlobalTheory.RunDepth.RunTie}
\formallabel{Formalism.UnitaryBoundIntegral.thl}
\end{definition}

\subsubsection{Pairs of runs and record conditions}

\begin{definition}[Pairs of runs]\label{pf:onto3:tworuns}
Let $J$ be a core. A \emph{run datum} for $J$ consists of a group datum, nabla data at every index, auxiliary prime
data $d_\delta(j,n)$ for all $j,n$, a character $\xi$ of the finite ideles of $E$ and an ideal $\mathfrak h$ of $\OO_E$;
together with $J$ it forms a top structure (the \emph{full structure} of the run). A \emph{pair of runs} $r$ for $J$
and a group $G$ consists of two run data $r_1,r_2$ and, for every $\iota<m$, a map $\ell_\iota$ from $G$ to the
$\OO_E$-linear automorphisms of the lattice module of $J$ at $\iota+1$. For a place $v$ of $\mathit{LL}$ the run
\emph{used at $v$} is $r_v=r_2$ if $p_v=p_{\mathrm{de}}$ and $r_v=r_1$ otherwise, and the \emph{own prime} of $v$ is
$\mathrm{own}(v)=p'_{\mathrm{de}}$ if $p_v=p_{\mathrm{de}}$ and $\mathrm{own}(v)=p_{\mathrm{de}}$ otherwise.
\formallabel{GlobalTheory.BigOn.TwoRuns}
\formallabel{GlobalTheory.RealIndexBound.RunData}
\end{definition}

\begin{definition}[Admissible and pinned pairs]\label{pf:onto3:runsok}
Let $r$ be a pair of runs for the core $J$ of a real datum $R$ with weight families $\kappa$ and a tie $\mathrm{fw}$.
The pair is \emph{admissible} if:
\begin{enumerate}[label=(\roman*)]
\item for every $v$ the full structure of $r_v$ satisfies the \emph{full laws}: the laws of its top structure and of its
core, the short laws, canonicity of the core, $d_\delta(j,n)$ has level $1$ and ideal $\mathfrak h$ for all $j$ in
range and all $n$, uniqueness of the short generators under a change of the nabla data of the auxiliary data, and the
$\xi$-value conditions;
\item (one cover) for every $v$ and $\iota$, each $d_\delta(\iota+1,n)$ of $r_v$ is obtained from $d_\delta(\iota+1,0)$
by a change of coefficients;
\item for every $v$ and $\iota$, the character $\xi$ of $r_v$ is $1$ on the determinants of the elements of the rational
subgroup of the space $V(\iota+1)$;
\item (run pins) for every twisted place $v$ of $R$ and every $\iota$: the prime of $d_\delta(\iota+1,n_{\kappa_\iota})$ of
$r_v$ is $\mathrm{own}(v)$, the ideal $\mathfrak h$ is contained in every place of type $\mu(v)$ above $\mathrm{own}(v)$, and
there is a unitary character $\chi$ of the ideles of $E$ matching the twist character of $R$ at $p_v$ such that
$\xi(\det g)=\chi(x)$ whenever $x$ is a determinant root of $g$;
\item (level pins) for every $v$ and $\iota$, the lifts of the isomorphisms $\ell_\iota(g)$ realise the level structure of
$R$ at $\iota$;
\item (action pins) for every $v$ and $\iota$, the action pin of \cref{pf:onto3:ontie} holds for $r_v$, $\ell_\iota$ and
$\mathrm{fw}$;
\item the level ideal of $J$ is the product of the places in the set $P_g$ of $R$.
\end{enumerate}
The pair is \emph{pinned} if for every place $v$, every $j$ in range and every $n$ the prime of the auxiliary datum
$d_\delta(j,n)$ of $r_v$ is $\mathrm{own}(v)$.
\formallabel{GlobalTheory.BigOn.RunsOk}
\formallabel{GlobalTheory.BigOn.OwnPin}
\formallabel{Formalism.L2.Pj.LawsFull}
\formallabel{GlobalTheory.FormTies.OneCover}
\formallabel{GlobalTheory.OnTies.RunPin}
\formallabel{GlobalTheory.OnTies.LevelPin}
\end{definition}

\begin{definition}[Record conditions]\label{pf:onto3:isrec}
Let $T_p$ be parameters of $(D,L)$ with run weights $\kappa R,lR,fR$, $\mathrm{EX}$ an error datum, $F$ a face of a real
datum $R$ with core $J$, $\mathrm{fw}$ a tie of weight families $\kappa$, $r$ a pair of runs for $J$ and the group $G$ of
$R$, and $O$ a record object over the false side of the complex of $(R,F)$, the group
\begin{equation}\label{eq:pf:onto3:Mir}
\mathrm{Mir}(R)=\prod_{v\in P_g}\mathrm{Mir}_{r}(k_v),\qquad r=k+m-1,
\end{equation}
the product of the mirabolic subgroups of $\GL_r$ over the residue fields $k_v$ at the level places $P_g$ of $R$ (the
places above the two type primes of the host), and constants $c$. The \emph{record
conditions} for $(r,O)$ are:
\begin{enumerate}[label=(\roman*)]
\item $r$ is admissible;
\item $\kappa_\iota=\kappa R(\iota+1)$ for all $\iota<m$;
\item the core $J$ satisfies its laws, the map $\OO\to\mathit{LL}$ is injective, and $O$ satisfies the level compatibility
with the full structures of the runs $r_v$, the primes $\mathrm{own}(v)$, the weights $lR(\iota+1)$, $fR(\iota+1)$, the
integer $T_O$, the level isomorphisms $\ell_\iota$ of $r$, the phases $1$ and the inclusion $\mathrm{Mir}(R)\to G$;
\item for some involution on $\OO_E$, the depth tie of $(O,\iota\mapsto\iota+1)$ holds for $J$, $\mathrm{EX}$ and $T_p$;
\item $D_F=\max_\iota\ontothreeDF(lR(\iota+1))$, $r_r=r$, $N_c(O)=N_c(L)$, $T_2=|T_2(\mathrm{host})|$ and
$p_{\mathrm{de}}(O)=p_{\mathrm{de}}$.
\end{enumerate}
For a point $p$, a \emph{record object of $p$ at the constants $c$} is a record object over
$(S_2(p),\mathrm{Mir}(R_2(p)),c)$, and $O\in\mathcal O(p,r,c)$ means that the record conditions hold for the parameters of
the family, the data of the two-sided reduction of $p$, $r$ and $O$, and that $r$ is pinned.
\formallabel{GlobalTheory.BigOn.IsRec}
\formallabel{GlobalTheory.BigOn.RecAt}
\formallabel{GlobalTheory.BigOn.Obj}
\formallabel{GlobalTheory.OntoChain.Mir}
\end{definition}

\begin{definition}[Closed statements over points and record objects]\label{pf:onto3:forall}
For a property $Q$ of points let $\mathbf Q_{\mathrm{pt}}(Q)$ be the statement
$\ontothreeClosed\bigl(\sfname{FamHyp},(\mathcal F,C_H)\mapsto\ontothreeEv(\mathcal F,p\mapsto\ontothreeMore(p)
\Rightarrow Q(p))\bigr)$, where $\ontothreeEv(\mathcal F,\cdot)$ means: there is $x_0$ such that the property holds for
all $x>x_1(\mathcal F)$ with $x\ge x_0$, all $i$ and all points $p$ over $(x,i)$ satisfying the point conditions; and
$\ontothreeMore(p)$ is the conjunction of: every place at which the complex of $p$ has a single index is twisted, the
core ties of $p$, and the record ties of $p$. A triple $c=(C,C_{b1},C_{b2})$ is \emph{admissible} if $14\le C_{b1}$ and
$66C+4\le C_{b2}$. For a property $Q$ of points, constants and record objects let $\mathbf Q_{\mathrm{rec}}(Q)$ be the
statement
\begin{multline*}
\ontothreeClosed\Bigl(\sfname{FamHyp},\ (\mathcal F,C_H)\mapsto\exists c_0\ \forall c\ \text{admissible with
}c_0\le C:\\ \ontothreeEv\bigl(\mathcal F,\ p\mapsto\ontothreeMore(p)\Rightarrow\forall r\ \forall O\in\mathcal
O(p,r,c):\ Q(p,c,O)\bigr)\Bigr).
\end{multline*}
\formallabel{GlobalTheory.BigOn.AtPtsAlt}
\formallabel{GlobalTheory.BigOn.ForallRecordObjects}
\formallabel{Formalism.L2.Indef.Tree.OnPage.ConstBad}
\end{definition}

% ---------------------------------------------------------------------------
\subsection{The surjectivity estimate}\label{sec:pf:onto3:path}

\subsubsection{Record statements}

Let $p$ be a point of a family. Its \emph{two-sided reduction} consists of the real datum $R_2(p)$ obtained from the
real datum of $p$ by the character of its face, the face $F_2(p)$ obtained from the level-two face of $p$ by
replacing its scale by $Z_p$, and the indefinite datum $S_2(p)$, the first side of the complex of
$(R_2(p),F_2(p))$. Its numbers are $N_p(C,C_H)$ again.

\begin{lemma}\label{pf:onto3:split}
For every point $p$, $D(S_2(p))=D(S_{\mathrm f})$, and $S_2(p)$ and $S_{\mathrm f}$ have the same local indices
\eqref{eq:pf:onto3:Iw}, twisted places and types. Consequently, for all $C,C_H$ and every function $\ontothreeerr$,
\begin{equation}\label{eq:pf:onto3:split}
\ontothreeONw(S_2(p),N_p(C,C_H),\ontothreeerr)\ \Longrightarrow\ \ontothreeONw(S_{\mathrm f},N_p(C,C_H),
\ontothreeerr)\quad\text{and}\quad D(S_{\mathrm f})=D(S_2(p)).
\end{equation}
\formallabel{GlobalTheory.BigOn.d_split_eq}
\formallabel{GlobalTheory.BigOn.iw_split_eq}
\formallabel{GlobalTheory.BigOn.onWith_of_split}
\formallabel{GlobalTheory.T.onClauses_with_error_of_splitPart}
\end{lemma}

\begin{proof}
The phases of the real datum of $p$ are all equal to $1$ (a field of the point). For such real data the complex
attached to the datum and the face coincides with the complex attached to the character-free realisation
(\lref{GlobalTheory.CharRealise.toCx3_eq_one}, which states this equality of complexes), and the first side of
the latter is $S_2(p)$ by definition; this gives the equalities. Since the numbers are those of $Z_p$ in both cases,
the three conditions (ON1)--(ON3) for $S_2(p)$ and for $S_{\mathrm f}$ are the same inequalities.
\end{proof}

Pairs of runs, admissible and pinned pairs, record objects of points, the record conditions and the
quantifiers $\mathbf Q_{\mathrm{pt}}$, $\mathbf Q_{\mathrm{rec}}$ are those of
\cref{pf:onto3:tworuns,pf:onto3:runsok,pf:onto3:isrec,pf:onto3:forall}; for a property $Q$ of points with pairs of runs
we put $\mathbf Q_{\mathrm{run}}(Q)=\mathbf Q_{\mathrm{pt}}\bigl(p\mapsto\forall r\text{ admissible and pinned}:\
Q(p,r)\bigr)$.

\begin{definition}[The record statements]\label{pf:onto3:recstmts}
The following statements are used (we write $N_p$ for $N_p(C,C_H)$).
\begin{itemize}
\item The \emph{runs statement} $\ontothreeSt{runs}=\mathbf Q_{\mathrm{pt}}(\exists r\text{ admissible and pinned})$.
\item The \emph{rationality statement} $\ontothreeSt{rat}=\mathbf Q_{\mathrm{pt}}(\mathrm{Rat})$, where
$\mathrm{Rat}(p)$ asserts that at every twisted place $w$ of $R_2(p)$ and every $j$ the class sums of $R_2(p)$ at the
type $\mu(w)$ (for the subgroup $P_0$, the character $\psi_0^{-1}$ and the phase $1$) are rational, i.e.\ a class
sum exists in the sense of the development.
\item The \emph{exported statement} $\ontothreeSt{ex}$ and the \emph{maximality statement} $\ontothreeSt{max}$,
of the form $\mathbf Q_{\mathrm{pt}}(\cdot)$, assert the export ties, resp.\ the class-span property, at all twisted
places and indices (\cref{pf:onto3:export,pf:onto3:runs}).
\item The \emph{geometric statement} $\ontothreeSt{geo}=\mathbf Q_{\mathrm{rec}}\bigl($\ref{pf:onto3:L:sup},
\ref{pf:onto3:L:span}, \ref{pf:onto3:L:lemS}, \ref{pf:onto3:L:data}, \ref{pf:onto3:L:idx} in existential form,
and for all $C,C_H$ the depth leaves \ref{pf:onto3:L:good}--\ref{pf:onto3:L:level} for $N_p$$\bigr)$.
\item The \emph{rank}, \emph{twist}, \emph{share} and \emph{bad-place statements}
$\ontothreeSt{rank}=\mathbf Q_{\mathrm{rec}}($\ref{pf:onto3:L:rank}$)$,
$\ontothreeSt{tw}=\mathbf Q_{\mathrm{rec}}($\ref{pf:onto3:L:tw1} in existential form$)$,
$\ontothreeSt{sh}=\mathbf Q_{\mathrm{rec}}($\ref{pf:onto3:L:share}$)$ and
$\ontothreeSt{bad}=\mathbf Q_{\mathrm{rec}}($\ref{pf:onto3:L:bad}$)$.
\item The \emph{exponent statement} $\ontothreeSt{exp}=\mathbf Q_{\mathrm{rec}}\bigl(\forall C,C_H$: the exponent
bound for $N_p$ and $O\bigr)$.
\item The \emph{page statement} $\ontothreeSt{page}=\mathbf Q_{\mathrm{rec}}\bigl($\ref{pf:onto3:L:page},
\ref{pf:onto3:L:sizes}, \ref{pf:onto3:L:x}, \ref{pf:onto3:L:badc}, (LS), (LY), $b(O)=b(L_{x,i})$,
$E_{1u}(Z_p)\ge0$, and for all $C,C_H$ the leaf \ref{pf:onto3:L:E1} for $N_p$ with $C_E=C\bigr)$.
\end{itemize}
Here (LS), the \emph{level-square condition}, asserts that $\gamma_1,\gamma_2$ are distinct primes, that every
residue field $k_v$ in \eqref{eq:pf:onto3:Mir} has $\gamma_1$ or $\gamma_2$ elements, and that each of the two
values occurs for exactly $f$ places, $f$ being the rank $f$ of the datum $D_{x,i}$; and (LY) asserts
$2f(O)/(\min(\gamma_1,\gamma_2)-1)\le y$.
\formallabel{GlobalTheory.BigOn.GeometricStatement'}
\formallabel{GlobalTheory.BigOn.ExponentStatement}
\formallabel{GlobalTheory.BigOn.PageStatement}
\end{definition}

\subsubsection{Record objects of points}

\begin{lemma}[Mirabolic orders]\label{pf:onto3:mirorders}
Let $G=\prod_{v\in\Upsilon}\mathrm{Mir}_{n+1}(k_v)$ and $U=\prod_vN_{n+1}(k_v)$ for finite fields $k_v$ with
$q_v$ elements. Then $|U|=\prod_vq_v^{n(n+1)/2}$ and $[G:U]=\prod_v\prod_{i=1}^n(q_v^i-1)$. If
$\sum_v1/(q_v-1)\le y$, then $(1-y)|U|\le[G:U]$. If every $q_v$ lies in $\{\gamma_1,\gamma_2\}$ for distinct
primes $\gamma_1,\gamma_2$, each value being taken exactly $f_N$ times, then for every finite place $w$ of $K$,
$v_w(|U|)=e_w f_Nn(n+1)/2$ if $p_w\in\{\gamma_1,\gamma_2\}$ and $v_w(|U|)=0$ otherwise.
\formallabel{Formalism.L2.Indef.Tree.leafIndexU_of}
\formallabel{Formalism.L2.Indef.Tree.leafUnip_of}
\end{lemma}

\begin{proof}
$\mathrm{Mir}_{n+1}(k)\cong\GL_n(k)\ltimes k^n$ has order $q^n\cdot q^{n(n-1)/2}\prod_{i=1}^n(q^i-1)$, and
$|N_{n+1}(k)|=q^{n(n+1)/2}$. Hence $[G:U]/|U|=\prod_v\prod_{i=1}^n(1-q_v^{-i})\ge1-\sum_v\sum_{i\ge1}q_v^{-i}=
1-\sum_v1/(q_v-1)\ge1-y$, by the inequality $\prod_k(1-a_k)\ge1-\sum_ka_k$ for $a_k\in[0,1]$ (induction on the
number of factors). In the last case $|U|=(\gamma_1\gamma_2)^{f_Nn(n+1)/2}$, and $v_w(\gamma_i)=e_w$ if
$p_w=\gamma_i$ and $0$ otherwise.
\end{proof}

\begin{lemma}[Link conditions of the two-sided reduction]\label{pf:onto3:links2}
Let $p$ be a point satisfying the point conditions such that every place at which the complex of $p$ has a single index
is twisted. Then the \emph{link conditions} of the two-sided reduction hold:
\begin{enumerate}[label=(\roman*)]
\item the character condition of $(R_2(p),F_2(p))$;
\item the family of forms of $p$ satisfies its first admissibility condition;
\item $R_2(p)$ is uniform, the base links of $(R_2(p),F_2(p))$ hold, and $S_2(p)$ has rank rows ($D(S_2(p))>0$ and
$d_W(j)=D(S_2(p))$ for all $j$);
\item the bar tie of the tie of $p$ and, at every index, the theta row pins and the phase tie with all phases $1$.
\end{enumerate}
\formallabel{GlobalTheory.T.links2_of_single_index_twisted}
\formallabel{GlobalTheory.BigOn.links2_of_rows_alt}
\formallabel{GlobalTheory.BigOn.Links2}
\formallabel{GlobalTheory.T.l_unif_base_links_rank_of_pt_split}
\formallabel{GlobalTheory.BigOn.SplitRows.uniform_and_base_links_and_rank_of_split_data}
\formallabel{GlobalTheory.BigOn.SplitRows.lUnif_of_rows_alt}
\formallabel{GlobalTheory.BigOn.SplitRows.base_links_of_split_data}
\formallabel{GlobalTheory.BigOn.SplitRows.rank_conditions_of_split_data}
\formallabel{GlobalTheory.BigOn.SplitRows.split_lifts_bijective_of_split_data}
\formallabel{GlobalTheory.BigOn.SplitRows.rat_theta_of_split_data}
\formallabel{GlobalTheory.BigOn.SplitRows.level_proviso_of_split_data}
\formallabel{GlobalTheory.T.bar_tie_om_pins_phase_tie_of_pt_split}
\formallabel{GlobalTheory.BigOn.ompins_onto_alt}
\formallabel{GlobalTheory.BigOn.conjugateReadingsVanishing_splitPart}
\formallabel{GlobalTheory.BigOn.thetaRowData_pins_splitPart}
\formallabel{GlobalTheory.BigOn.phaseCompatible_splitPart_of_liftUniformisation}
\end{lemma}

\begin{proof}
(i) The real datum of $p$ has the existence property at $p$ (a point condition), and its weak characterisation is part of
the point. By \lref{GlobalTheory.BigOn.splitPartRealisationWithScalars_char_of_weak}, the characterisation of weak real
data (\lref{Formalism.IndefR.Links.RealDat3.characterisation_of_weak}), together with the fact that the places with
non-zero index of the split part are twisted (\lref{GlobalTheory.T.splitPart_nonzeroIndexPlaces_twisted}, which uses
the hypothesis on single-index places), gives the character condition through
\lref{GlobalTheory.BigOn.char_charOf}. (ii) is one of the point conditions.
(iii) Uniformity is the primitive uniformisation of the point conditions. By
\lref{Formalism.IndefR.Links.complex_data_conditions_base_level_of_primitive_links}, the admissibility of the family, the
separation condition, the weak characterisation and the primitive links of the point imply the conditions of the complex
data at the base level; since all phases of the real datum of $p$ are $1$, the complex of $(R_2(p),F_2(p))$ is that of the
character-free realisation (\lref{GlobalTheory.CharRealise.toCx3_eq_one}). Among these conditions are the three base
links at the false side (bijectivity of the split lifts, rationality of the theta data, and the level proviso) and the
equalities $D>0$, $d_W(j)=D$ for the indefinite datum at the base level, which is $S_2(p)$.
(iv) The bar tie and the theta row pins are fields of the point. The phase tie at $i$ is a field of the point for the
effective phases of the real datum, under uniformity of lifts (which holds by (iii)); the effective phases are all equal to
$1$, another field of the point.
\end{proof}

\begin{proposition}[Twisted places bound for record objects of points]\label{pf:onto3:chain}
Let $p$ be a point, $r$ a pair of runs and $O$ a record object of $p$ at constants $c$ such that $(p,r,O)$ satisfies
the record conditions. Assume the link conditions of the real datum $R_2(p)$ and the face $F_2(p)$ (the condition on
the family of forms, the level-group, $\chi_W$, model, conjugation and theta-tie conditions, the uniformity
condition, the base links, and $\mathrm{Rat}(p)$), that $S_2(p)$ has rank rows, and that $O$ satisfies (LS), (LY),
\ref{pf:onto3:L:sup}, \ref{pf:onto3:L:span}, \ref{pf:onto3:L:lemS}, \ref{pf:onto3:L:idx} (existential),
\ref{pf:onto3:L:data}, \ref{pf:onto3:L:rank}, \ref{pf:onto3:L:ident}, \ref{pf:onto3:L:tw2}, \ref{pf:onto3:L:tw1}
(existential), \ref{pf:onto3:L:share}, \ref{pf:onto3:L:bad}, \ref{pf:onto3:L:good}--\ref{pf:onto3:L:badc}, the
exponent bound and \ref{pf:onto3:L:E1}, for numbers $N$ and a constant $C_E\ge C_E(c)$. Then
$\ontothreeONw(S_2(p),N,\ontothreeerr)$ holds for some function $\ontothreeerr$.
\formallabel{GlobalTheory.T.twistedPlacesBoundChain_twoRunChain}
\formallabel{GlobalTheory.OnAtPlacesBase.on_base_atPlaces}
\end{proposition}

\begin{proof}
The level compatibility contained in the record conditions identifies $U$ with $\prod_{v\in P_g}N_r(k_v)$ and
$\psi$ on $U$ with $u\mapsto\prod_v\chi_v(\sum_iu_{i,i+1})$ for the level character of $R_2(p)$, whose components
$\chi_v$ are non-trivial, and it identifies $e_j$ with an idempotent with image $V_0(j)$
(\lref{GlobalTheory.RealIndexBound.level_of_tied}, which states these identifications together with the
consequences \ref{pf:onto3:L:mackey} and \ref{pf:onto3:L:proj}). By \cref{pf:onto3:mirabolic} this gives
\ref{pf:onto3:L:mackey}, and \ref{pf:onto3:L:proj} holds. The depth tie gives $f(O)=f$, the rank $f$ of the datum,
and $r_r=r=n+1$ with $n=k+m-2$. By (LS), $\sum_v1/(q_v-1)\le2f/(\min(\gamma_1,\gamma_2)-1)$, which is at most $y$ by
(LY); so \cref{pf:onto3:mirorders} (with $f_N=f$) gives \ref{pf:onto3:L:indexU} and \ref{pf:onto3:L:unip}. By
\lref{GlobalTheory.BigOnBase.exists_tieCoord_base}, from the level compatibility, the link conditions,
$\mathrm{Rat}(p)$, \ref{pf:onto3:L:ident} and \ref{pf:onto3:L:tw2} one obtains new coordinate maps
$\varphi_j$ (isomorphisms onto spaces of forms composed with the projection) such that the record object with
these $\varphi_j$ satisfies \ref{pf:onto3:L:coord}; the other leaves do not involve $\varphi_j$. Now
\cref{pf:onto3:oncore} applies.
\end{proof}

\begin{proposition}[The onto path]\label{pf:onto3:path}
The runs, rationality, geometric, rank, twist, share, bad-place, exponent and page statements
$\ontothreeSt{runs}$, $\ontothreeSt{rat}$, $\ontothreeSt{geo}$, $\ontothreeSt{rank}$, $\ontothreeSt{tw}$,
$\ontothreeSt{sh}$, $\ontothreeSt{bad}$, $\ontothreeSt{exp}$, $\ontothreeSt{page}$ imply $\sfname{SurjectiveOntoX}$.
\formallabel{GlobalTheory.BigOn.OntoPath.onto_path_alt}
\formallabel{GlobalTheory.T.onto_x_of_all_estimates}
\end{proposition}

\begin{proof}
By the closure properties of $\ontothreeClosed$ the conjunction of the nine statements is closed, and it suffices
to show: if $\eps>0$, $\mathcal F\in\sfname{FamP}(\eps)$, $C_H\in\R$ and the nine statements hold for
$(\mathcal F,C_H)$ (in the form inside $\ontothreeClosed$), then the property $Q_{\mathrm{ON}}(\mathcal F,C_H)$ of
\cref{pf:onto3:stmts} holds. (Its hypothesis $\sfname{DIVinN}(\mathcal F)$ is not needed.)

Let $x_R,x_Q$ be the thresholds of the two statements of the form $\mathbf Q_{\mathrm{pt}}$, and let $c_G$, $c_K$,
$c_T$, $c_S$, $c_B$, $c_E$, $c_P$ be the constants $c_0$ of the seven statements of the form $\mathbf Q_{\mathrm{rec}}$.
Let $c_0$ be the maximum of these seven constants and $1$, and $c=(c_0,14,66c_0+4)$; then $c$ is admissible and all
seven statements apply to $c$, with thresholds $x_1,\dots,x_7$. Put $C_0=2\max(C_E(c),0)\ge0$ and let $x_0$ be the
maximum of $x_R,x_Q,x_1,\dots,x_7$. Let $x>x_1$ with $x\ge x_0$, $i\in I_x$, $p\in\mathrm{Pt}(\mathcal F,x,i)$ with
$\ontothreePC(p)$ and $\ontothreeMore(p)$, and $C\ge C_0$.
\begin{enumerate}[label=(\arabic*)]
\item The runs statement gives an admissible pinned pair of runs $r$, and the rationality statement gives
$\mathrm{Rat}(p)$.
\item By \lref{GlobalTheory.BigOn.exists_isRec}, which constructs a record object satisfying the record
conditions from an admissible pair of runs and the canonical-core and half-error conditions contained in
$\ontothreeMore(p)$, there is a record object $O$ of $p$ at $c$ such that $(p,r,O)$ satisfies the record
conditions; hence $O\in\mathcal O(p,r,c)$.
\item The seven record statements give: \ref{pf:onto3:L:sup}, \ref{pf:onto3:L:span}, \ref{pf:onto3:L:lemS},
\ref{pf:onto3:L:data}, \ref{pf:onto3:L:idx} (existential) and the depth leaves for $N_p(C_E(c),C_H)$;
\ref{pf:onto3:L:rank}; \ref{pf:onto3:L:tw1} (existential); \ref{pf:onto3:L:share}; \ref{pf:onto3:L:bad}; the
exponent bound for $N_p(C',C_H')$ for all $C',C_H'$; and \ref{pf:onto3:L:page}--\ref{pf:onto3:L:badc}, (LS), (LY),
$E_{1u}(Z_p)\ge0$ and \ref{pf:onto3:L:E1} for $N_p(C_E(c),C_H)$ with $C_E=C_E(c)$.
\item From $\ontothreePC(p)$ and the first auxiliary condition (single-index places are twisted),
\cref{pf:onto3:links2} gives the link conditions of $(R_2(p),F_2(p))$
(characterisation, the condition on the family of forms, uniformity, base links) and the rank rows of $S_2(p)$.
With $\mathrm{Rat}(p)$ these give \ref{pf:onto3:L:ident} and \ref{pf:onto3:L:tw2}
(\lref{GlobalTheory.ClassSumIdentity.leafIdentity_of_isRec} and
\lref{GlobalTheory.BaseTwistIncl.leafTW2_of_isRec}, which deduce these two leaves from the record conditions,
the link conditions and $\mathrm{Rat}(p)$).
\item \Cref{pf:onto3:chain}, applied with $N=N_p(C_E(c),C_H)$ and $C_E=C_E(c)$, gives a function
$\ontothreeerr$ with $\ontothreeONw(S_2(p),N_p(C_E(c),C_H),\ontothreeerr)$. By \eqref{eq:pf:onto3:split},
$\ontothreeONw(S_{\mathrm f},N_p(C_E(c),C_H),\ontothreeerr)$ holds and $D(S_{\mathrm f})=D(S_2(p))>0$.
\item The error function of $S_{\mathrm f}$ is the canonical error \eqref{eq:pf:onto3:erreq}, $E_{1u}(Z_p)\ge0$,
and $J_0\ge0$ at the twisted places by the exponent bound for $N_p(C,C_H)$. Since $C\ge C_0\ge2C_E(c)$,
\cref{pf:onto3:onwith} gives that $S_{\mathrm f}$ satisfies $\ontothreeON(N_p(C,C_H))$.\qedhere
\end{enumerate}
\end{proof}

\subsubsection{The nine statements}

We now establish the nine statements, except the geometric and the exponent statements:
\lref{GlobalTheory.BigOn.geoSt_of_pageSt_alt} states that the page statement implies the geometric statement,
and \lref{GlobalTheory.T.expSt_of_base_change_and_crystalline_lattices} states that the arithmetic inputs listed
in \cref{pf:onto3:ontoX} imply the exponent statement. Several of the proofs below are assemblies of results whose
proofs are long; we state precisely what is used.

\begin{proposition}[Runs, rationality, maximality]\label{pf:onto3:runs}
The runs statement $\ontothreeSt{runs}$, the rationality statement $\ontothreeSt{rat}$ and the maximality
statement $\ontothreeSt{max}$ hold. Here $\ontothreeSt{max}$ asserts, at every point $p$ with $\ontothreePC(p)$ and
$\ontothreeMore(p)$, every twisted place $v$ and every $j$, the \emph{class-span property}: the base theta lattice
of $R_2(p)$ at $(v,j)$ (for the own prime of $v$) is the $\OO_{\mathit{LL}}$-span of the theta series attached to
any system of representatives of the classes, for every base kernel.
\formallabel{GlobalTheory.T.runs_estimate_statement'}
\formallabel{GlobalTheory.BigOn.bigOn_runsSt}
\formallabel{GlobalTheory.T.rational_estimate_statement'}
\formallabel{GlobalTheory.BigOn.bigOn_ratSt}
\formallabel{GlobalTheory.BigOn.maxSt_alt}
\formallabel{GlobalTheory.T.maximum_estimate_statement'}
\end{proposition}

\begin{proof}
All three hold for every family, with $\eps_1=1$ and trivial thresholds, and for all $x$:
\begin{itemize}
\item At every point $p$ with $\ontothreePC(p)$ and $\ontothreeMore(p)$ there is an admissible pinned pair of runs
(\lref{GlobalTheory.BigOn.RunsExist.exists_runs_at_alt}). The two runs are produced, one for each of the two
primes $\ell\in\{p_{\mathrm{de}},p'_{\mathrm{de}}\}$, by
\lref{GlobalTheory.T.exists_parameterRun_with_laws_and_pins_at_special_prime}, which builds from the core of
$p$ a run with all laws whose auxiliary prime data are located at $\ell$ and verifies the pins; the level
isomorphisms come from \lref{GlobalTheory.T.finiteGroup_levelStructure_action_fixedRuns}, an action of the
finite group on the level structures fixing the runs.
\item At every such point $\mathrm{Rat}(p)$ holds (\lref{GlobalTheory.BigOn.twistedRational_R2_F2}, which
states exactly this).
\item The class-span property follows from the link conditions of the two-sided reduction
(\cref{pf:onto3:links2}), the uniqueness of the base theta kernel series
(\lref{GlobalTheory.T.baseThetaKernel_series_unique}) and the span of the kernel theta series over class
representatives (\lref{GlobalTheory.T.kernel_thetaSeries_span_over_class_representatives}), whose last
hypothesis (the values of the characters $\delta$ at the place lie in the image of $\mathit{LL}$) is part of the
core ties in $\ontothreeMore(p)$; \lref{GlobalTheory.TraceClassRank.spanOverClasses_of} combines the two.\qedhere
\end{itemize}
\end{proof}

\begin{proposition}[Rank at places]\label{pf:onto3:rank}
Let the \emph{comparison statement} $\ontothreeSt{cmp}$ be the statement of the form $\mathbf Q_{\mathrm{run}}$
asserting: for every twisted place $v$ and every $j$, and every finite place $\lambda$ of $\mathit{LL}$ whose residue
characteristic is not the own prime of $v$, the localisation at $\lambda$ of the trace lattice of the run used at
$v$ is contained in the localisation at $\lambda$ of the base theta lattice. Then $\ontothreeSt{cmp}$ holds, and
$\ontothreeSt{max}$, $\ontothreeSt{geo}$ and $\ontothreeSt{cmp}$ imply $\ontothreeSt{rank}$.
\formallabel{GlobalTheory.BigOn.compareOffOwn_atRuns_alt}
\formallabel{GlobalTheory.T.bigOn_comparison_at_twisted_places}
\formallabel{GlobalTheory.BigOn.rankLocSt_of_max_geo_runs}
\formallabel{GlobalTheory.T.rank_at_twisted_places_of_max_geo}
\end{proposition}

\begin{proof}
The comparison holds closedly with the threshold of the numerical conditions along families
(\cref{pf:onto3:pagenum}). At a point with $\ontothreePC$ and $\ontothreeMore$ and runs as above, the link
conditions of the two-sided reduction, the span of the class sums
(\lref{GlobalTheory.ClassSumSpanCore.classSumsSpan_core}, which states that the class sums span the relevant
core lattice) and the pins of the runs are the hypotheses of
\lref{GlobalTheory.T.traceLattice_le_span_thetaSeries_awayFromOwnPrime}, which gives the inclusion of the trace
lattice in the span of the theta series away from the own prime. For the second assertion, combine the three
closed statements, take the maximum of the thresholds, and apply at each $O\in\mathcal O(p,r,c)$ the lemma
\lref{GlobalTheory.BigOn.rank_at_twisted_places_of_rec_span}, which deduces \ref{pf:onto3:L:rank} from the
class-span property at the twisted places, the leaves \ref{pf:onto3:L:span} and \ref{pf:onto3:L:data} of the
geometric statement, and the comparison.
\end{proof}

\begin{proposition}[Twist statement]\label{pf:onto3:twist}
The exported, geometric and rank statements imply the twist statement $\ontothreeSt{tw}$.
\formallabel{GlobalTheory.BigOn.TwistBoundAtPlace.twistSt_of_alt}
\formallabel{GlobalTheory.BigOn.TwistBoundAtPlace.hasTwistBound_at}
\end{proposition}

\begin{proof}
Combine the three closed statements and take the maximum of the thresholds, and of the threshold of the twist
goodness $\sfname{TwGood}$, which is a condition of $\sfname{FamHyp}$: eventually, every member $D_{x,i}$ is
twist-good at each of its twisted places for its norm choice and its twisting datum. Let $O\in\mathcal O(p,r,c)$.
Every twisted place $v$ of $R_2(p)$ is twisted for the member, and the conjugate of a twisted place is twisted.
The level compatibility of the record conditions and the export tie at $(v,j)$ (\cref{pf:onto3:export}) form a tie
between $O$ and the twisting datum at $v$. Then \lref{GlobalTheory.TwistIndex.exists_leafTW1_of_twGoodAt}
gives numbers $\eps_0(v,j)\ge0$ with $\sum_j\eps_0(v,j)\le p_v^{-2}$ and the index inequality of \ref{pf:onto3:L:tw1}
at $v$: the twist goodness provides such numbers and inequalities for the lattices of the twisting datum, and the
tie, together with \ref{pf:onto3:L:sup} and the lattice part of \ref{pf:onto3:L:rank}, transports them to $O$.
Since $p_{\bar w}=p_w$, summing over $v\in\{w,\bar w\}$ gives \ref{pf:onto3:L:tw1}.
\end{proof}

Let $p$ be a point and $i$ an index; let $n_W(i)$ be the dimension of the $i$-th space of $p$ and $\mathrm{sp}_i$ its
subspace on which $P_0$ acts by the character $\mu=\chi_W|_{P_0}\cdot\psi_0^{-1}$. For $y,\delta\ge0$ consider the
\emph{share inequality} $\mathrm{Sh}(p,i,y)$: $(1-y)\,n_W(i)\le|P_0|\cdot\dim\mathrm{sp}_i$, and the \emph{bad-share
inequality} $\mathrm{Bd}(p,i,\delta)$: the intersection of $\mathrm{sp}_i$ with the kernel of the good projection
$e_{\mathrm{Good}}$ has dimension at most $\delta\dim\mathrm{sp}_i$. The \emph{model condition}
$\mathrm{Mod}(p,i,T)$ asserts the existence of: a finite $G_p$-set, a two-level datum $M$ over the host with $f$
places $\Psi$ and a weight $\kappa$, a representation $\rho$ of $G_p$ on the space $M.\mathrm{cusp}(\kappa)$ of cusp
forms (a pull-back representation), an embedding $\sigma\colon\mathit{LL}\to\C$ and a $\sigma$-semilinear,
$G_p$-equivariant map from the $i$-th space onto $M.\mathrm{cusp}(\kappa)$ carrying a basis to a basis, level
hypotheses for the subgroups attached to half-unipotent subgroups, and either $p_0(M)\in\{0,r\}$ (the definite
case) or $0<p_0(M)<r$ together with an archimedean condition $\mathrm{Arch}(M,\kappa,T)$.

\begin{lemma}[Shares from models]\label{pf:onto3:sharemodel}
Let $T\in\N$ and $y\ge0$. Assume $\mathrm{Mod}(p,i,T)$, that $P_0$ is half-unipotent for $p$, that
$2r+50r^2fp_0(r-p_0)+2\le T$ for all $0\le p_0\le r$, and that $d_Y(d_Y+6)/T\le y$ with $d_Y=f\,r(r-1)/2$. Then
$\mathrm{Sh}(p,i,y)$ holds.
\formallabel{GlobalTheory.ShareOn.shareW_of_modelAt}
\end{lemma}

\begin{proof}
Through the semilinear map of $\mathrm{Mod}(p,i,T)$, $\dim_\C$ of the $\sigma\circ\mu$-eigenspace of $P_0$ in
$M.\mathrm{cusp}(\kappa)$ equals $\dim\mathrm{sp}_i$, and $\dim_\C M.\mathrm{cusp}(\kappa)=n_W(i)$. In the definite
case, \lref{GlobalTheory.T.definite_cuspForms_finite_regularRepresentation} shows that $M.\mathrm{cusp}(\kappa)$
is, as a representation of every subgroup satisfying the level hypothesis (in particular of $P_0$), a multiple of
the regular representation, so that $|P_0|$ times the dimension of every eigenspace is $n_W(i)$; the claim follows
from $y\ge0$. In the indefinite case the hypothesis on $T$ bounds the size parameter of $M$ by $T$;
\lref{GlobalTheory.T.twoLevelFlagDomain_cuspForm_ranks_shifted_weights} (ranks of cusp forms on two-level flag
domains with shifted weights, under the archimedean condition) and
\lref{GlobalTheory.T.twoLevelFlagDomain_eigenspace_share} (the resulting lower bound for eigenspace dimensions
of subgroups satisfying the level hypothesis) give $(1-d_Y(d_Y+6)/T)\,n_W(i)\le|P_0|\dim\mathrm{sp}_i$, and
$d_Y(d_Y+6)/T\le y$.
\end{proof}

\begin{proposition}[Model, share and bad-place statements]\label{pf:onto3:share}
Let the \emph{model statement} $\ontothreeSt{mod}$ be the statement of the form $\mathbf Q_{\mathrm{run}}$ asserting
$\mathrm{Mod}(p,i,T_O)$ for all $i$, where $T_O$ is the parameter of the parameter choice of the family. Then
$\ontothreeSt{mod}$ holds; it implies the share statement $\ontothreeSt{sh}$; and together with Green's
parametrisation \ref{pf:onto3:P:green} it implies the bad-place statement $\ontothreeSt{bad}$.
\formallabel{GlobalTheory.ShareOn.modelSt_alt}
\formallabel{GlobalTheory.ShareOn.shareWSt_of_modelSt_alt}
\formallabel{GlobalTheory.ShareOn.shareSt_of_shareWSt_alt}
\formallabel{GlobalTheory.BigOn.badSt_of_badW_alt}
\formallabel{GlobalTheory.BadWOn.badWSt_of_greenCuspidal_of_modelSt}
\end{proposition}

\begin{proof}
\emph{Model statement.} At a point with $\ontothreePC$, $\ontothreeMore$ and runs as above,
\lref{GlobalTheory.ShareOn.exists_runModel} produces a two-level model $M$ with an equivariant identification
of the $i$-th space with a space of cusp forms (through a level-group transport and uniformisation);
\lref{GlobalTheory.ShareOn.Reading.reading_of_runModel} gives the representation, the embedding and the
semilinear map; \lref{GlobalTheory.ShareOn.levelCondition_restrict_of_halfUnipotent} gives the level
hypotheses for half-unipotent subgroups, using the closed statement
\lref{GlobalTheory.BigOn.typePrimes_gt_three_r_alt} (the type primes exceed $3r$); and the closed statement
\lref{GlobalTheory.ShareOn.p0_extreme_or_archOn_of_runModel} gives the dichotomy $p_0\in\{0,r\}$ or the
archimedean condition for every such model.

\emph{Shares.} Let $\ontothreeSt{shw}$ be the statement of the form $\mathbf Q_{\mathrm{run}}$ asserting
$\mathrm{Sh}(p,i,1/(4(fk)^3))$ for all $i$. It follows from the model statement and \cref{pf:onto3:sharemodel}: the
subgroup $P_0$ is half-unipotent (\lref{GlobalTheory.T.exists_whittaker_lam_psi_half_unip}, which states this
for the points with $\ontothreePC$), and the two conditions on $T=T_O$ hold eventually along every family
(\lref{GlobalTheory.T.shareSize_bounds_along_family}). Finally
\lref{GlobalTheory.ShareOn.Comparison₂.leafShareI_of_shareW} deduces \ref{pf:onto3:L:share} for every
$O\in\mathcal O(p,r,c)$ from $\mathrm{Sh}(p,j,y)$, $y=1/(4(fk)^3)$, for all $j$.

\emph{Bad places.} Let $\ontothreeSt{bw}$ be the statement of the form $\mathbf Q_{\mathrm{run}}$ asserting
$\mathrm{Bd}(p,i,\delta_x)$ for all $i$, with $\delta_x=\exp(-x^{\eps_z/7})$. It follows from Green's
parametrisation, the model statement, the bounds \lref{GlobalTheory.T.badShareSize_bounds_along_family} on
bad-share sizes along every family and the bound \lref{GlobalTheory.T.hodgeLineMultiple_le_familyThreshold}
on multiples of the Hodge line, by \lref{GlobalTheory.BadWOn.badW_of_del_bound_of_T_bound}, which deduces the
bad-share inequality from these inputs at a point. Since eventually $\delta_x\le\frac14$
(\lref{GlobalTheory.T.badShareSize_bounds_along_family}), \lref{GlobalTheory.BadOnRec.bad_rec_alt}, which
deduces $(1-\delta_x)d_0(j)\le D(S_2(p))$ for $O\in\mathcal O(p,r,c)$ from $\mathrm{Bd}(p,j,\delta_x)$, gives
\ref{pf:onto3:L:bad} with $\delta=\delta_x$.
\end{proof}

\begin{proposition}[Page statement]\label{pf:onto3:page}
The page statement $\ontothreeSt{page}$ holds.
\formallabel{GlobalTheory.BigOn.pageSt_alt}
\formallabel{GlobalTheory.T.page_estimate_statement'}
\end{proposition}

\begin{proof}
Combine the closed statements \lref{GlobalTheory.T.indexBoundParameterRequirements_eventually_family}, which
states size conditions on the parameters eventually along every family, and
\lref{GlobalTheory.T.leafLevelSq_and_leafYBound}, which states (LS) and (LY) for record objects of points; take
$c_0=\max(c_1,1)$ and the threshold $\max(x_s,x_2,3C)$. For $O\in\mathcal O(p,r,c)$,
\lref{GlobalTheory.BigOn.Check.twoRunDatum_boundUnit_at_indefinitePoint} gives \ref{pf:onto3:L:E1} for
$N_p(C',C_H')$ with $C_E=C'$, for all $C',C_H'$, and $b(O)=b(L_{x,i})$; and
\lref{GlobalTheory.IndexAssembly.numericLeaves} gives \ref{pf:onto3:L:page}--\ref{pf:onto3:L:badc} from the
depth tie, the numerical identities of the record conditions, the size conditions, $1\le C$, the admissibility of
$c$ and $3C\le x$. Finally $E_{1u}(Z_p)=x^{\eps}\bigl(fmb\,x^{\eps}+N_c\bigr)\ge0$ with $b=e^{x^\eps}>0$ and
$N_c\ge0$.
\end{proof}

Let (RS) denote the rationality statement for theta series deduced from global Artin reciprocity in
\lref{GlobalTheory.seriesRatS_and_serInt_of_globalArtin}: for every real datum satisfying the standard laws
and ties and every lift index with standard frames there is a complex number of absolute value $1$ (a phase) such
that the theta series attached to every good cover are rational up to this phase.

\begin{proposition}[Exported statement]\label{pf:onto3:export}
Assume rapid decay of cusp forms \ref{pf:onto3:P:decay}, global Artin reciprocity, the existence of ray class
fields, the classification of the discrete series \ref{pf:onto3:P:ds}, Casselman--Wallach globalisation
\ref{pf:onto3:P:cw} and (RS). Then the exported statement $\ontothreeSt{ex}$ holds: eventually, at every point $p$
with $\ontothreePC(p)$ and $\ontothreeMore(p)$, every twisted place $v$ and every $j$, there is a semilinear reading
map from the $j$-th space of $R_2(p)$ to the forms of the twisting datum at $v$ which reads the base theta lattice
at $v$ (for the own prime of $v$) and the theta tuple at $\mu(v)$ into the corresponding objects of the twisting
datum (the \emph{export tie} at $(v,j)$).
\formallabel{GlobalTheory.T.exportSt_of_reciprocity_and_casselman_wallach}
\formallabel{GlobalTheory.BigOn.exportSt_of_reading_alt}
\formallabel{GlobalTheory.T.exportSt_of_discrete_series_cover}
\formallabel{GlobalTheory.BigOn.Exp.exportSt_of_hypotheses_alt}
\formallabel{GlobalTheory.BigOn.Exp.InductionConditions.exportSt}
\end{proposition}

\begin{proof}
The statement follows from thirteen closed conditions on twisted places: the canonical existence of reading maps
and the canonical projection; twist admissibility; a condition at the place; rationality of the compatibility; a
recursion on the side; existence of adapted data; conjugation statements for the base and for the twists; and upper
and lower bounds for the base and for the twists. Indeed, combine them (taking the maximum of the thresholds); at a
twisted place the canonical existence statement provides the reading map, and each of the two required reading
identities is obtained as two opposite inequalities (the base upper and lower bounds, resp.\ the twist upper and
lower bounds), combined with the conjugation statements. The thirteen conditions are supplied as follows:
\begin{itemize}
\item twist admissibility, the canonical existence and projection and the existence of adapted data hold
unconditionally (\lref{GlobalTheory.T.twist_admissible'}, \lref{GlobalTheory.T.bigOn_canEx},
\lref{GlobalTheory.T.bigOn_canProj}, \lref{GlobalTheory.T.adaptEx'});
\item the twist characters come from Artin reciprocity (\lref{GlobalTheory.T.twist_character_of_artin'});
\item the condition at the place follows from rapid decay, Artin reciprocity, ray class fields, (RS) and a
low-weight lower bound at the place (\lref{GlobalTheory.T.at_statement_of_low_weight_twAdm},
\lref{GlobalTheory.T.lowAt}); the rational compatibility and the recursive side follow from it and twist
admissibility (\lref{GlobalTheory.T.tieRat_of_twAdm_and_at}, \lref{GlobalTheory.T.reciprocity_side_of_twAdm});
\item the conjugation of the base follows from rapid decay, reciprocity, ray class fields and (RS)
(\lref{GlobalTheory.T.conjBase_of_ray_class_field_and_seriesRat}), and the conjugation of the twists in
addition from the discrete series and Casselman--Wallach
(\lref{GlobalTheory.T.conjTw_of_discrete_series_and_casselman_wallach});
\item the four bounds follow from rapid decay, reciprocity, ray class fields and, for the lower bounds, (RS), and
for the twist bounds in addition from twist admissibility and the twist characters
(\lref{GlobalTheory.T.baseLe_of_cuspRapidDecay_artin_rayClassField},
\lref{GlobalTheory.T.baseGe_of_cuspRapidDecay_artin_rayClassField_seriesRat},
\lref{GlobalTheory.T.twist_upper_bound_of_prints}, \lref{GlobalTheory.T.twist_lower_bound_of_prints}).\qedhere
\end{itemize}
\end{proof}

\subsubsection{Assembly}

\begin{theorem}\label{pf:onto3:fivest}
The exported, twist, share, bad-place and exponent statements imply $\sfname{SurjectiveOntoX}$.
\formallabel{GlobalTheory.BigOn.ontoX_of_twist_shares_exp_alt}
\formallabel{GlobalTheory.T.geoSt_of_pageSt'}
\end{theorem}

\begin{proof}
By \cref{pf:onto3:page} the page statement holds, hence the geometric statement
(\lref{GlobalTheory.BigOn.geoSt_of_pageSt_alt}); by \cref{pf:onto3:runs,pf:onto3:rank} the runs, rationality,
maximality and comparison statements hold, hence the rank statement. Now \cref{pf:onto3:path} applies. The exported
statement is not used here (it enters through the twist statement).
\end{proof}

\begin{theorem}[Surjectivity estimate]\label{pf:onto3:ontoX}
Assume \ref{pf:onto3:P:green}--\ref{pf:onto3:P:cw}, that is: Green's parametrisation of cuspidal characters,
Clozel's theorem, the $\ell$-adic characters of algebraic Hecke characters, Chebotarev's theorem with primes of
degree one, the Galois representations of polarised cuspidal representations, base change for definite unitary
groups, the realisation of crystalline lattices by prismatic $F$-crystals, the local properties above $\ell$ of the
$\ell$-adic characters, rapid decay of cusp forms, $\sfname{RecTExists}$, the classification of the discrete series
and Casselman--Wallach globalisation, each in the formulation of the development. Then $\sfname{SurjectiveOntoX}$
holds.
\formallabel{GlobalTheory.BigOn.ontoX_of_hypotheses_alt_alt}
\end{theorem}

\begin{proof}
The $\ell$-adic character statement \ref{pf:onto3:P:hecke} contains the existence of the global Artin maps, hence
global Artin reciprocity (\lref{GlobalTheory.Published.HeckeLAdicCharacter.globalArtin}, which extracts it);
hence the existence of ray class fields (\lref{GlobalTheory.Published.GlobalArtinReciprocity.rayClassFieldSt},
which deduces it from Artin reciprocity) and (RS). \Cref{pf:onto3:export} gives the exported statement. The page
statement holds (\cref{pf:onto3:page}) and implies the geometric statement
(\lref{GlobalTheory.BigOn.geoSt_of_pageSt_alt}); with \cref{pf:onto3:runs,pf:onto3:rank} this gives the rank
statement, and \cref{pf:onto3:twist} the twist statement. \Cref{pf:onto3:share} gives the share and bad-place
statements (the latter using \ref{pf:onto3:P:green}), and
\lref{GlobalTheory.T.expSt_of_base_change_and_crystalline_lattices}, which states that
\ref{pf:onto3:P:clozel}--\ref{pf:onto3:P:above} and \ref{pf:onto3:P:llc} imply the exponent statement, gives the
latter. \Cref{pf:onto3:fivest} concludes.
\end{proof}

% ---------------------------------------------------------------------------
\subsection{The low-weight bound}\label{sec:pf:onto3:low}

\paragraph{Sizes and the lower bound.}
Let $D$ be a datum (\cref{sec:families:theta}) with ranks $f,r,n_V,k,m,\dots$, set of embedding types
$\mathcal T$, host with type primes $g_0,g_1$ and root discriminant $r_E$, and let $i<m$ and $b\in\R$. The
\emph{size condition} $\mathrm{Sz}(D,i,b)$ means: for every $\tau\in\mathcal T$ the rational numbers $e_i(\tau)$,
$v_l(\tau)$, $b_l(\tau)$ attached to $D$ (the infinity types of the lifts) satisfy $0<e_i(\tau)\le2b$,
$0\le v_l(\tau)\le b$ and $|b_l(\tau)-e_i(\tau)|\le b$ for all $l$, and $(kb)^8\le g_l$ for $l=0,1$. Let $h(\cdot,\cdot)\ge1$
denote the hook product of the development (the product of the hook lengths of the Young diagram attached to a
highest weight; its positivity is \lref{GlobalTheory.HookDiagram.hookOf_pos_log}), and put
\begin{gather*}
L_p(D,b)=\log\bigl(b\,r_E\,g_0g_1\,n_V\bigr),\qquad
H^+(D,i)=\sum_{\rho\in\mathcal T}\log h\bigl(p_W(D,i,\rho),a^+(D,i,\rho)\bigr),\\
B_{\mathrm{low}}(D,i,b)=H^+(D,i)-8fn_V^2L_p(D,b),\\
H_\sigma(D,i)=\sum_{\rho\in\mathcal T}\log\Bigl(h\bigl(p_W(D,i,\rho),a^+_\sigma(D,i,\rho)\bigr)\,
h\bigl(q_W(D,i,\rho),a^-_\sigma(D,i,\rho)\bigr)\Bigr)
\end{gather*}
for embeddings $\sigma\colon\mathit{LL}\to\C$, where $p_W,q_W$ are the signature data of the spaces $W_c(i)$ at
$\rho$ and $a^\pm,a^\pm_\sigma$ the corresponding weights. We call $H_\sigma(D,i)$ the \emph{main term}.

\paragraph{The scalars of the second side.}
Let $p$ be a point. Its real datum $R$ carries, for every lift index $j$, type $t$, index $p'\in\mathrm{Ix}$ and
infinite place $v$, the numbers $\lambda_{j,t,p',v}$ of the second side $S_{\mathrm t}$. The \emph{weak
characterisation} of $R$ (a field of the point) contains the following conventions, in which $\theta_{G,j,t}$ is the
theta lift of index $j$ and type $t$, $W_{p'}(j)$ the type-$p'$ part of its target at the base point, and $l_{6m}(j)$
that of the scale of $p$:
\begin{enumerate}[label=(\alph*)]
\item if $W_{p'}(j)\ne0$ and $\theta_{G,j,t}\theta_{G,j,t}^*$ acts on $W_{p'}(j)$ by a scalar $l$, then
$\lambda_{j,t,p',v_0}=l$; if $W_{p'}(j)=0$, then $\lambda_{j,t,p',v_0}=\exp(l_{6m}(j))$;
\item for $v\ne v_0$: call the type-$p'$ part \emph{non-degenerate at $v$} if some vector $x$ in it (lifted to $v$)
has $\mathrm{maj}_v(x,x)\ne0$; if it is non-degenerate at $v$ and the theta images satisfy
$\langle\theta x,\theta y\rangle_v=l\,\mathrm{maj}_v(x,y)$ for $x,y$ in it, then $\lambda_{j,\mu_0,p',v}=l$; if it is
degenerate at $v$, then $\lambda_{j,\mu_0,p',v}=\exp(l_{6m}(j))$;
\end{enumerate}
here $\mathrm{maj}_v$ is the hermitian majorant at $v$ of the side. Moreover the primitive rows of $p$ contain the
positivity $\lambda>0$ and the eigenvalue property at the base place ($\theta_G\theta_G^*$ acts on each $W_{p'}(j)$
by $\lambda_{j,t,p',v_0}$). By definition of the scale of a point, $l_{6u}=0$ and
$l_{6m}(j)=\Sigma_h(j)-C_{\mathrm{sz}}fn_V^2L_p(D,b)$, where $b=b(L_{x,i})$, $C_{\mathrm{sz}}$ is the size constant
of the family and $\Sigma_h(j)$ is a sum over $\rho\in\mathcal T$ of logarithms of hook products attached to the
record of $p$; hence $l_6(j)=l_{6m}(j)$ for all $C$.

\begin{definition}\label{pf:onto3:adelcdef}
For a real datum $R$ and a face $F$, the \emph{adelic constant bound} $\ontothreeAdel(R,F)$ asserts: under the core
hypotheses on $(R,F)$, for every $i$ and $b$ with $\mathrm{Sz}(D,i,b)$, every type $t$ prime to the host, every
complex embedding $\sigma$ over the reflex field, every local level $P$, all characters $\chi_\mu,\chi_\eta$ and
local characters conjugate in the prescribed way, every Weil datum $\omega$ for the pair $(V_c(i),W_c(i))$ with
these splitting characters, all orthonormal bases at the real places, level vectors $f_0,f$, coefficient data,
Haar measures with fundamental sets, class representatives and hermitian module structures on the adelic forms:
there is a theta lift $\mathrm{th}$, and for every theta lift $\mathrm{th}$ and every non-zero type part $X$ of its
target, $\mathrm{th}\circ\mathrm{th}^*$ (adjoint for the level) acts on $X$ by a scalar $l>0$ with
\[
H_\sigma(D,i)-8fn_V^2L_p(D,b)\le\log l .
\]
\end{definition}

\begin{proposition}[Inner product bound]\label{pf:onto3:adelc}
Assume \ref{pf:onto3:P:decay} and \ref{pf:onto3:P:ws}--\ref{pf:onto3:P:fd}: rapid decay of cusp forms, the
Weil--Siegel formula, the second Siegel--Weil statement, the local zeta integrals of the doubling method and their
unramified computation, Bernstein--Zelevinsky finiteness and Harish-Chandra's formal degrees. Let $D$ be a datum
admitting a parameter choice and with \emph{standard frames} (the frames of the hermitian spaces $V_c(i)$ are
orthonormal and span their standard lattices). Then $\ontothreeAdel(R,F)$ holds for every real datum $R$ and face
$F$ over $D$.
\formallabel{GlobalTheory.LowRowAdel.adelC_of_siegelWeil_of_doubling_alt}
\formallabel{GlobalTheory.T.adelC_of_parameter_choice_stdV}
\end{proposition}

\begin{proof}
By \lref{GlobalTheory.T.adelicLift_scalar_of_innerProductFormula}, $\ontothreeAdel(R,F)$ follows from nine
component statements listed below. Given the data of \cref{pf:onto3:adelcdef}, a theta lift exists by the first
component; for a theta lift and a non-zero type part, the inner product formula expresses
$\langle\mathrm{th}^*U_1,\mathrm{th}^*U_2\rangle$ as a constant times a sum of Petersson products of lifts of $U_1$
and $U_2$, which shows that $\mathrm{th}\,\mathrm{th}^*$ acts on the type part by the explicit scalar
\[
l=\Bigl(\prod_\rho\dim K_\rho\Bigr)\cdot d\cdot\vol(\mathcal F)\cdot\prod_\rho h_\rho^+h_\rho^-\cdot\|f\|^2\cdot
d_{\mathrm{off}}^{-1}\cdot\prod_\rho\int|\langle v_\rho,\pi_\rho(g)v_\rho\rangle|^2\,dg\cdot|E_{\mathrm{off}}(x)|,
\]
a product of the dimensions of the lowest $K$-types, the dimension $d\ge1$ of the coefficient space, the volume of
the fundamental set, the hook products (whose logarithm is the main term), the norm of the level Schwartz function,
the inverse of a local factor at the finitely many bad places, the integrals of the squared matrix coefficients of
the discrete series at the real places, and a partial Euler product evaluated at the Satake parameters $x$ of the
type part (\lref{GlobalTheory.LowRow.adelC_of_innerProductFormula}, which carries out this deduction). The bound
$\log l\ge H_\sigma(D,i)-8fn_V^2L_p$ then follows from the identity volume $=$ level index $\times$ mass, the lower
bound for the level index times the Schwartz norm, the lower bound for the mass of the unitary group from
orthonormal lattice frames, the upper bound for formal degrees (equivalently, lower bounds for the integrals of
squared matrix coefficients), the bound $|x_{w,a}|\le\Nm w^{(r-1)/2-(1+t)/2}$ on the Satake parameters, which bounds
the Euler product from below, and the size conditions, by an explicit numerical inequality. The nine components are:
\begin{enumerate}[label=(\arabic*)]
\item the lift lands in the cuspidal multiplicity space
(\lref{GlobalTheory.T.adelicLift_mem_cuspForm_multiplicitySpace}, from rapid decay);
\item a law for the source-level character (\lref{GlobalTheory.T.finiteDatum_sourceLevelCharacter_law}, from a
single-block local integral, \lref{GlobalTheory.SingleBlock.localIntegralSingleBlock}, and restricted-tensor place
models, \lref{GlobalTheory.PlaceTensor.placeModels});
\item the inner product formula for the twisted theta lift
(\lref{GlobalTheory.T.twistedThetaLift_innerProductFormula_of_hypotheses}, from rapid decay, the Weil--Siegel
formula, the parameter choice, and the continuation of Siegel Eisenstein series, which
\lref{ResidueKilled.sec_of_dm} obtains from the doubling method);
\item a formula for matrix coefficients of Fock vectors
(\lref{GlobalTheory.T.fockVector_matrixCoefficient_formula}, from the harmonic module statement
\lref{GlobalTheory.HarmModule.harmonicModule_of}, which follows from the decay of Fock-space functions,
\lref{GlobalTheory.FockDecay.harmonicDecay}, and an integral of powers of $\cosh$,
\lref{GlobalTheory.CartanInt.coshIntegral});
\item the $\Delta$-norm identities (\lref{GlobalTheory.LowRowAdel.deltaNorm_all}: products of minors of Fock
forms are products of hook products);
\item an Euler product for the finite zeta integral
(\lref{GlobalTheory.T.finiteZetaIntegral_eulerProduct_of_partialL}), from its identification with a partial
$L$-value (\lref{GlobalTheory.T.finiteZetaIntegral_thetaLift_eq_partialLValue}, from rapid decay,
Bernstein--Zelevinsky finiteness and the unramified doubling integral);
\item a formula for the product of the theta level index and the Schwartz norm
(\lref{GlobalTheory.T.thetaLevelIndex_mul_schwartzNorm_formula});
\item a lower bound for the mass of the unitary group from orthonormal lattice frames
(\lref{GlobalTheory.T.unitaryMassLowerBound_of_orthonormalLatticeFrames}, using the standard frames);
\item an upper bound for formal degrees at the real places
(\lref{GlobalTheory.T.formalDegreeRealPlaceUpperBound_of_harishChandra}, from Harish-Chandra's formula).
\end{enumerate}
The second Siegel--Weil statement \ref{pf:onto3:P:sw} is not used.
\end{proof}

\begin{lemma}\label{pf:onto3:lowplaces}
Let $p$ be a point with real datum $R$ (second side $S_{\mathrm t}$), $b\in\R$, and assume
$l_{6m}(j)\le B_{\mathrm{low}}(D,j,b)$ for all $j$.
\begin{enumerate}[label=(\roman*)]
\item (Base place.) Assume the \emph{base-place bound} $\mathrm{Low}_0(R,b)$: for every $j$, every non-zero $x$ in
the target of $\theta_{G,j,\mu_0}$ at the base point and every $l>0$ with $\theta_{G,j,\mu_0}\theta_{G,j,\mu_0}^*x=lx$,
one has $B_{\mathrm{low}}(D,j,b)\le\log l$. Then $l_{6m}(j)\le\log\lambda_{j,\mu_0,p',v_0}$ for all $j,p'$.
\item (Other places.) Assume the \emph{off-base bound} $\mathrm{Low}_\infty(R,b)$: for every $j$, every infinite
place $v\ne v_0$ and every vector $x$ of the base type, $e^{B_{\mathrm{low}}(D,j,b)}\Re\,\mathrm{maj}_v(x,x)\le
\|\theta x\|_v^2$; and assume that the theta images satisfy $\langle\theta x,\theta y\rangle_v=
\lambda_{j,\mu_0,p',v}\mathrm{maj}_v(x,y)$ on the type-$p'$ part. Then $l_{6m}(j)\le\log\lambda_{j,\mu_0,p',v}$ for
all $j,p'$ and $v\ne v_0$.
\end{enumerate}
\formallabel{GlobalTheory.T.log_liftEigenvalue_lowerBound_basePlace}
\formallabel{GlobalTheory.LowRow.l6m_le_log_lam3_base_place}
\formallabel{GlobalTheory.T.log_liftEigenvalue_lowerBound_otherPlaces}
\formallabel{GlobalTheory.LowOff.l6m_le_log_lam3_other_places}
\formallabel{GlobalTheory.LowOff.le_log_of_inner_eq}
\end{lemma}

\begin{proof}
(i) If $W_{p'}(j)=0$, then $\lambda_{j,\mu_0,p',v_0}=\exp(l_{6m}(j))$ by convention (a), and the claim is an
equality. Otherwise take $x\ne0$ in $W_{p'}(j)$; by the eigenvalue property, $\theta_G\theta_G^*x=\lambda x$ with
$\lambda=\lambda_{j,\mu_0,p',v_0}>0$, and $\mathrm{Low}_0(R,b)$ gives $l_{6m}(j)\le B_{\mathrm{low}}(D,j,b)\le
\log\lambda$.
(ii) If the type part is degenerate at $v$, $\lambda=\exp(l_{6m}(j))$ by convention (b). Otherwise choose $x$ with
$\mu:=\mathrm{maj}_v(x,x)\ne0$ and put $u=\theta x$, $\lambda=\lambda_{j,\mu_0,p',v}>0$. Then
$\|u\|^2=\langle u,u\rangle=\lambda\mu$, so $\mu$ is real, and positive because $\lambda>0$ and $\mu\ne0$; and
$e^{B_{\mathrm{low}}}\mu\le\|u\|^2=\lambda\mu$ gives $B_{\mathrm{low}}\le\log\lambda$, whence the claim.
\end{proof}

\begin{proposition}[Lift eigenvalue lower bound]\label{pf:onto3:lamweak}
Let $p$ be a point satisfying $\ontothreePC(p)$, $b\in\R$ with $\mathrm{Sz}(D,j,b)$ for all $j$, and assume
$\ontothreeAdel(R,F)$, that the base type $\mu_0$ is prime to the host, $l_{6m}(j)\le B_{\mathrm{low}}(D,j,b)$ and
$l_6(j)\le l_{6m}(j)$ for all $j$. Then $S_{\mathrm t}$ satisfies $\ontothreeLsix(N_p(C,C_H))$.
\formallabel{GlobalTheory.T.liftEigenvalueLowerBound_of_weakCharacterisation}
\formallabel{GlobalTheory.LowRow.lamLow_of_weak}
\end{proposition}

\begin{proof}
From the primitive rows, the condition on the family of forms, the separation of types, the weak characterisation,
the existence condition at the infinite places and phases equal to $1$ one obtains the core hypotheses, the family
condition and the base links of the level-two face (\lref{GlobalTheory.T.splitPart_baseLinks_of_primitiveRows},
which states this deduction). \Lref{GlobalTheory.T.conjugatesOfRationalThetaLift_from_characterisation} gives
the Galois conjugates of the rational theta lift at every index, from the core, the standard character, the ties
$\mathrm{fw}$ and $\mathrm{bar}$, the extension condition, the uniformity, the pins and the phase ties.
\Lref{GlobalTheory.T.liftAdjoint_eigenvalue_lower_bound_base_embedding} deduces $\mathrm{Low}_0(R,b)$ from
$\ontothreeAdel$ at the base embedding and the sizes, and
\lref{GlobalTheory.T.rationalLift_lowerBound_off_basePlace} deduces $\mathrm{Low}_\infty(R,b)$ from
$\ontothreeAdel$, the Galois conjugates and the existence of extremal vectors (provided by $\mathrm{fw}$); the
pairing identity of \cref{pf:onto3:lowplaces}(ii) is part of the primitive rows. By
\cref{pf:onto3:lowplaces}, $l_6(j)\le l_{6m}(j)\le\log\lambda_{j,\mu_0,p',v}$ for all $j,p',v$, which is
$\ontothreeLsix(N_p(C,C_H))$.
\end{proof}

\begin{lemma}\label{pf:onto3:hooksum}
If $C_{\mathrm{sz}}\ge8$ and $\mathrm{Sz}(D,j,b)$, then
$\Sigma_h(j)-C_{\mathrm{sz}}fn_V^2L_p(D,b)\le B_{\mathrm{low}}(D,j,b)$.
\formallabel{GlobalTheory.T.hookLogSum_sub_sizeTerm_le_lowerBound}
\formallabel{GlobalTheory.LowMainTerm.main_le_lowB}
\end{lemma}

\begin{proof}
By \lref{GlobalTheory.LowMainTerm.sum_logHook_eq}, which identifies the two sums of logarithms of hook products,
$\Sigma_h(j)=H^+(D,j)$; under $\mathrm{Sz}(D,j,b)$ one has $L_p(D,b)\ge0$
(\lref{GlobalTheory.AdelicScalar.SizeBounds.sizes}, which collects elementary consequences of the size
condition). Hence $C_{\mathrm{sz}}fn_V^2L_p\ge8fn_V^2L_p$.
\end{proof}

\begin{theorem}[Low-weight bound]\label{pf:onto3:lamlow}
Under the hypotheses of \cref{pf:onto3:adelc}, $\sfname{LambdaLowerBound}_2$ holds.
\formallabel{GlobalTheory.BigTop.lamLow_of_hypotheses_alt_alt}
\formallabel{GlobalTheory.T.finiteness_inputs_real_data_three_alt}
\formallabel{GlobalTheory.LowRow.lamLow_closing_alt}
\end{theorem}

\begin{proof}
The property $C_{\mathrm{sz}}(\mathcal F)\ge8$ is closed on every class: take $\eps_1=1$ and thresholds with
$K_{\mathrm{mid}}\equiv8$; since $C_{\mathrm{sz}}$ is one of the middle constants, every $K$-large family has
$C_{\mathrm{sz}}\ge8$. By the closure
properties it suffices to show $Q_{\mathrm{L6}}(\mathcal F,C_H)$ for every family with $C_{\mathrm{sz}}\ge8$. By
\lref{GlobalTheory.T.thetaLift_weightSizes_eventually_along_family} there is $x_0$ such that
$\mathrm{Sz}(D_{x,i},j,b(L_{x,i}))$ holds for all $x\ge x_0$, $i\in I_x$ and $j$, for every family. Let $p$ be a point
over such $(x,i)$ with $\ontothreePC(p)$, and $C\in\R$. The real datum of $p$ satisfies the adelic constant bound by
\cref{pf:onto3:adelc}, since $D_{x,i}$ has the parameter choice of the family and standard frames (a field of the
point). The core hypotheses give that the bad places of the base tuple are prime to the host
(\lref{GlobalTheory.T.badPlacesPrimeToHost_baseTuple}, which states this). The two numerical hypotheses of
\cref{pf:onto3:lamweak} hold: $l_{6m}(j)\le B_{\mathrm{low}}(D,j,b)$ by \cref{pf:onto3:hooksum} and the formula for
$l_{6m}$, and $l_6(j)=l_{6m}(j)$ since $l_{6u}=0$. \Cref{pf:onto3:lamweak} gives $\ontothreeLsix(N_p(C,C_H))$ for
$S_{\mathrm t}$.
\end{proof}

% ---------------------------------------------------------------------------
\subsection{Run weights, page conditions and generating families}\label{sec:pf:onto3:runw}

This subsection proves the results on runs, pages and records which are used in the proofs of the index bound
along runs (\lref{GlobalTheory.FinalJoin.runHolds_of_parOK}) and of the geometric and exponent statements
(\lref{GlobalTheory.BigOn.geoSt_of_pageSt_alt},
\lref{GlobalTheory.T.expSt_of_base_change_and_crystalline_lattices}). We keep the notation of
\cref{sec:pf:onto3:ties}; in addition $\mathrm{rd}_E$ is the root discriminant parameter of the host, $\delta_2$ the
thick density of $L$, $p_j=k+j-1$ and $q_j=m-j$ the signature at $j$ (so $p_j+q_j=r$), $n_V=r+t$, $f=2n$, and
\begin{gather*}
z_B=fr\mathrm L,\qquad \ontothreePP(C,A)=C\,z_B^{A},\qquad \ontothreeRR(C)=\mathrm{rd}_E^{\,C},\\
q_\wedge=\min(g_0,g_1),\qquad c_0=2\cdot5^4r,\qquad \lambda^*=\log(g_0g_1)+\log(f^2qd).
\end{gather*}
For a sequence $a$ indexed by $\{0,\dots,p-1\}$ the \emph{gaps} are $\mathrm{gap}_i(a)=a(i)-a(i+1)$ for $i<p-1$ and
$\mathrm{gap}_{p-1}(a)=a(p-1)$; a non-increasing sequence of natural numbers is the sequence of partial sums
$a(i)=\sum_{i'\ge i}\mathrm{gap}_{i'}(a)$ of its gaps.

\subsubsection{The run weights}

Let $T_p$ be parameters of a parameter choice for $(D,L)$; they provide integers $T_O$, $n_P=8fr+4$ and $n_W$ (with
$n_W$ divisible by $8$), $D_p=f\binom r2$, and the relations $r+t=n_P+T_On_W$ and $40(D_p+2)\le T_O$. Let $\kappa_{\mathrm
def}$ be the weight definition of $D$, with big and small entries $\kappa^{\mathrm{big}}_j(\tau,i)\ge0$ and
$\kappa^{\mathrm{sm}}_j(\tau,i)\le0$. For $j$ in range put $T'=T_O+\delta_j+2$.

\begin{definition}[Run weights]\label{pf:onto3:runweights}
For $j$ in range the \emph{run weights} are the weights $\kappa R(j)$, $lR(j)$, $fR(j)$ with
\begin{enumerate}[label=(\roman*)]
\item $\kappa R(j)$: $a^+_\tau(i)=\kappa^{\mathrm{big}}_j(\tau,i)$, $a^-_\tau(c)=-\kappa^{\mathrm{sm}}_j(\tau,q_j-1-c)$,
$n=n_V$;
\item $lR(j)$: $n=n_W$, and the gaps of $a^\pm_\tau$ are $\max\bigl(0,\lfloor(\mathrm{gap}_i(\kappa R(j)^\pm_\tau)+2)/T'
\rfloor\bigr)$ for non-last $i$, $\max\bigl(0,\lfloor(\kappa R(j)^+_\tau(p_j-1)-(\delta_j+2)n_W-2r)/T'\rfloor\bigr)$
for the last index of $a^+$ and $\max(0,\lfloor\kappa R(j)^-_\tau(q_j-1)/T'\rfloor)$ for the last index of $a^-$;
\item $fR(j)$: $n=n_P$, and the gaps of $a^\pm_\tau$ are $\max\bigl(0,\mathrm{gap}_i(\kappa R(j)^\pm_\tau)-T_O\,\mathrm{gap}_i
(lR(j)^\pm_\tau)\bigr)$.
\end{enumerate}
Outside the range, $\kappa R(j)$ is the scalar weight with $n=n_V$.
\formallabel{Formalism.MainNotions.DIV.Cons.fR}
\formallabel{Formalism.MainNotions.DIV.Cons.fWt}
\end{definition}

\begin{lemma}[Weight relations]\label{pf:onto3:runsum}
For $j$ in range:
\begin{enumerate}[label=(\roman*)]
\item $n(\kappa R(j))=n_V$ (for every $j$);
\item $fR(j)^\pm_\tau(i)=\kappa R(j)^\pm_\tau(i)-T_O\,lR(j)^\pm_\tau(i)$ for all $\tau,i$, and $n(fR(j))=n_P$; in particular
$T_O\,lR(j)^\pm_\tau(i)\le\kappa R(j)^\pm_\tau(i)$ and $T_O\ontothreeDF(lR(j))\le\ontothreeDF(\kappa R(j))$;
\item $n(\kappa R(j))=n(fR(j))+T_O\,n(lR(j))$, $\delta_j\le T_O$ and $r\le n(fR(j))$.
\end{enumerate}
\formallabel{Formalism.MainNotions.DIV.Cons.kaWt_n}
\formallabel{GlobalTheory.RunDepth.run_weights_sum}
\formallabel{GlobalTheory.RunDepth.mul_rows_lR_le_kaR}
\formallabel{GlobalTheory.RunWeightTest.isF'_fWt}
\formallabel{GlobalTheory.RunDepth.mul_lWt_aP_le}
\formallabel{GlobalTheory.RunDepth.mul_lWt_aM_le}
\formallabel{GlobalTheory.RunDepth.mul_DF_lWt_le}
\formallabel{GlobalTheory.DepthBoundTest.dimY_eq_flag}
\end{lemma}

\begin{proof}
(i) Both cases of the definition have $n=n_V$. (ii) Write $\kappa=\kappa R(j)$, $l=lR(j)$, $g=\mathrm{gap}_i(\kappa^\pm_\tau)
\ge0$ and $g'=\mathrm{gap}_i(l^\pm_\tau)$. We claim $T_Og'\le g$. For a non-last index, if $g+2<T'$ then $g'=0$;
otherwise $g'\ge1$ and $g'T'\le g+2$, so $T_Og'\le g+2-g'(\delta_j+2)\le g$. For the last index of $a^+$, $g'T'\le
g-(\delta_j+2)n_W-2r\le g$ if $g'>0$, and similarly for $a^-$. Hence the gaps of $fR(j)$ are exactly $g-T_Og'$, and
summing gaps from $i$ to the end gives $fR(j)^\pm_\tau(i)=\kappa^\pm_\tau(i)-T_Ol^\pm_\tau(i)$; summing over $\tau,i$
gives the inequality for $\ontothreeDF$. (iii) $n(fR)=n_P$, $n(lR)=n_W$ and $n_V=r+t=n_P+T_On_W$. Since $p_j+q_j=k+m-1
=r$, $\delta_j=f\bigl(p_jq_j+\binom{p_j}2+\binom{q_j}2\bigr)=f\binom r2=D_p\le T_O$. Finally $f\ge1$, so $r\le8fr+4=n_P$.
\end{proof}

\begin{lemma}[Sizes of the cloud]\label{pf:onto3:cloud}
Let $w_j(\tau)$ denote the weights of the datum at $j$ and the type $\tau$, assume $3\le x$ and $kb'\le b$, and let a
parameter choice for $(D,L)$ be given. Then $S_c+\eta\le2b$, and for $j$ in range: $0\le w_j(\tau)\le2b$ for all
$\tau$; $w_j(\tau)\le2b'$ for the types $\tau$ attached to the first real places $(0,t)$; and the weights of the
numerical datum $\mathrm{num}(D,L)$ at $(j,z)$ are $w_j$ at the type attached to $z$.
\formallabel{GlobalTheory.RunDepth.sizes_le_two_mul}
\formallabel{GlobalTheory.T.cloudSizes_between_zero_and_twice_b}
\end{lemma}

\begin{proof}
Recall $b'=xb/X'$ with $X'=\exp(x^{\eps/3})$, $w''(n)=b'/(xn)$, and that $k\ge5$, $n\ge1$. From $kb'\le b$ we get
$5b'\le b$, and $w''(n)\le b'/3$. The parameter choice gives $t\le\eta$, $t\eta\le w''(n)/x$ (hence, as $t\ge1$ by the
size data, $1\le\eta\le b'/9$), $kS_c\le w''(n)$ (hence $S_c\le b'/3$), and $b_T(t)\le b'+b'/4$ for every $t$, the last
from $b_T(t)\le b'+N^*x^{\eps}\Lambda_Gb'/k$ together with $\rho\le x^{-2}\le1$ for the quantity
$\rho=4N^*x^{\eps}\Lambda_G/k+C_\rho(\log x)\Lambda_G/k$ of the parameter choice, whose second term is non-negative
($\Lambda_G>0$ because $0<4S_c\le\Lambda_Gb'/k$), so that $N^*x^\eps\Lambda_G/k\le\frac14$. Therefore $S_c+\eta\le b'/3+b'/9\le2b$. The size data put $w_j(\tau)$ in windows: $b\le w_j(\tau)\le b+\eta+1$ at
the thick places, $b_T(t)+b'/x\le w_j(\tau)\le b_T(t)+b'/x+\eta$ at the top place attached to $t$, and $w''(n)\le
w_j(\tau)\le w''(n)+\eta$ at the thin places. With $b_T(t)\ge b'>0$, $\eta\le b'/9$ and $b'/x\le b'/3$, all three
windows lie in $[0,2b]$, and the last two in $[0,2b']$. The places $(0,t)$ are not thick (the thick places are attached
to the second real places), so they are top or thin places. The last assertion is the definition of $\mathrm{num}(D,L)$.
\end{proof}

\begin{lemma}[Degree of the run weight]\label{pf:onto3:dfka}
If $3\le x$ and $kb'\le b$, then for $j$ in range and $i=0,1$,
\[
\ontothreeDF(\kappa R(j))\le k\bigl(2fb+f(S_c+\eta)\bigr)\qquad\text{and}\qquad\ontothreeDF(\kappa R(j))<\frac{g_i-2}2 .
\]
\formallabel{GlobalTheory.RunDepth.DF_kaR_le}
\formallabel{GlobalTheory.RunDepth.four_sq_lt_of_eighth}
\formallabel{GlobalTheory.RunDepth.sum_fin_le_two_blocks}
\end{lemma}

\begin{proof}
By \lref{GlobalTheory.WeightRows.kappa_entries_le}, which bounds the entries of $\kappa R(j)$ by the sizes, for
every type $\tau$: $\kappa R(j)^+_\tau(i)\le w_j(\tau)+S_c$ for $i<k-1$, $\kappa R(j)^+_\tau(i)\le\eta$ for $i\ge k-1$,
and $\kappa R(j)^-_\tau(i)\le\eta$. The row $a^+$ has $p_j=(k-1)+j$ entries and $a^-$ has $q_j=m-j$ entries, so by
\cref{pf:onto3:cloud},
\[
\sum_i\kappa R(j)^+_\tau(i)+\sum_i\kappa R(j)^-_\tau(i)\le(k-1)(2b+S_c)+j\eta+(m-j)\eta\le k(2b+S_c+\eta),
\]
using $m\le k$. There are $f$ types, which gives the first bound. For the second, $S_c+\eta\le2b$ and $f\le k$ give
$\ontothreeDF(\kappa R(j))\le k\cdot4kb=4k^2b$. The parameter choice provides $Y$ with $(kb)^8\le Y\le g_i$. For real
$k\ge5$, $b\ge1$ and $(kb)^8\le g$ one has $4k^2b<(g-2)/2$: with $y=kb\ge5$, $k^2b\le y^2$ and $8y^2+2<y^8$.
\end{proof}

\subsubsection{Page conditions}

The following definition makes explicit the conditions described in \cref{pf:onto1:hyps,pf:onto1:parconds}. Let $X$ be
a valued form algebra over $B$ at $j$, $\mathbf H=(C_K,C_A,C_1,A_1,C_c,A_c,C_{Rc})$ a record of constants, $C_2,A_2,C_R,
C_T\in\R$, weights $\kappa,w_L,w_F$, $T,n_P,n_W\in\N$ and $h_0\in\R$; put $\delta=\delta_j$, $T'=T+\delta+2$, $w_0=w_F+
\delta w_L$, and
\begin{gather*}
c=\frac{\ontothreeRR(C_{Rc})\ontothreePP(C_c,A_c)\ontothreeDF(w_L)^2}{q_\wedge},\qquad
n_0(j)=r+2c_0(fp_jq_j+1),\qquad \mathrm{PN}=2\mathrm L+\lambda^*+4,\\
\ell_{rd}=\frac{\log|d_E|}{2f},\qquad\mathrm{PU}=fr\bigl(2\ell_{rd}+2\log(fr)+3\bigr),\qquad
\mathrm{PO}=19fr^2\mathrm L+198f^2\mathrm L,\\
D_W(j)=72\cdot5^4\delta^2f^2r^2m^4,\\
\mathrm{sc}(n_W)=\max\bigl(90(fr+1),\,360fm^3,\,2(3r+3+r\log r),\,n_W,\,c_0,\,fr^2+6,\,11\bigr),\\
\mathrm{fl}(n_W)=\max\bigl(16200\pi n_Wm^4u^4,\ u^4n_W+2u^3,\ ru^3\bigr)\quad\text{with }u=2fm,
\end{gather*}
and $\mathrm{nc}(w,\tau)=a^+_\tau(p_j-1)+a^-_\tau(q_j-1)+n_w-r+1$.

\begin{definition}[Page conditions]\label{pf:onto3:pagedef}
\begin{enumerate}[label=(\roman*),leftmargin=*]
\item (Parameter bounds, (H1) for $C_A$.) $x\le fr\le x^{C_A}$; $\log|d_E|\le C_Afr\mathrm L$; $x^{-\eps}\le\delta_2$;
$\log g_i\le9\mathrm L+1$; $|\log\mathrm{Der}|\le x^{\eps}$; $|\log\Omega|\le C_A\mathrm L$; $\log p_{\mathrm{de}},\log
p'_{\mathrm{de}}\le\mathrm L$; $2fm\le\mathrm{rd}_E$.
\item (Field bounds.) $\log[K_0:\Q]\le\mathrm{PO}$ and $\log|d_{K_0}|/[K_0:\Q]\le\mathrm{PO}$.
\item (Domination for $C_P$.) $1000\le C_P$, $C_P^{20}\le z_B$, $\lambda^*,\mathrm{PN},\log(N_{\mathrm{cond}}g_0g_1),\mathrm{PU}
\le C_Pz_B$, $\mathrm{PO}\le C_Pz_B^2$, $n_0(j)\le C_Pz_B^3$, $\mathrm{sc}(n_W)\le C_Pz_B^4$ and $D_W(j)\le C_Pz_B^{10}$, where
$N_{\mathrm{cond}}$ is the conductor integer of the host.
\item (Weight conditions, (H2) for $C_K,C_1$.) $n_\kappa=n_V$; $T'(T'+2r)\le\mathrm{gap}_i(\kappa^\pm_\tau)+2$ for all
non-last $i$; $4T^2+t+1\le h_0$; $h_0/4-t\le\mathrm{nc}(\kappa,\tau)$ and $h_0/2-(t+1)/2\le\kappa^\pm_\tau(\text{last})$ for
all $\tau$; $\ontothreesize(\kappa)\le C_Kkfb$; $n_P$ even, $8fr+4\le n_P\le n_W\le\mathrm{sc}(n_W)$, $fr(r+1)\le n_W$, $n_V=
n_P+Tn_W$; $n(w_L)=n_W$ and the gaps of $w_L$ are $\lfloor(\mathrm{gap}_i(\kappa^\pm_\tau)+2)/T'\rfloor$ for non-last $i$,
$\lfloor(\kappa^+_\tau(p_j-1)-(\delta+2)n_W-2r)/T'\rfloor$ and $\lfloor\kappa^-_\tau(q_j-1)/T'\rfloor$ for the last indices
(without truncation at $0$, unlike \cref{pf:onto3:runweights}(ii)); $n(w_F)=n_P$ and $w_F=\kappa-Tw_L$ entrywise;
$\ontothreesize(w_L)\le C_Kkfb/T$; $0<T$, $\delta\le T$,
$40(\delta+2)\le T$, $(\delta+2)n_W+2r\le T$, $40\le n_W$, $n_0(j)\le T$, $11\delta^2\le T$; the field bounds (ii); and the
domination (iii) for $C_P=C_1$.
\item (Numerical part of (H6) for $C_2,A_2,C_R$.) $\ontothreeRR(C_R)^4\ontothreePP(C_2,A_2)^3\ontothreeDF(w_L)^2\le
q_\wedge$; $Tc_0\ontothreeRR(C_R)\ontothreePP(C_2,A_2)\ontothreeDF(w_L)^2\le q_\wedge$;
$\ontothreePP(C_2,A_2)^4\delta^3kf/(m\delta_2)\le T$; $Tcc_0\le1$; $fm^3\ontothreeRR(C_R)^2c\le\frac12$;
$(\delta+4)\ontothreeDF(w_L)\le q_\wedge-2$.
\item (Size conditions.) $b\le\ontothreesize(\kappa)$; $T\ontothreesize(w_L)\le2\ontothreesize(\kappa)$;
$\ontothreesize(w_F)\le(\delta+2)\ontothreesize(w_L)+5fr^2T$; $\ontothreePP(C_2,A_2)^3T^2\le b$; $\ontothreeRR(C_R)
\ontothreePP(C_2,A_2)\le b$; $2\ontothreeDF(\kappa)+2<g_i$ for $i=0,1$.
\end{enumerate}
The \emph{page hypotheses} $\mathcal H_{\mathrm{page}}$ are (i) for $C_A$, (iv) for $(C_K,C_1)$, (v), (vi) and $T\le
x^{C_T}$. \emph{Page parameters} for $(C_2,A_2)$ are a parameter choice for $(D,L)$ with parameters $T_p$ such that $L$
satisfies the size conditions and $(C_P,e_P)(T_p)=(C_2,A_2)$. The \emph{page conditions} $\mathcal C_{\mathrm{page}}$ are
$\mathcal H_{\mathrm{page}}$ with $C_R=7$, page parameters for $(C_2,A_2)$, $C_1=C_c=C_2$, $A_1=A_c=A_2$, $C_{Rc}=7$ and
$12\le A_2$. The \emph{strengthened page conditions} $\mathcal C^+_{\mathrm{page}}$ are $\mathcal C_{\mathrm{page}}$ together
with $\mathrm{fl}(n_W)\le\ontothreePP(C_2,A_2)$, $f^2\le\mathrm{rd}_E$, $8\le r$, $3\le\ontothreeDF(w_L)$ and
$4C_K\le\ontothreePP(C_2,A_2)^2f$. The \emph{page support} for a scalar quadruple $S$ consists of
$\mathrm{fl}(n_W)\le\ontothreePP(C_2,A_2)$, $\frac12\ontothreeDF(w_L)\log(g_0g_1)\le\log|\sigma(s_W)|$ for all
$\sigma$, $f^2\le\mathrm{rd}_E$, $8\le r$ and $3\le\ontothreeDF(w_L)$. The \emph{page-row bounds} at a weight $w$ are
$\log g_i\le9\mathrm L+1$, $\ell_{rd}\le2\mathrm L$, $\log d\le2\mathrm L$, $\log\Nm\mathfrak c_K\le2\mathrm L$, $\log f\le2
\mathrm L$, $\log p_{\mathrm{de}},\log p'_{\mathrm{de}}\le\mathrm L$, $\log(1+\ontothreeDF(w))\le2\mathrm L$ and $\log n_w\le
\mathrm L$. Finally, for $P_\Lambda\in\R$ and $p\in\N$,
\begin{multline*}
\mathrm{Gc}(P_\Lambda,p)=P_\Lambda+3\mathrm L+2\max\Bigl(\log[K_0:\Q]+\frac{\log|d_{K_0}|}{[K_0:\Q]},\ \log[E:\Q]+
\frac{\log|d_E|}{[E:\Q]}\Bigr)\\+2fr(\lambda^*+2\log p)+31z_B .
\end{multline*}
\formallabel{Formalism.MainNotions.WeightAlgebra.ParA}
\formallabel{Formalism.MainNotions.WeightAlgebra.OPO}
\formallabel{Formalism.MainNotions.WeightAlgebra.DomAt}
\formallabel{Formalism.MainNotions.WeightAlgebra.SetC}
\formallabel{Formalism.MainNotions.WeightAlgebra.H3}
\formallabel{Formalism.MainNotions.Gaps.FamilyHypotheses.FourthHypothesisSizes}
\formallabel{Formalism.MainNotions.Gaps.FamilyHypotheses.HypPage}
\formallabel{Formalism.MainNotions.WeightAlgebra.PageAt}
\formallabel{Formalism.MainNotions.WeightAlgebra.TieC}
\formallabel{GlobalTheory.FinalJoin.PageRun}
\formallabel{GlobalTheory.RunJoin.PageOK}
\formallabel{Formalism.MainNotions.Gaps.EtaleWeightGaps.PageSupp}
\formallabel{Formalism.MainNotions.Gaps.EtaleWeightGaps.flipOneThree}
\formallabel{ReciprocityTheory.ChartRestriction.G.PullbackSections.SupBoundPageRows}
\formallabel{Formalism.MainNotions.Gaps.FamilyHypotheses.SupBoundPageRows}
\formallabel{Formalism.MainNotions.WeightAlgebra.Gc}
\end{definition}

The two formal versions of the page-row bounds have the same content. The page hypotheses do not involve the points,
the valuations or the theta forms of $X$, only $B$ (through the field bounds) and the weights.

\begin{lemma}[Page rows]\label{pf:onto3:pagerows}
Assume $\mathcal H_{\mathrm{page}}$.
\begin{enumerate}[label=(\roman*)]
\item If the page-row bounds hold at $\kappa$, then $2\le\mathrm L$, $2\le f$, $6\le r$, $\lambda^*\ge0$, and $\delta_j\ge1$
if $j\ge1$.
\item If page parameters for $(C_2,A_2)$ exist and $4C_K\le\ontothreePP(C_2,A_2)^2f$, then the page-row bounds hold at
$\kappa$ and $16\le f$.
\item If the page-row bounds hold at $\kappa$, $\mathrm L\ge0$, $\ontothreeDF(w)\le\ontothreeDF(\kappa)$ and $n_w\le n_\kappa$,
then they hold at $w$.
\end{enumerate}
\formallabel{GlobalTheory.page_two_le_L_and_signs}
\formallabel{GlobalTheory.pageRows_of_pageAt}
\formallabel{GlobalTheory.pageRows_of_le}
\end{lemma}

\begin{proof}
(i) $m\ge2$ and $k\ge m+3$ give $r=k+m-1\ge6$, and $f=2n\ge2$. If $\mathrm L<2$, the page rows give $\log f\le4$ and
$\log n_\kappa\le4$, so $f\le e^4\le81$ and $r\le n_\kappa=n_V\le81$; then $z_B=fr\mathrm L\le2\cdot81^2<1000^{20}\le
C_1^{20}$, contradicting the domination (iii) for $C_1$. $\lambda^*\ge0$ because $g_0g_1\ge1$ and $f^2qd\ge1$. If
$j\ge1$ then $p_j=k+j-1\ge5$, so $\binom{p_j}2\ge1$ and $\delta_j\ge f\ge1$.
(ii) Let $(Q,T_p)$ be page parameters. The parameter choice gives $m^4\le f$ (hence $f\ge16$, as $m\ge2$), $\log f\le
\mathrm L/72$, $\log r\le\mathrm L/300$, $\log t\le\mathrm L/300$, $\log(Q_Sqd)\le\mathrm L/72$, $\log\Nm\mathfrak c_K\le2\mathrm L$
and $\ell_{rd}\le2\mathrm L$ (here $\ell_{rd}=\log\mathrm{rd}_E$), and the ranks satisfy $k\le r\le t$. Together with (i) of
\cref{pf:onto3:pagedef} ($\log g_i\le9\mathrm L+1$, $\log p_{\mathrm{de}},\log p'_{\mathrm{de}}\le\mathrm L$) and $d\le Q_Sqd$,
these give all page rows except the last two. By (iv) and (vi), $1+\ontothreeDF(\kappa)\le\ontothreesize(\kappa)\le C_Kkfb$,
and $C_K\le\ontothreePP^2f$, $k\le r$, so $1+\ontothreeDF(\kappa)\le\ontothreePP^2f^2rb$ with $\ontothreePP=\ontothreePP(C_2,A_2)$.
From $\ontothreePP^3T^2\le b$ and $T\ge1$ we get $3\log\ontothreePP\le\mathrm L$, hence $\log(1+\ontothreeDF(\kappa))\le
\frac23\mathrm L+\frac2{72}\mathrm L+\frac1{300}\mathrm L+\mathrm L\le2\mathrm L$. Finally $n_\kappa=n_V=r+t\le2t\le t^2$, so
$\log n_\kappa\le2\log t\le\mathrm L$.
(iii) Only the last two page rows involve the weight, and $\log$ is non-decreasing on $(0,\infty)$; if $n_w=0$ the
bound $\log0=0\le\mathrm L$ is the assumption $\mathrm L\ge0$.
\end{proof}

\begin{lemma}[Weights below $\kappa$]\label{pf:onto3:below}
Assume $j\ge1$ and $\mathcal H_{\mathrm{page}}$. Then:
\begin{enumerate}[label=(\roman*)]
\item for $w\in\{w_L,w_F,w_0,\kappa\}$: $\ontothreeDF(w)\le\ontothreeDF(\kappa)$, $n_w\le n_\kappa$, $r\le n_w$, and
$a^+_\tau(p_j-1)\ge1$ for all $\tau$;
\item $8fr+4\le n_w$ for $w\in\{w_F,w_0,\kappa\}$;
\item for every weight $w_2$ with $w_2+w_F=(\delta+4)w_L$: $\ontothreeDF(w_2)\le\ontothreeDF(\kappa)$, $n_{w_2}\le n_\kappa$
and $\ontothreesize(w_2)\le(\delta+4)\ontothreesize(w_L)$;
\item for every weight $w_1$ with $a^\pm(w_1)=2a^\pm(w_L)$ and $n_{w_1}+2=2n_W$: $\ontothreeDF(w_1)\le\ontothreeDF(\kappa)$,
$n_{w_1}\le n_\kappa$ and $\ontothreesize(w_1)\le2\ontothreesize(w_L)$.
\end{enumerate}
\formallabel{GlobalTheory.RunWeights.weights_below}
\formallabel{GlobalTheory.RunWeights.size_le_mul}
\end{lemma}

\begin{proof}
From (iv) of \cref{pf:onto3:pagedef}: $\kappa=w_F+Tw_L$ entrywise, $n_\kappa=n_P+Tn_W$ with $n_{w_L}=n_W$, $n_{w_F}=n_P$;
$T\ge40(\delta+2)\ge80$, so $\delta+4\le T$; $f\ge1$ (from $C_1^{20}\le z_B$, so $z_B>0$), hence $r\le8fr+4\le n_P\le n_W$.
The last entries satisfy $\kappa^+_\tau(p_j-1)\ge h_0/2-(t+1)/2\ge2T^2$. The last gap of $w_L$ is
$\lfloor(\kappa^+_\tau(p_j-1)-(\delta+2)n_W-2r)/T'\rfloor$, and $\kappa^+_\tau(p_j-1)-(\delta+2)n_W-2r\ge2T^2-T\ge T'$
because $(\delta+2)n_W+2r\le T$, $\delta\le T$ and $T\ge80$; hence $w_L^+(\tau,p_j-1)\ge1$, and
$T'w_L^+(\tau,p_j-1)\le\kappa^+_\tau(p_j-1)-(\delta+2)n_W-2r$, so $w_F^+(\tau,p_j-1)=\kappa^+_\tau(p_j-1)-Tw_L^+(\tau,p_j-1)
\ge(\delta+2)n_W\ge1$. For $w=w_L,w_F,w_0$ the entries are bounded by those of $\kappa=w_F+Tw_L$ (for $w_0=w_F+\delta
w_L$ use $\delta\le T$), and the integers $n_w$ by $n_\kappa=n_P+Tn_W$; this gives (i) and (ii), the last entry of $w_0$
being at least that of $w_F$. For (iii), $w_2\le(\delta+4)w_L$ entrywise and $n_{w_2}\le(\delta+4)n_W$, and $(\delta+4)
w_L\le Tw_L\le\kappa$; similarly (iv) with $2\le T$. The size bounds follow from the fact that if $x\le cz$ entrywise and
$n_x\le cn_z$ with $c\ge1$, then $\ontothreesize(x)\le c\ontothreesize(z)$, since $\ontothreesize$ is $1$ plus a linear
function of the entries and of $n$ with non-negative coefficients.
\end{proof}

\begin{lemma}[Generator constants]\label{pf:onto3:gc}
Let $w$ be a weight with $n_w>0$, $p\in\N$, and assume $\mathrm L\ge2$. With $|w|=\ontothreeDF(w)$ and $\mathrm{Gc}=
\mathrm{Gc}(1009z_B,p)$:
\begin{gather*}
140r\mathrm L|w|+10\mathrm L|w|+200frn_w\mathrm L\le400z_B\ontothreesize(w),\\
2(\lambda^*+2\log p)\ontothreesize(w)+400z_B\ontothreesize(w)+\lambda^*\ontothreesize(w)\le\mathrm{Gc}\,\ontothreesize(w),\\
(\lambda^*+2\log p)\ontothreesize(w)+400z_B\ontothreesize(w)+3frn_w\mathrm L+5\log[K_0:\Q]+\frac{\log|d_{K_0}|}{[K_0:\Q]}+2
\le\mathrm{Gc}\,\ontothreesize(w).
\end{gather*}
\formallabel{GlobalTheory.RunIn.sup_bounds_le_Gc}
\formallabel{GlobalTheory.RunIn.gc_arith}
\end{lemma}

\begin{proof}
Put $s=\ontothreesize(w)$, $F=fr$ and let $M\ge0$ be the maximum in $\mathrm{Gc}$; then $\mathrm{Gc}\,s=1040Fs\mathrm L+3\mathrm Ls+
2Ms+2F\lambda^*s+4F(\log p)s$. We have $f\ge2$, $r\ge6$, so $F\ge12$ and $r\le F$; $s\ge|w|+Fn_w+1$, so $|w|,n_w,Fn_w\le s$
and $s\ge13$. First inequality: $140r\mathrm L|w|+10\mathrm L|w|+200F n_w\mathrm L\le(140+10+200)F\mathrm Ls\le400F\mathrm Ls$.
Second: $3\lambda^*s+4(\log p)s+400F\mathrm Ls\le2F\lambda^*s+4F(\log p)s+1040F\mathrm Ls$, since $F\ge12$. Third: $3Fn_w\mathrm L\le3F\mathrm Ls$, $2\le3\mathrm Ls$, and $5\log[K_0:\Q]+\log|d_{K_0}|/[K_0:\Q]\le5M\le2Ms$ (the two
terms are non-negative with sum at most $M$, and $s\ge13$); the remaining terms are absorbed as in the second inequality.
\end{proof}

\begin{lemma}[Theta run parameters at the record]\label{pf:onto3:runrec}
Assume $\mathcal H_{\mathrm{page}}$, the ties of constants $(\mathrm K)$ of \cref{pf:onto1:hyps}, page parameters for
$(C_2,A_2)$ and the page support for a scalar quadruple $S$. Then:
\begin{enumerate}[label=(\roman*)]
\item the run parameters $(q,R_0,m)=\bigl(g_0g_1,\sqrt{90(fr+\ontothreeDF(w_L))},m\bigr)$ are admissible for every
principal torus datum of rank $r$ with a type of dimension at most $m$ and the degree $\ontothreeDF(w_L)$, i.e.\ for
every such datum there is a run parameter with these values;
\item $8100\pi n_{w_L}m^4(2fm)^4\mathrm{rd}_E^{7}/(g_0g_1)\le c$;
\item $p_{14}\le\frac12\bigl(frn_W\log(g_0g_1)+n_W(\log24+fm\log(4fm))+(r+2)\log2\bigr)$, where
$p_{14}=\frac12\bigl(\log\pi+n_{w_L}(\log24+fm\log(4fm))+(r+2)\log2\bigr)$;
\item $fr(r+1)\le2n_{w_L}$.
\end{enumerate}
\formallabel{GlobalTheory.RunParametersAtRecord.run_at_record}
\end{lemma}

\begin{proof}
By $(\mathrm K)$, $C_R=7$, $C_1=C_2$ and $A_2\ge12$, so $\ontothreeRR(C_R)=\mathrm{rd}_E^7$. Since $z_B\ge1$ and $A_2\ge4$, the
domination (iii) for $C_1=C_2$ gives $\mathrm{sc}(n_W)\le C_2z_B^4\le\ontothreePP$ ($\ontothreePP=\ontothreePP(C_2,A_2)$), in
particular $90(fr+1)\le\ontothreePP$. The page support gives $\mathrm{fl}(n_W)\le\ontothreePP$, $f^2\le\mathrm{rd}_E$, $8\le r$
and $3\le\ontothreeDF(w_L)$; the page parameters give $m^4\le f$, and (i) of \cref{pf:onto3:pagedef} gives
$\mathrm{rd}_E\ge2fm\ge16$. From the first inequality of (v), $\ontothreeRR(7)^4\ontothreePP^3\ontothreeDF(w_L)^2\le q_\wedge$;
as $(\delta+4)\ontothreeDF(w_L)\le q_\wedge-2$, $q_\wedge\ge2$, so $g_0,g_1\ge2$ and $q_\wedge\le g_0g_1$.
(i) By \lref{GlobalTheory.RunParametersOfLower.runOK_of_le}, which states that the run parameters
$(q,\sqrt{90(fr+\mathrm{DF})},m)$ are admissible as soon as $1\le\mathrm{DF}$, $2\le fr$, $1\le m$, $m^4\le f$, $f^2\le
\mathrm{rd}_E$, $16\le\mathrm{rd}_E$ and $\mathrm{rd}_E^{28}\cdot90(fr+\mathrm{DF})\cdot\mathrm{DF}\cdot r\le q$, it suffices to verify
the last inequality; it follows from $\mathrm{rd}_E^{28}\ontothreePP^3\ontothreeDF(w_L)^2\le q_\wedge\le g_0g_1$ and
$90(fr+\ontothreeDF(w_L))\ontothreeDF(w_L)r\le\ontothreePP^3\ontothreeDF(w_L)^2$, which holds because $90(fr+1)\le
\ontothreePP$, $r\le\ontothreePP$ and $\ontothreeDF(w_L)\ge1$.
(ii) $n_{w_L}=n_W$ and $8100\pi n_Wm^4(2fm)^4\le\mathrm{fl}(n_W)\le\ontothreePP\le\ontothreePP\ontothreeDF(w_L)^2$; multiplying
by $\mathrm{rd}_E^7$ and dividing by $g_0g_1\ge q_\wedge$ gives (ii), since $c=\mathrm{rd}_E^7\ontothreePP\ontothreeDF(w_L)^2/
q_\wedge$.
(iii) $\log\pi\le\log4\le\log(g_0g_1)\le frn_W\log(g_0g_1)$, since $frn_W\ge1$. (iv) is $fr(r+1)\le n_W$ of (iv) of
\cref{pf:onto3:pagedef}.
\end{proof}

\begin{lemma}[The numerical threshold]\label{pf:onto3:pageok}
For every function $Z_0\colon\R^2\to\R$ put $x_0(a,b)=\max(Z_0(a,b),4b)$. Let $j$ be in range and assume
$\mathcal C_{\mathrm{page}}$ and $x_0(C_A,C_K)\le x$. Then $\mathcal C^+_{\mathrm{page}}$ holds and $Z_0(C_A,C_K)\le z_B$.
\formallabel{GlobalTheory.FinalJoin.pageOK_of_pageRun_beyond}
\formallabel{GlobalTheory.FinalJoin.indexBoundConstant_le_pagePolynomial_of_runConditions}
\formallabel{GlobalTheory.FinalJoin.indexBound_threeTerm_constant_le_real}
\formallabel{GlobalTheory.FinalJoin.three_le_DF_of_pageRun}
\formallabel{GlobalTheory.FinalJoin.three_le_DF_of_SetC}
\formallabel{GlobalTheory.FinalJoin.pageRun_sq_le_rd_and_eight_le_rank}
\end{lemma}

\begin{proof}
Write $\ontothreePP=\ontothreePP(C_2,A_2)$. By $\mathcal C_{\mathrm{page}}$, $C_1=C_2$, so the domination gives $1000\le C_2$
and $C_2^{20}\le z_B$; in particular $z_B\ge2000$. As $\mathrm L\ge1$, $fr\le z_B$, and by (i) of
\cref{pf:onto3:pagedef} $x\le fr\le z_B$. Since $A_2\ge12\ge1$ and $C_2\ge1$, $z_B\le\ontothreePP$, so $1\le\ontothreePP\le
\ontothreePP^2f$ ($f\ge8$ by the ranks).

$\mathrm{fl}(n_W)\le\ontothreePP$: put $Z=z_B$, $u=2fm\ge1$; then $f,r\le Z$, $m^4\le f\le Z$ (page parameters), $u^4=16f^4m^4
\le16Z^5$, and $n_W\le\mathrm{sc}(n_W)\le C_2Z^4$. Elementary estimates with $Z\ge2000$ and $\pi\le4$ give $16200\pi n_Wm^4u^4
\le1036800C_2Z^{10}$, $u^4n_W+2u^3\le16C_2Z^9+32Z^5$ and $ru^3\le16Z^6$, all at most $C_2Z^{12}\le C_2Z^{A_2}$.

$f^2\le\mathrm{rd}_E$ is a property of page parameters, and $8\le r$ follows from $m^4\le f$, $m\ge2$ and the relations
between the ranks ($f\le k$, $r=k+m-1$). For $3\le\ontothreeDF(w_L)$: fix a type $\tau$ and the last index $i=p_j-1$; then
$\kappa^+_\tau(i)\ge2T^2$ (\cref{pf:onto3:below}), $(\delta+2)n_W+2r\le T$ and $T\ge80$, so
$\kappa^+_\tau(i)-(\delta+2)n_W-2r\ge2T^2-T\ge3T'$, and the last gap $w_L^+(\tau,i)=\lfloor(\cdots)/T'\rfloor\ge3$; hence
$\ontothreeDF(w_L)\ge3$. Finally $4C_K\le x\le z_B\le\ontothreePP\le\ontothreePP^2f$ and $Z_0(C_A,C_K)\le x\le z_B$, by the
choice of $x_0$.
\end{proof}

\begin{proposition}[Page conditions at the record]\label{pf:onto3:pagerun}
There is a constant $C_K$ (namely $C_K=10$) with the following property. Let $T_p$ be parameters of a parameter choice for
$(D,L)$ such that $L$ satisfies the size conditions, the parameter $C_{nW}$ of the choice is $10$, $\log g_i\le9\mathrm L+1$,
$10^{60}\le x$, $1000\le C_P$ and $C_P^{20}\le z_B$. Then for every form provider $B$ with the field bounds (ii) of
\cref{pf:onto3:pagedef}, every $j$ in range and every valued form algebra $X$ over $B$ at $j$, the page conditions
$\mathcal C_{\mathrm{page}}$ hold for $X$, the constants $\mathbf H=(C_K,20,C_P,e_P,C_P,e_P,7)$, $(C_2,A_2)=(C_P,e_P)$,
$C_T=C_t(L)$, the run weights $(\kappa,w_L,w_F)=(\kappa R(j),lR(j),fR(j))$, $(T,n_P,n_W)=(T_O,n_P,n_W)$ and $h_0=h_0(L)-4$.
\formallabel{GlobalTheory.FinalJoin.RecordRun.exists_pageRun_record_sub_four}
\formallabel{GlobalTheory.FinalJoin.RecordRun.pageRun_record}
\end{proposition}

\begin{proof}
From $1000^{20}=10^{60}$ and $C_P\ge1000$ we get $z_B\ge10^{60}$, and $x\ge3$. We verify the components of
\cref{pf:onto3:pagedef}, using the following numerical consequences of the definitions of parameter choices and of the
size conditions, proved in the formal development:
\lref{Formalism.MainNotions.Gaps.EtaleWeightGaps.Num.rn_of} (explicit inequalities between $r,f,m,k,b,\mathrm{rd}_E$,
$T_O$ and $\ontothreePP=\ontothreePP(C_P,e_P)$; among them, for an explicit $s_b$ with $s_b^2=b$: $T_O\le s_b$,
$\ontothreePP\le s_b$ and $\mathrm{rd}_E^7\le s_b$);
\lref{Formalism.MainNotions.Gaps.EtaleWeightGaps.Num.Pm.domAt_of} and \lref{GlobalTheory.RunJoin.InhAll.domAt_index}
(the domination (iii) at the index $m$ for some constant, hence at every $j$ in range, and, by monotonicity in the
constant, for $C_P$, since $C_P^{20}\le z_B$);
\lref{GlobalTheory.FinalJoin.RecordRun.setC_of} (given the domination for $C_P$ and the field bounds, the weight
conditions (iv) hold for the run weights with $C_K=10$ and $C_1=C_P$);
\lref{GlobalTheory.FinalJoin.RecordRun.DF_wL_le} ($\ontothreeDF(lR(j))\le2frb$) and
\lref{Formalism.MainNotions.Gaps.EtaleWeightGaps.Num.qq_ge} ($q_\wedge\ge(kb)^8$), which give, by the elementary
inequalities \lref{GlobalTheory.RunJoin.Inh.nums}, the first two inequalities of (v) and the inequalities
$T_O\cdot1250r\cdot\mathrm{rd}_E^7\ontothreePP\ontothreeDF(lR(j))^2\le q_\wedge$ and $2fm^3\mathrm{rd}_E^{21}\ontothreePP
\ontothreeDF(lR(j))^2\le q_\wedge$, which are $Tcc_0\le1$ and $fm^3\ontothreeRR(7)^2c\le\frac12$ (recall $c_0=1250r$);
\lref{Formalism.MainNotions.Gaps.EtaleWeightGaps.Num.third_hypothesis_of_params} and
\lref{Formalism.MainNotions.Gaps.EtaleWeightGaps.Num.fourth_hypothesis_of_params} (the third inequality of (v), with
$\delta_j=D_p$, and $\ontothreePP^3T_O^2\le b$);
\lref{GlobalTheory.FinalJoin.RecordRun.b_le_size}, \lref{GlobalTheory.FinalJoin.RecordRun.T_wL_le} and
\lref{GlobalTheory.FinalJoin.RecordRun.wF_size_le} ($b\le\ontothreesize(\kappa R(j))$, $T_O\ontothreesize(lR(j))\le
\ontothreesize(\kappa R(j))+T_O$, and the bound (vi) for $\ontothreesize(fR(j))$); and
\lref{Formalism.MainNotions.Gaps.EtaleWeightGaps.Num.Pm.parA_of} (the parameter bounds (i) for $C_A=20$).
The remaining inequalities follow: $(\delta_j+4)\ontothreeDF(lR(j))\le q_\wedge-2$ from $\ontothreeDF(lR(j))\le2frb$,
$\delta_j=D_p$ and $q_\wedge\ge(kb)^8$ (an inequality of the first cited result); $T_O\ontothreesize(lR(j))\le2
\ontothreesize(\kappa R(j))$ from $T_O\le s_b\le b\le\ontothreesize(\kappa R(j))$; $\ontothreeRR(7)\ontothreePP\le s_b^2=b$; and
$2\ontothreeDF(\kappa R(j))+2<g_i$ from \cref{pf:onto3:dfka}, which applies since $kb'\le b$ for parameter choices with
$x\ge3$ (\lref{GlobalTheory.FinalJoin.RecordRun.k_b'_le}). The bound $T_O\le x^{C_t}$ is a property of parameters.
The page parameters are $(Q,T_p)$ themselves, and the equalities of constants hold by the choice of $\mathbf H$, with
$e_P\ge12$.
\end{proof}

\begin{lemma}[Field smallness]\label{pf:onto3:fieldsmall}
For a top structure $T$ with laws, the field smallness of $(D,L)$ (the field bounds (ii) of \cref{pf:onto3:pagedef} for
the coefficient field $F_{\OO}$ of the datum) holds if and only if the field bounds hold for $B_T$.
\formallabel{GlobalTheory.FinalJoin2.fieldSmall_iff_opo}
\end{lemma}

\begin{proof}
By the pins of the laws, the coefficient field of $T$ is the field $F_{\OO}$ of the datum
(\lref{GlobalTheory.FinalJoin2.FO_eq_coeff}), and the coefficient field of $B_T$ is that of $T$. Isomorphic number
fields have the same degree and the same discriminant.
\end{proof}

\begin{lemma}[Parameter conditions from laws]\label{pf:onto3:parok}
Let $J$ be a core with its laws and $T$ a top structure with the laws of top structures and the short laws, whose core is
$J$ and is canonical. Let $\mathrm{EX}$ be an error datum which is half-compatible with the core, and $p$ a prime at which
the field conditions hold (in particular $p\in\{p_{\mathrm{de}},p'_{\mathrm{de}}\}$). If every auxiliary datum
$d_\delta(j,n)$ of $T$ with $j$ in range has level $1$ and prime $p$, then the parameter conditions
$\mathcal C_{\mathrm{par}}(T,\mathrm{EX},p,j)$ of \cref{pf:onto1:parconds} hold for every $j$ in range.
\formallabel{GlobalTheory.FinalJoin2.liftIndexParameterConditions_of_coreAndShortLaws}
\formallabel{GlobalTheory.PtTie.halfEX_iff_oriented}
\formallabel{GlobalTheory.TrivLevel.badN_pos_iff}
\end{lemma}

\begin{proof}
We go through the conditions. The short laws are assumed, and the levels of the shortening data are $1$ by
hypothesis. The components of the unfolding at $j$ share their group and character data: the model laws of $J$ identify
the group and the character of every component with those of a model, and $J$ is the core of $T$. The bad integer
$N_{\mathrm{bad}}=Q_Sqd\prod_pp$ (product over the level primes of the host) is positive, because $Q_S,q,d\ge1$ and the
field conditions state that $0$ is not a level prime (conversely $N_{\mathrm{bad}}=0$ if $0$ is a level prime). All
shortening data at $j$ have prime $p$ by hypothesis. Orientation is equivalent to half-compatibility: by
\lref{GlobalTheory.Survey.PiTwistNegation.halfEX_iff_le} the error datum is half-compatible with the core if and only
if the level ideal is contained in the principal ideals generated by the conjugates $\bar\pi_i$ of the elements $\pi_i$
of $\mathrm{EX}$, and this is the definition of orientation. Finally, the pins give level $M=g_0g_1$ and $A=1$, so $M$
divides $g_0g_1p$, and \lref{GlobalTheory.FinalJoin.FieldAlg.fnAlgClosed_of_pins}, which derives the algebraic
closedness of the coefficient field in the function fields of the moduli spaces from the pins, canonicity and the curve
part of the field conditions for a multiple of the level, gives the last condition.
\end{proof}

\begin{lemma}[Numerical conditions along families]\label{pf:onto3:pagenum}
For every $x_1\in\R$ the following holds closedly over $\sfname{FamHyp}$: there is $y$ such that for all $x>x_1(\mathcal F)$
with $x\ge y$ and all $i\in I_x$, the size data $L=L_{x,i}$, datum $D=D_{x,i}$ and parameters $T_p$ of the family at
$(x,i)$ satisfy: the size conditions, $C_{nW}=10$, $\log g_i\le9\mathrm L+1$, $1000\le C_P$, $C_P^{20}\le z_B$ and $x_1\le
x$ (the \emph{page-run conditions}); $3\le x$; $kb'\le b$; $3\le x^{\eps_z/5}$; and $8fr+4\le n_V$.
\formallabel{GlobalTheory.T.runNumericalConditions_hold_along_family}
\formallabel{GlobalTheory.OntoNum.pageNum_along}
\end{lemma}

\begin{proof}
By \lref{GlobalTheory.T.indexBoundParameterRequirements_eventually_family}, there are $\eps_1>0$ and, for every
$\eps\in(0,\eps_1]$, thresholds $K$ such that every family of $\sfname{FamHyp}\cap\sfname{FamP}(\eps)$ which is $K$-large
satisfies, for all large $x$, a list of parameter requirements, among them $kb'\le b$ and $11\le x^{\eps_z/5}$, which
gives $3\le x^{\eps_z/5}$ and $3\le x$ (\lref{GlobalTheory.IndexAssembly.three_le_scale_of_pow_bound}). Take the
thresholds $\max(K,K_0)$ with $K_0$ requiring the pre-constant $C_P$
to be at least $1000$. For a family large with respect to these thresholds, $C_P\ge1000$, and by
\lref{GlobalTheory.FinalRows.rows_of_pageFix} (eventually along a family of $\sfname{FamHyp}$ with $C_P\ge1000$: the size
conditions, $C_{nW}=10$, $\log g_i\le9\mathrm L+1$, $C_P^{20}\le z_B$ and $x(L_{x,i})=x$) the remaining page-run conditions hold
for large $x$, with $x_1\le x$ as soon as $x\ge x_1$. The inequality $8fr+4\le n_V$ is a property of the parameters. Taking
$y$ the maximum of the thresholds and $x_1$ gives the claim.
\end{proof}

\subsubsection{Theta lattices of records}

\begin{lemma}[Theta lattices are integral]\label{pf:onto3:thetamsup}
Let $X$ be a valued form algebra over $B$ at $j$, $S$ a scalar quadruple and $\kappa,w_L,w_F$ weights, $T\in\N$, with
\begin{enumerate}[label=(\alph*)]
\item $W_{s_W}(w_L)\subseteq\mathcal M_{w_L}$, $\mathrm{Th}_{s_F}(w_F)\subseteq\mathcal M_{w_F}$, $\mathrm{Th}_{s_0}(w_0)\subseteq
\mathcal M_{w_0}$ and $\mathrm{Th}_{s_K}(\kappa)\subseteq\mathcal M_\kappa$, where $w_0=w_F+\delta_jw_L$;
\item $\kappa=w_F+Tw_L$ entrywise and $n_\kappa=n_{w_F}+Tn_{w_L}$, and $\delta_j\le T$.
\end{enumerate}
Then $\Theta(S,\kappa,w_L,w_F,T)\subseteq\iota(\mathcal M_\kappa)$.
\formallabel{Formalism.MainNotions.Gaps.T0.theta_le_msup}
\end{lemma}

\begin{proof}
If $F\in\mathcal M_w$ and $G\in\mathcal M_{w'}$, then $\iota_w(F)\iota_{w'}(G)=\iota_{w+w'}(H)$ for some $H$
(\cref{pf:onto1:PJ}(i)), and $v_{\mathfrak l,y}(\iota H)=v_{\mathfrak l,y}(\iota F)+v_{\mathfrak l,y}(\iota G)\ge0$ at all points, so
$H\in\mathcal M_{w+w'}$ by the integrality condition. Hence $\iota(\mathcal M_w)\cdot\iota(\mathcal M_{w'})\subseteq\iota(
\mathcal M_{w+w'})$, and by induction $\iota(\mathcal M_{w_L})^n\cdot\iota(\mathcal M_w)\subseteq\iota(\mathcal M_{nw_L+w})$ for
all $n$ (products of $\OO_K$-submodules are monotone). By (a), the three summands of $\Theta$ are contained in
$\iota(\mathcal M_{Tw_L+w_F})$, $\iota(\mathcal M_{(T-\delta_j)w_L+w_0})$ and $\iota(\mathcal M_\kappa)$, and by (b) both
$Tw_L+w_F$ and $(T-\delta_j)w_L+w_0$ equal $\kappa$.
\end{proof}

\begin{lemma}[Scaled theta forms are integral]\label{pf:onto3:thint}
Let $T$ be a top structure with laws, $j$ in range, $\OO\to\mathit{LL}$ injective and $\mathrm{EX}$ an error datum for which
$T$ is \emph{oriented}, i.e.\ the level ideal of $T$ is contained in $\bar\pi_i\OO_E$ for $i=0,1$. Then
$W_{s(w)}(w)=s(w)\mathcal T(w)\subseteq\mathcal M_w$ for every weight $w$, where $\mathcal T(w)$ are the theta forms of
$X_{T,j}$.
\formallabel{GlobalTheory.T.scaledThetaLattice_le_supLattice_every_weight}
\formallabel{GlobalTheory.PlaceTop.thInt_of_oriented}
\formallabel{GlobalTheory.RunIn.ThInt}
\formallabel{GlobalTheory.PtTie.Oriented}
\end{lemma}

\begin{proof}
By \cref{pf:onto3:sof}, $s=s(w)\ne0$. By the integrality condition of $X_{T,j}$ it suffices to show
$v_{\lambda,y}(\iota_w F)\ge-v_\lambda(s)$ for every $F\in\mathcal T(w)$, every place $\lambda$ of $K_0$ and every point $y$
over $\lambda$, and since $v_{\lambda,y}$ satisfies the ultrametric inequality and $v_{\lambda,y}(aF)=v_\lambda(a)+v_{\lambda,y}
(F)\ge v_{\lambda,y}(F)$ for $a\in\OO_{K_0}$, it suffices to show this for the theta rules $g$ of $T$ at $(j,w)$, which span
$\mathcal T(w)$. For these, \lref{GlobalTheory.T.thetaGenerator_place_valuation_bound} (a valuation bound for theta
rules at integral points of oriented top structures) gives $v_y(g)\ge-\bigl(\ell_{\mathrm{ram}}(\lambda,w)+v_\lambda(
\mathrm{nn}(w))\bigr)$, where $\ell_{\mathrm{ram}}(\lambda,w)=\sum_{p\mid N_{\mathrm{cond}}}v_\lambda(p)v^H_d(p)
\ontothreeDF(w)/2$, and \lref{GlobalTheory.T.ramified_prime_loss_le_valuation_prime_power} gives $\ell_{\mathrm{ram}}
(\lambda,w)\le v_\lambda\bigl(\mathrm{pp}(\ontothreeDF(w),c_{1/2}(w),c_{1/2}(w))\bigr)$. As $v_\lambda(s)=v_\lambda(\mathrm{nn}
(w))+v_\lambda(\mathrm{pp}(\cdots))$, the claim follows.
\end{proof}

\begin{lemma}[Principal congruence subgroups]\label{pf:onto3:gamat}
Let $T$ be a core with its laws. Then $T$ has the \emph{congruence property} $(\Gamma)$: for every $j$ in range, every good
embedding $\sigma$, every component $k$ of the uniformisation at $\sigma$ and every $\gamma\in\GL_r(E)$ preserving the
hermitian form and the lattice $\Lambda_k$ of the component with $\gamma x-x\in M\Lambda_k$ for all $x\in\Lambda_k$ ($M$ the
level), $\gamma$ lies in the arithmetic group $\Gamma_k$ of the component.
\formallabel{GlobalTheory.RunIn.gamAt_of_lawsCore}
\formallabel{GlobalTheory.T.coreArithmeticGroup_contains_principalCongruence}
\formallabel{GlobalTheory.PrincCore.hasPrinc_core}
\end{lemma}

\begin{proof}
By \lref{GlobalTheory.unf_group_contains_congruence_upto_rootsOfUnity} (for a unitary uniformisation with $A=1$ and
$r>0$), there are $\gamma'\in\Gamma_k$ and a root of unity $u\in E$ with $\gamma'=u\gamma$. The pins of the laws give $A=1$
and $M\ge3$ (the prime $2$ is good, so the level is at least $3$), and the uniformity of the dictionary and the unit
property of the laws allow \lref{GlobalTheory.PrincCore.root_eq_one} (a root of unity $u$ with $u\gamma\in\Gamma_k$ for
$\gamma\equiv1$ modulo $M\ge3$ is $1$). Hence $\gamma=\gamma'\in\Gamma_k$.
\end{proof}

\begin{proposition}[Lower bound for integral forms]\label{pf:onto3:lamw}
There is $C_0>0$ such that for every core $T$ with its laws and the congruence property $(\Gamma)$ of
\cref{pf:onto3:gamat}, all size data $L$ with $\log g_i\le9\mathrm L+1$ ($i=0,1$), all $j$, all weights $w$ and all
$C\ge C_0$, $A\ge2$, every non-zero integral form $F\in\mathcal M_w$ of $B_T$ satisfies
\[
\frac1{[K_0:\Q]}\sum_{\sigma}\log\|F\|_\sigma\ \ge\ -\ontothreePP_0(C,A)\bigl(1+\log\ontothreesize(w)\bigr)
\ontothreesize(w),
\]
where $\ontothreePP_0(C,A)=C\bigl(fr+\mathrm L+\log d_E+|\log\Omega|\bigr)^A$, $\log d_E$ is the logarithmic discriminant
parameter of the host and $\log\Omega$ the period of $L$.
\formallabel{GlobalTheory.T.meanLogSupNorm_lower_bound_of_coreLaws}
\formallabel{GlobalTheory.lamwAt_of_lawsCore}
\formallabel{GlobalTheory.log_two_mul_mul_le}
\formallabel{Formalism.MainNotions.Gaps.EtaleWeightGaps.LAMwAt}
\end{proposition}

\begin{proof}
Let $C_0>0$ be the constant of \lref{GlobalTheory.T.exists_constant_cmEvaluation_fits_room}, an explicit numerical
inequality: for $f,r\ge1$, $p_0\le r$, $d\ge1$, $\mathrm L\ge0$, $M,k_0>0$, $S\ge1$ with $\log(Mk_0)\le19\mathrm L+3+\log d$, the
frame basis radius $\varrho$ attached to $(r,p_0,f,d,M,k_0,S)$ is at least $1$, and for $s,t\ge0$ with $s+t\le S$ and all
$C\ge C_0$, $A\ge2$ the CM evaluation exponent $(\mathrm{p}_{\deg}(\Omega)+\frac12\log\varrho)(s+t)+s\log(r\varrho)$ is at
most $C(fr+\mathrm L+\log d+|\Omega|)^A(1+\log S)S$. Let $F\ne0$ be integral of weight $w$; then $j$ is in range
(otherwise the forms are $0$). We work at the good embedding $\sigma_0$ of the core and its uniformisation, whose
components carry frame data.
\begin{enumerate}[label=(\arabic*)]
\item By the laws, $F$ is a holomorphic section of the automorphic bundle of weight $w$ on every component, its
pointwise norm is bounded, and its supremum is positive since $F\ne0$. Hence there are a component $k$ and a frame
point $z_0$ in its domain with $\sup|F|<2|F|(z_0)$ (\lref{Formalism.L2.Pj.abelian₂.half_max}).
\item By the congruence property $(\Gamma)$, the frame data of the component $k$ extend to lattice data $X_x$ with
level $g=M$ and lattice $\Lambda$ spanned by the multiples of $L'$ by an explicit set of algebraic integers
(\lref{GlobalTheory.LatticeData.dataXOfModel_of_gamAt}); for these data
\lref{GlobalTheory.TransvModel.transv_model} gives $M_kk_0=2g_0g_1\sqrt d$ with $M_k=g_0g_1\Nm(\mathfrak c_K)d$, and,
for every isotropic $v\in L'$ and every real algebraic integer $b$, an element of the arithmetic group acting as the
flow along $v$ by $M_kb$, on which the character of the weight is trivial (a law of the core).
\item By \lref{Formalism.L2.Pj.abelian₂.cm_frame_of_eq}, these data produce a CM frame point $x$ in the domain and
a lattice basis $\lambda$ of radius at most $\varrho$ at which the pointwise norm of $F$ is controlled by its value at
$z_0$. The weight sizes of the core (\lref{GlobalTheory.WtSizes.sizes}: $p+q=r$, $\sum_\tau\mathrm{spread}\le
\ontothreesize(w)$, $\sqrt d\le d_E$, $1\le d_E$) provide the inputs, and since $\log2\le1$, $\log g_i\le9\mathrm L+1$
and $\sqrt d\le d_E$, $\log(M_kk_0)=\log(2g_0g_1\sqrt d)\le19\mathrm L+3+\log d_E$.
\item By \lref{GlobalTheory.LamT.exists_kronecker_logOm_core}, $\log\Omega$ is the real part of the logarithmic period of
$E$ at a Kronecker character $\chi$. The sup norms of $B_T$ at $\sigma$ are the sup norms of the uniformisation at the
good embedding attached to $\sigma$, all positive. By
\lref{GlobalTheory.T.meanLogSupNorm_lowerBound_cmFrame_nonvanishing}, the lower bound for the mean of $\log$ sup
norms of an integral cusp form with a CM frame point as in (3) (a non-vanishing statement at CM points with the
period $\Omega$), the mean $\frac1{[K_0:\Q]}\sum_\sigma\log\|F\|_\sigma$ is at least minus the CM evaluation exponent with
$S=\ontothreesize(w)$.
\end{enumerate}
By the inequality of $C_0$ (with $d=d_E$, $f=[E:\Q]/2$ and $\Omega=\log\Omega$ of $L$) the exponent is at most
$\ontothreePP_0(C,A)(1+\log\ontothreesize(w))\ontothreesize(w)$.
\end{proof}

\begin{lemma}[Integral forms form a lattice]\label{pf:onto3:latat}
Let $T$ be a core with its laws and the congruence property $(\Gamma)$, $L$ size data with the parameter bounds (i) of
\cref{pf:onto3:pagedef} for some $C_A$, $X$ a valued form algebra over $B_T$ at $j$, $S$ a scalar quadruple with $s_K\ne0$
and $\kappa,w_L,w_F$ weights such that the short theta forms $\mathcal T^{\mathrm{sh}}(\kappa)$ span $\mathcal F_\kappa$ over
$K_0$ and condition (a) of \cref{pf:onto3:thetamsup} holds. Then $\mathcal M_\kappa$ is a lattice in $\mathcal F_\kappa$.
\formallabel{GlobalTheory.latAt_record}
\formallabel{ReciprocityTheory.ChartRestriction.G.PullbackSections.latOfLam}
\formallabel{Formalism.MainNotions.Gaps.EtaleWeightGaps.LatAt}
\end{lemma}

\begin{proof}
Finite generation: by \cref{pf:onto3:lamw} with $(C,A)=(C_0,2)$ (the hypothesis $\log g_i\le9\mathrm L+1$ is part of (i)),
every non-zero $F\in\mathcal M_\kappa$ satisfies $\sum_\sigma\log\|F\|_\sigma\ge-A_0$ with $A_0=[K_0:\Q]\ontothreePP_0(C_0,2)
(1+\log\ontothreesize(\kappa))\ontothreesize(\kappa)$. By \lref{discriminant_of_product_lower} an $\OO_{K_0}$-submodule of a
$K_0$-space with a family of norms whose products are bounded below on non-zero elements is uniformly discrete, and by
\lref{formalGroup_of_discriminant} a uniformly discrete submodule of a finite-dimensional space is finitely
generated. Spanning: let $c=(g_0g_1)^{frn_\kappa}\ne0$. For $x\in\mathcal T^{\mathrm{sh}}(\kappa)$, $s_Kc\,x\in\mathrm{Th}_{s_K}
(\kappa)\subseteq\mathcal M_\kappa$, so $x=(s_Kc)^{-1}(s_Kcx)$ lies in the $K_0$-span of $\mathcal M_\kappa$; as these $x$ span
$\mathcal F_\kappa$, so does $\mathcal M_\kappa$.
\end{proof}



\begin{proposition}[Chart independence of harmonic values]\label{pf:onto3:chartindep}
Let $\mathcal X$ be a moduli space of a CM host with a coefficient datum $c$, the involution of $\OO_E$ being the host
conjugation. For every $c.F$-algebra $S$ containing $\Q$, every framed point $y$ of $\mathcal X$ over $S$, every two charts
$C,C'$ at $y$, all $k,n$, all sections $F_1,\dots,F_n$ of the $k$-th power of the automorphic line bundle with germs $f_b$ in
$C$ and $f'_b$ in $C'$, every weight vector $a$ and every harmonic polynomial $P$ of weight $a$ in $n$ variables, the
harmonic values of $P$ on the germs in the two charts coincide.
\formallabel{GlobalTheory.T.harmonic_germ_value_chart_independent}
\formallabel{GlobalTheory.harmonicValue_chart_independent}
\end{proposition}

\begin{proof}
The statement is assembled from results proved elsewhere in the development. By
\lref{Formalism.MainNotions.ModuliIntegralCharts.H1.chartIndep2_verb}, it follows from four statements: chart
independence of truncated values, fixedness of truncations under changes of frames (in balanced form), fixedness of
germs, and harmonic invariance. Fixedness of germs is \lref{Formalism.MainNotions.ModuliIntegralCharts.H1.germFix_holds}.
Fixedness of truncations holds in the unbalanced form (\lref{Formalism.MainNotions.ModuliIntegralCharts.H1.frTruncFix_holds}),
and the two forms are equivalent by \lref{GlobalTheory.BalancedTwist.frTruncFix_iff}. The remaining two statements
rest on the exponential action of the chart directions: by \lref{GlobalTheory.EigDir.linEig_holds} the integral chart
directions are eigenvectors of the action of $\OO_E$, which gives the exponential action
(\lref{Formalism.MainNotions.ModuliIntegralCharts.H1.expActOf_of_eig}); from it follow harmonic invariance
(\lref{Formalism.MainNotions.ModuliIntegralCharts.H1.harmInv_of}), exponential charts
(\lref{Formalism.MainNotions.ModuliIntegralCharts.H1.expChartOf_of}) and exponential gauges
(\lref{Formalism.MainNotions.ModuliIntegralCharts.H1.ExpGauge.expGaugeOf_of}). Chart independence of truncated values
follows (\lref{Formalism.MainNotions.ModuliIntegralCharts.H1.chartIndepTrunc_of}) from exponential charts and gauges,
the tie statement (\lref{Formalism.MainNotions.ModuliIntegralCharts.H1.tie_holds}), uniqueness of tangents
(\lref{Formalism.MainNotions.ModuliIntegralCharts.H1.ExpUniqChart.expUniqTan}), uniqueness of gauges, which follows from
the statement on global functions (\lref{Formalism.MainNotions.ModuliIntegralCharts.H1.gaugeUniq_of_global},
\lref{GlobalTheory.globalFns_holds}), and the vanishing of additive valuations, which follows from the
$\psi$-covers of the host (\lref{Formalism.MainNotions.ModuliIntegralCharts.H1.killAddVal_of},
\lref{Formalism.MainNotions.ModuliIntegralCharts.H1.psiCover_holds}).
\end{proof}

\begin{proposition}[Generators without common zero]\label{pf:onto3:ncz}
Assume chart independence (\cref{pf:onto3:chartindep}). Let $T$ be a top structure with laws, $j$ in range, $S$ a scalar
quadruple for $B_T$ with $s_W\ne0$, and assume $\mathcal H_{\mathrm{page}}$ for $X_{T,j}$. Let $G$ be a global family datum
for $(X_{T,j},\kappa,w_L,w_F)$ satisfying $(\mathrm G1)$ (of which only the fact that its forms $u_1,\dots,u_N$ span
$W_{s_W}(w_L)$ over $\OO_{K_0}$ is used). Then
the $u_k$ have no common zero: for every algebraically closed field $k$ over $c.F$ and every $k$-point of $\mathcal X_j$,
some $u_k$ does not vanish there.
\formallabel{GlobalTheory.T.scaledThetaLattice_generators_noCommonZero}
\formallabel{GlobalTheory.RunIn.noCommonZero_of_tiedW}
\formallabel{GlobalTheory.FrameZero.noCommonZero_of_readings}
\formallabel{GlobalTheory.RunIn.eq_zero_of_span_eq_WWc}
\formallabel{GlobalTheory.FrameZero.exists_complex_point}
\end{proposition}

\begin{proof}
\emph{Reduction to frame points.} Let $\mathrm{rd}_{k,x}(F)\in\C$ denote the reading of a rule $F$ at a frame point $x$ of
the component $k$ of the uniformisation at the good embedding $\sigma_0$. We claim that if at every component $k$ and every
frame point $x$ of its domain some $u_i$ has non-zero reading, then the $u_i$ have no common zero. Suppose all $u_i$
vanish at a geometric point $y$. The atlas of the moduli space (\lref{GlobalTheory.Reg.evLaws_genFr_atlas}) provides an
\'etale chart $C$ with a homomorphism $\psi$ to the field of definition of $y$ and a frame $z$ over it; the evaluations
of the $u_i$ at $z$ vanish (an evaluation at a frame is the product of the evaluation at the point with a non-zero
weight character), so the ideal generated by the evaluations $\mathrm{ev}(u_i)$ in the coordinate ring of $C$ lies in
$\ker\psi$ and is proper. Hence it has a complex point $\psi'$ over $\sigma_0$: a maximal ideal containing it has as
quotient a finite extension of the coefficient field (Nullstellensatz), which embeds into $\C$ over $\sigma_0$. By base
change and the uniformisation of the laws, $\psi'$
corresponds to a frame point $x$ of some component $k$; at $x$ every $u_i$ has zero reading, a contradiction.

\emph{Non-vanishing of a theta rule.} Let $k,x$ be given. By \cref{pf:onto3:below}, $\ontothreeDF(w_L)\le\ontothreeDF(\kappa)$,
$r\le n_{w_L}$ and the last entries of $w_L$ are positive; the size condition $2\ontothreeDF(\kappa)+2<g_i$ and $M=g_0g_1$
give $4\ontothreeDF(w_L)+4\le M$. By \lref{GlobalTheory.noCommonZero_of_chartIndep} (under chart independence, these
conditions and the norm of the half-level generator: some theta rule of weight $w_L$ has non-zero reading at every frame
point) there is a theta rule $f$ of $T$ at $(j,w_L)$ with $\mathrm{rd}_{k,x}(f)\ne0$; it is a cusp form, hence $f\in
\mathcal T(w_L)$.

\emph{Conclusion.} The reading $\varphi=\mathrm{rd}_{k,x}$ is additive and $\sigma_0$-semilinear. If $\varphi(u_i)=0$ for all
$i$, then $\varphi$ vanishes on the $\OO_{K_0}$-span $W_{s_W}(w_L)=s_W\mathcal T(w_L)$, in particular $\sigma_0(s_W)\varphi(f)=
\varphi(s_Wf)=0$; as $s_W\ne0$, $\varphi(f)=0$, a contradiction.
\end{proof}

\subsubsection{Auxiliary forms and generating families}

\begin{proposition}[Auxiliary form]\label{pf:onto3:etaok}
Let $T$ be a top structure with laws, $j$ in range with $j<m$, and assume $\mathcal C^+_{\mathrm{page}}$ for $X_{T,j}$. Then
there are $i\in\N$ and $\eta\in\mathcal F_{\epsilon_i}$ satisfying the auxiliary form condition of
\cref{pf:onto1:genok-def}: $i>0$; $\eta$ is a non-zero cusp form which is not a zero divisor; $v_{\mathfrak l,y}(\iota\eta)\ge
2i\phi_C(\mathfrak l,y)$ for all $\mathfrak l,y$; $\log\|\eta\|_\sigma/i\le\ontothreePP(200000,2)$ for all $\sigma$; the
archimedean bound at frames with the constant $200000z_B^2$; and $\eta\in\mathcal M_{\epsilon_i}$.
\formallabel{GlobalTheory.EtaBoundForm.exists_etaOK}
\formallabel{GlobalTheory.RunJoin.EtaOK}
\formallabel{GlobalTheory.EtaInt.etVal_integral}
\end{proposition}

\begin{proof}
Let $(Q,T_p)$ be the page parameters of $\mathcal C_{\mathrm{page}}$. By \lref{GlobalTheory.EtaAn.page_rows}, the
parameter bounds (i) and the size conditions give $2\le\mathrm L$, $\log r\le\mathrm L$, $\log\mathrm{rd}_E\le2\mathrm L$,
$\log f\le2\mathrm L$, $\log g_i\le9\mathrm L+1$ and $16\le f$. The ties of $\mathcal C_{\mathrm{page}}$ give $(\mathrm K)$ for
$\mathbf K=(1,2,1,2,1,2,1,4,7)$ and $(C_h,A_h,C_q,A_q)=(1,4,0,11)$, and the page support holds for the scalar quadruple
$((g_0g_1)^{\ontothreeDF(w_L)},1,1,1)$: its first entry has $\log|\sigma(\cdot)|=\ontothreeDF(w_L)\log(g_0g_1)\ge
\frac12\ontothreeDF(w_L)\log(g_0g_1)$, and the other items are part of $\mathcal C^+_{\mathrm{page}}$. Hence
\cref{pf:onto3:runrec}(i) shows that $(q,R_0,m)=(g_0g_1,\sqrt{90(fr+\ontothreeDF(w_L))},m)$ are admissible run parameters.
Moreover $1\le g_0g_1$ and $\log(g_0g_1)\le18\mathrm L+2\le70\mathrm L$. By
\lref{GlobalTheory.T.etaleExtension_degreeTwo_constants} (whose proof combines
\lref{GlobalTheory.T.exists_cuspForm_four_clauses_explicitConstant}, the existence of a cusp form with the four
clauses below for admissible run parameters, with \lref{GlobalTheory.T.log_card_components_le_of_rows_bound}, a bound
for the number of components, and an explicit numerical estimate), under these bounds there are $i>0$ and a cusp form
$\eta$ of weight $\epsilon_i$ which is not a zero divisor, has $v_{\mathfrak l,y}(\iota\eta)\ge2i\phi_C(\mathfrak l,y)$, satisfies
the sup bound with $\ontothreePP(200000,2)=200000z_B^2$ and the archimedean bound at frames for the theta constant of
$(q,R_0)$ with the constant $200000z_B^2$. Finally $\eta$ is integral: the valuation condition gives
$v_{\mathfrak l,y}(\iota\eta)\ge0$ at all points of $X_{T,j}$, because $\phi_C\ge0$
(\lref{GlobalTheory.EtaValuationTest.phiC_nonneg}), and the integrality condition of $X_{T,j}$ gives $\eta\in
\mathcal M_{\epsilon_i}$.
\end{proof}

\begin{lemma}[Refilling generating families]\label{pf:onto3:refill}
Admissibility of a global family datum (\cref{pf:onto1:genok-def}) depends only on its components $N$, $u_k$ and $m_x$;
replacing all other components (the forms $\theta'_i$, their number $\Delta$, the zero family, the auxiliary form and
the remaining data) gives an admissible datum.
\formallabel{GlobalTheory.RunJoin.genOK_refill}
\formallabel{GlobalTheory.RunJoin.GenOK}
\end{lemma}

\begin{proof}
The four conditions of admissibility involve only $N$, the $u_k$ and $m_x$, and these are unchanged by the
replacement.
\end{proof}

% ---------------------------------------------------------------------------
\subsection{Remarks on the formal statements}\label{sec:pf:onto3:remarks}

\begin{remark}\label{pf:onto3:rem:divinn}
The hypothesis $\sfname{DIVinN}(\mathcal F)$ in the definition of $\sfname{SurjectiveOntoX}$ is not used: in the
proof of \cref{pf:onto3:path} the inequality $\sfname{ON}$ is established for all families of $\sfname{FamHyp}$ which
are large in the sense of \cref{def:families:closed}, whether or not they satisfy $\sfname{DIVinN}$.
\end{remark}

\begin{remark}\label{pf:onto3:rem:unused}
Several hypotheses of the formal statements are not used. (a) The proof of $\sfname{RecTExists}$ in
\cref{pf:onto3:main} passes twelve fields of $P$ to a lemma which returns the first one; the other eleven (local
converse theorem, depth-zero induction, Brauer's induction theorem in degree zero, Chevalley's theorem on units,
global Artin reciprocity, exhaustion and multisegments for $\GL_n$, standard cosocles, Zelevinsky's polynomial ring,
the geometric lemma, Bernstein--Zelevinsky finiteness) play no role there. (b) In \cref{pf:onto3:fivest} the exported statement is a
hypothesis which is not used; it enters the proof of \cref{pf:onto3:ontoX} through the twist statement
(\cref{pf:onto3:twist}). The page, geometric and rank statements are derived twice, once in the proof of
\cref{pf:onto3:ontoX} and once in that of \cref{pf:onto3:fivest}. (c) The second Siegel--Weil statement
\ref{pf:onto3:P:sw} is a hypothesis of \cref{pf:onto3:adelc} which is not used; the low-weight bound uses the other
six results listed in the proof of \cref{pf:onto3:main}.
\end{remark}

\begin{remark}\label{pf:onto3:rem:canon}
Since the error function of the first side of a point is the canonical one \eqref{eq:pf:onto3:erreq}, conditions
(ON1) and (ON2) of $\sfname{ON}$ hold automatically as soon as $D(S_{\mathrm f})>0$ (this is the case $\ontothreeerr=
\ontothreeerr_0$ of the first two parts of \cref{pf:onto3:canonical}); the content of $\sfname{SurjectiveOntoX}$ is
the global inequality (ON3) for the canonical error, i.e.\ a bound for the weighted sum over the twisted places of
the excess of the local indices over the local budget. The canonical error is defined with division by
$D(S_{\mathrm f})$, which is read as $0$ if $D(S_{\mathrm f})=0$; in the proof $D(S_{\mathrm f})>0$ is established
from the rank rows.
\end{remark}

\begin{remark}\label{pf:onto3:rem:l6}
In $\sfname{LambdaLowerBound}_2$ the numbers $l_6(j)=l_{6m}(j)-C\,l_{6u}$ do not depend on $C$, because the scale of
a point has $l_{6u}=0$ by definition; the statement is the inequality $l_{6m}(j)\le\log\lambda_{j,\mu_0,p',v}$. By the
conventions (a) and (b) of \cref{sec:pf:onto3:low}, the scalar $\lambda$ is set equal to $\exp(l_{6m}(j))$ when the
corresponding type part is zero (at $v_0$) or degenerate (at $v\ne v_0$); in these cases the inequality is an
equality and carries no information.
\end{remark}

\begin{remark}\label{pf:onto3:rem:scope}
All statements of this section are closed statements over $\sfname{FamHyp}$, and their conclusions concern the
points satisfying the point conditions (and the auxiliary conditions). As noted after \cref{def:families:closed},
they say nothing about classes without large members, and the existence of points and of record objects with the
record conditions is not part of them, except that \lref{GlobalTheory.BigOn.exists_isRec} provides a record
object for every point with runs. The runs, rationality and maximality statements (\cref{pf:onto3:runs}) hold for
every family and every scale, without largeness.
\end{remark}

\begin{remark}\label{pf:onto3:rem:mackey}
The local index bound (\cref{pf:onto3:lcw}) is a statement about lattices with an action of a finite group and an
eigenspace selected by a Mackey pair; it is not specific to theta lifts. The only property of the mirabolic groups
used is \cref{pf:onto3:mirabolic}, together with their orders (\cref{pf:onto3:mirorders}), which enter through the
leaves \ref{pf:onto3:L:indexU} and \ref{pf:onto3:L:unip}. The numerical constants of the leaves ($2.04$, $1/8$,
$4C^2+25C+C_b+44$, the conditions $14\le C_{b1}$ and $66C+4\le C_{b2}$) are those of the development.
\end{remark}

\begin{remark}\label{pf:onto3:rem:ties}
In the definitions of \cref{sec:pf:onto3:ties,sec:pf:onto3:runw}: (a) the class condition imposed on the valued points
of the enlargement $E_{\mathrm{all}}$ is the true proposition, so $E_{\mathrm{all}}$ consists of all valued points of the
moduli space; (b) in the normalised depth an infinite local depth is replaced by $0$ (the development reads an extended
real number $\pm\infty$ as the real number $0$), and the exponent $e_\Lambda$ is truncated at $0$; (c) the record
conditions ask for the depth tie for \emph{some} involution of $\OO_E$, but the depth tie forces this involution to be
the host conjugation; (d) $\mathrm{th}_k(\ell)=0$ only for $\ell\le1$, since for $\ell\ge2$ odd powers $\ell^i$ exceed
$k^{3/2}$; (e) the two formal versions of the page-row bounds have identical definitions.
\end{remark}
